\documentclass{article}
\usepackage[utf8]{inputenc}
\usepackage[T2A,T1]{fontenc}
\usepackage[english,russian]{babel}
\usepackage{amsmath}
\usepackage{amssymb,amsfonts,textcomp}
\usepackage{fancyhdr}
\usepackage{array}
\usepackage{longtable}
\usepackage{hhline}
\usepackage[pdftex]{graphicx}
\usepackage[colorlinks=true,linkcolor=blue,unicode, pdftex]{hyperref}
\makeatletter
\newcommand\arraybslash{\let\\\@arraycr}
\makeatother
\setlength\tabcolsep{1mm}
\renewcommand\arraystretch{1.3}
\newcounter{Table}[section]
\renewcommand\theTable{\thesection.\arabic{Table}}
\parindent=1cm
\usepackage[letterpaper, left=3cm, right=2.5cm, top=2.5cm, bottom=2.5cm]{geometry}
\usepackage{color}
\usepackage{framed}
\usepackage{titlesec}
\usepackage{fancybox}
\usepackage[tikz]{bclogo}
\usepackage{xhfill}
\usepackage{booktabs}
\usepackage{listingsutf8}
\usepackage{float}
\usepackage{tablefootnote}
\usepackage{colortbl} 
\usepackage{type1ec} % для четкости шрифта
\usepackage[numbered]{bookmark} % для корректного отображения оглавления в ридере
\usepackage{arydshln} %для пунктирных линий
\usepackage[framemethod=tikz]{mdframed}
\usepackage{multirow}


\mdfsetup{%
	middlelinewidth=0.5pt,
	backgroundcolor=white!98!black,	
	roundcorner=5pt,
	innerleftmargin = 10pt,
	innertopmargin=5pt,
	skipabove = 10pt,
	skipbelow = 5pt,
	shadow=true,
	shadowsize=4pt,
	nobreak=false,
	rightmargin=0pt,
	align=right, 
}
\global\mdfdefinestyle{importantstyle}{%
	linecolor=red,linewidth=4pt,%
	topline=false,bottomline=false, rightline = false, %
	backgroundcolor=red!8,%
	shadow=false
}

\global\mdfdefinestyle{examplestyle}{%
	middlelinewidth=0.5pt,
	backgroundcolor=white!98!black,	
	roundcorner=5pt,
	innerleftmargin = 10pt,
	innertopmargin=5pt,
	skipabove = 10pt,
	skipbelow = 5pt,
	shadow=true,
	shadowsize=4pt,
	nobreak=false,
	rightmargin=0pt,
	align=right, 
}


\global\mdfdefinestyle{remarkstyle}{%
	linecolor=green!50!black,linewidth=4pt,%
	topline=false,bottomline=false, rightline = false, %
	backgroundcolor=green!8,%
	shadow=false
}

\mdfdefinestyle{theoremstyle}{%
	frametitlebackgroundcolor=white!40!black,
	frametitlealignment = \raggedright,
	frametitlerule=true,
	frametitlefont={\normalfont\bfseries\color{white}},
	frametitlebelowskip = 4pt,
	frametitleaboveskip = 4pt,
}


\newcounter{redlisting}[section]
\renewcommand{\theredlisting}{\thesection.\arabic{redlisting}}
\newcommand{\li}{\refstepcounter{redlisting} Листинг \theredlisting. }
\newcommand{\enli}{\refstepcounter{redlisting} Listing \theredlisting. }


\newenvironment{redlisting}[1]{\begin{mdframed}[style=theoremstyle,frametitle={\li #1}]\flushleft}{\end{mdframed}}

\newenvironment{redexample}{\begin{mdframed}\flushleft}{\end{mdframed}}

\newenvironment{redimportant}{\begin{mdframed}[style=importantstyle]}{\end{mdframed}}

\newenvironment{redremark}{\begin{mdframed}[style=remarkstyle]}{\end{mdframed}}


\usepackage{comment}

%\usepackage{parskip}
\usepackage{enumitem}

\usepackage{indentfirst}  % абзацный отступ после заголовка
\graphicspath{{SQLGuide-img/}} %

% другой счетчик для приложений (по \ref выдает A,B,C...)
\newcounter{application}
\renewcommand{\theapplication}{\Asbuk{application}}
\newcommand{\appcount}{\refstepcounter{application}}

\newcounter{test}
\renewcommand{\thetest}{\Asbuk{test}}
\newcommand{\apptest}{\refstepcounter{test}}

% чтобы в оглавлении было не присто А Б, а Приложение А, Приложение Б
\renewcommand\appendixname{Приложение}

\makeatletter
\def\redeflsection{\def\l@section{\@dottedtocline{1}{0em}{8em}}}
\renewcommand\appendix{\par
	\setcounter{application}{0}%
	\setcounter{subsection}{0}%
	\def\@chapapp{\appendixname}%
	\addtocontents{toc}{\protect\redeflsection}	
	\def\thesection{\appendixname\hspace{0.2cm}\Asbuk{application}}
	}
\makeatother


% чтобы разделы в оглавлении были не жирным а обычным шрифтом
\makeatletter 
\renewcommand*\l@section[2]{% 
	\ifnum \c@tocdepth >\z@ 
	\addpenalty\@secpenalty 
	\addvspace{1.0em \@plus\p@}% 
	\setlength\@tempdima{1.5em}% 
	\begingroup 
	\parindent \z@ \rightskip \@pnumwidth 
	\parfillskip -\@pnumwidth 
	\leavevmode 
	\advance\leftskip\@tempdima 
	\hskip -\leftskip 
	#1\nobreak\hfil \nobreak\hb@xt@\@pnumwidth{\hss #2}\par 
	\endgroup 
	\fi} 
\makeatother 


\emergencystretch=25pt %борьба с переполнением строк


\renewcommand{\thetable}{\thesection.\arabic{table}} % нумерация таблиц с номером section
\renewcommand{\thefigure}{\thesection.\arabic{figure}} % нумерация рисунков с номером section

\makeatletter
\@addtoreset{table}{section} % сбросит счетчик таблиц в начале каждого раздела. 
\makeatother

\makeatletter
\@addtoreset{figure}{section} % сбросит счетчик рисунков в начале каждого раздела. 
\makeatother

% пакет оформления блоков "Внимание" и "Пример"
\usepackage{tcolorbox}
\tcbuselibrary{most}

%пакет для работы с плавающими объектами
\usepackage{afterpage}

%оформление таблиц и рисунков по ГОСТ
\usepackage[tableposition=top]{caption}
\usepackage{subcaption}
\DeclareCaptionLabelFormat{gostfigure}{Рисунок #2} 
\DeclareCaptionLabelFormat{gosttable}{Таблица #2} 
\DeclareCaptionLabelSeparator{gost}{~---~}
\captionsetup{labelsep=gost}
\captionsetup[figure]{labelformat=gostfigure}
\captionsetup[table]{labelformat=gosttable,justification=raggedright,slc=off} 

\newcommand{\vspacebf}[1]{\vspace{5pt}\noindent \textbf{#1}}

%задает цвет гиперссылок 
\hypersetup{urlcolor=blue}

\tcbset{
	enhanced,
	colback=black!2,
	boxrule=0.4pt,
	fonttitle=\bfseries
}



%задает пустое пространство до и после таблицы
\setlength{\LTpre}{5pt}
\setlength{\LTpost}{5pt}

\newenvironment{redparam}
{\vspace{10pt}\bf\parindent=0cm}
{\par}

\newenvironment{redparam1}
{\vspace{5pt}\ttfamily\parindent=0cm}
{}


\newenvironment{rederror}
{\par\vspace{5pt}\ttfamily}
{\par\vspace{5pt}}

\newcommand{\sectionbreak}{\clearpage}

\titleformat{\section}
{\huge}
{\thesection}{1em}{}

\titleformat{\subsection}
{\LARGE}
{\thesubsection}{1em}{}

\titleformat{\subsubsection}
{\Large}
{\thesubsubsection}{1em}{}

\titleformat{\paragraph}
{\large}
{\theparagraph}{1em}{}

  
\title{}
\author{Red Soft Corporation}
\date{}

\setlength{\headheight}{24pt}
\pagestyle{fancy}


\newcommand{\anonsection}[1]{\phantomsection\section*{#1}\addcontentsline{toc}{section}{#1}}


\newcommand{\anonsubsection}[1]{\phantomsection\subsection*{#1}\addcontentsline{toc}{subsection}{#1}}

\newcommand{\anonsubsubsection}[2][]{\phantomsection%
	\subsubsection*{#2}%
	\ifthenelse{\equal{#1}{}}{\addcontentsline{toc}{subsubsection}{#2}}{\addcontentsline{toc}{subsubsection}{#1}}%
	}





\newcommand{\beginrudoc}[1]{
\fancyhf{}
\thispagestyle{empty}

\topskip0pt
\vspace*{\fill}

\begin{flushright}
\Huge {\xhrulefill{red}{2mm}\color{red} Red Expert} \\
\LARGE Версия 2020.04\\
#1
\end{flushright}

\vspace*{\fill}

\newpage

\fancyhead[R]{#1\\\rightmark}
\fancyfoot[C]{\xhrulefill{red}{2mm} Стр. \thepage}

\setcounter{tocdepth}{10}
\renewcommand\contentsname{Содержание}
\tableofcontents
}


%печать одинарных кавычек
\newcommand{\rr}[1]{{\textquotesingle}#1{\textquotesingle}}

%печать двойных кавычек
\newcommand{\pp}[1]{"#1{}"}

\newcommand{\ttt}{\texttt}

%для приложения 5 в SQL
\newcommand{\hl}[1]{\hyperlink{#1}{#1}}

\setlist[enumerate]{itemindent=1.3mm, topsep = 1mm, parsep=0pt, leftmargin=45pt} %чтобы списки были более сжатыми
\setlist[itemize]{itemindent=1.3mm, topsep = 1mm, parsep=0pt, leftmargin=45pt}   %чтобы списки были более сжатыми



\begin{document}

\beginrudoc{Руководство пользователя}

\anonsection{Введение}

Red Expert (Ред Эксперт) - это GUI-оболочка, предоставляющая интерфейс для работы с базами данных Firebird и Ред Базой Данных. Написанная полностью на Java, Red Expert является программой, независимой от операционной системы, для которой требуется только JDK не ниже 1.8.

Используя гибкость, предоставляемую Java Database Connectivity (JDBC), Red Expert обеспечивает простой способ взаимодействия с базой данных от простых запросов до создания таблиц и импорта / экспорта данных всей базе данных.

Red Expert был протестирован под Windows, Linux с использованием Java 1.8.

Red Expert обладает рядом полезных компонентов: визуальным редактором для всех объектов базы данных, редактором SQL запросов и исполнителем SQL скриптов и др.

\newpage

\section{Установка Red Expert}

\begin{enumerate}[leftmargin=15pt]
	\item Установите JDK не ниже 1.8.
	\item Скачайте дистрибутив Red Expert для соответствующей операционной системы и ее разрядности с официального сайта СУБД Ред Базы Данных --- \href{http://reddatabase.ru/downloads/redexpert/}{reddatabase.ru}. Загрузка доступна только авторизованному пользователю.
	\item Запустите установку Red Expert с помощью файла \ttt{red\_expert\_installer-<версия>-xxxxxx}. Инсталляция Red Expert осуществляется с помощью стандартного мастера установки программ.
	
	В начале Вам предложат выбрать язык установки. Предусмотрена установка на русском и английском языках. Далее Вы увидите подобные окна:

\begin{figure}[H]
	\begin{center}
		\begin{minipage}[h]{0.49\linewidth}
			\includegraphics[width=1\linewidth]{img/install_1.png}
			\caption{Приветственное окно установки} 
		\end{minipage}
		\hfill 
		\begin{minipage}[h]{0.49\linewidth}
			\includegraphics[width=1\linewidth]{img/install_2.png}
			\caption{Лицензионное соглашение}
		\end{minipage}
	\end{center}
\end{figure}

\begin{figure}[H]
	\begin{center}
		\begin{minipage}[h]{0.49\linewidth}
			\includegraphics[width=1\linewidth]{img/install_3.png}
			\caption{Выбор каталога установки} 
		\end{minipage}
		\hfill 
		\begin{minipage}[h]{0.49\linewidth}
			\includegraphics[width=1\linewidth]{img/install_4.png}
			\caption{Окончание установки}
		\end{minipage}
	\end{center}
\end{figure}

\item По окончанию процесса установки появится иконка 
\begin{tikzpicture}
\pgftext{\includegraphics[width=14pt]{img/red_expert.png}} at (0pt,0pt)
\end{tikzpicture}
на рабочем столе.

\item Запустите Red Expert.

\begin{redremark}
Если при запуске приложению не удастся найти путь установки Java, то пользователю будет предложено выбрать путь вручную, либо запустить автоматическое скачивание Java:
\begin{itemize}[leftmargin=15pt]
	\item На Linux системах для автоматического скачивания Java должен быть установлен пакет \ttt{libcurl4}.
	\item На Windows Server для автоматического скачивания Java нужно отключить усиленную безопасность InternetExplorer.
	\item Если для вашей системы не предусмотрен инсталятор или возникают ошибки с запуском, то можно запустить приложение командой "java -jar RedExpert.jar" из каталога в котором находится RedExpert.jar
\end{itemize}
\end{redremark}
\end{enumerate}	

\newpage

\section{Обновление Red Expert}
\begin{enumerate}[leftmargin=25pt]
\item При выходе новой версии приложения в строке состояния (внизу главного окна)
будет появляться мигающая точка желтого цвета с надписью \ttt{Доступно обновление}:
\begin{figure}[H]
	\begin{flushright}
		\includegraphics[width = 0.94\linewidth]{img/status_bar.png}
	\end{flushright}
\end{figure}

\item Нажмите на эту панель и согласитесь на все предлагаемые действия.

\item По окончании обновления перезапустите Red Expert.
\end{enumerate}	

Проверять наличие обновлений можно также через пункт меню \textit{Справка} $\rightarrow$ \textit{Проверить обновления}.

\newpage
\section{Обзор графического интерфейса}\label{sec:overview}

В Red Expert используется распространенный для таких инструментов и даже большинства IDE интерфейс с вкладками и подключаемыми компонентами.

Все окно приложения можно визуально разделить на  \textit{центральную, левую, правую, верхнюю} и  \textit{нижнюю панели}.

На \textit{центральной панели} окна приложения в виде вкладок размещаются основные функциональные компоненты, такие как:

\anonsubsubsection[Браузер баз данных]{\hyperref[sec:connections]{Браузер баз данных}}

С его помощью выполняется подключение и отключение от базы данных, просматриваются все настройки подключения.
	\begin{figure}[H]
		\centering
		\includegraphics[width = 0.8\linewidth]{img/browser_db.png}
	\end{figure}
\anonsubsubsection[Компонент для установки JDBC драйверов]{\hyperref[sec:new_driver]{Компонент для установки JDBC драйверов}}

Предназначен для регистрации в Red Expert других JDBC драйверов помимо Jaybird Driver, а также просмотра и редактирования их настроек.
	\begin{figure}[H]
		\centering
		\includegraphics[width = 0.8\linewidth]{img/new_drivers.png}
	\end{figure}
	\begin{figure}[H]
		\centering
		\includegraphics[width = 0.8\linewidth]{img/driver_show.png}
	\end{figure}

\anonsubsubsection[Компонент для создания новой базы данных]{\hyperref[sec:create_db]{Компонент для создания новой базы данных}}

В нем настраиваются параметры подключения к создаваемой базе данных, в том числе расширенные.
\begin{figure}[H]
	\centering
	\includegraphics[width = 0.8\linewidth]{img/create_db.png}
\end{figure}

\anonsubsubsection[Обозреватели и редакторы объектов базы данных]{\hyperref[sec:objects]{Обозреватели и редакторы объектов базы данных}}

Каждый созданный объект базы данных (процедура, индекс, таблица и т.д.) можно просматривать и изменять через графический интерфейс.
	\begin{figure}[H]
		\centering
		\includegraphics[width = 0.8\linewidth]{img/table_explor.png}
	\end{figure}

\anonsubsubsection[Редактор запросов]{\hyperref[sec:query_editor]{Редактор запросов}}

Инструмент просмотра и выполнения операторов SQL.
	\begin{figure}[H]
		\centering
		\includegraphics[width = 0.8\linewidth]{img/query_editor.png}
	\end{figure}

\anonsubsubsection[Черновик]{\hyperref[sec:draft]{Черновик}}

Текстовый редактор для временного отображения или редактирования текста по мере необходимости.
	\begin{figure}[H]
		\centering
		\includegraphics[width = 0.8\linewidth]{img/scratch_pad.png}
	\end{figure}

\anonsubsubsection[Менеджер пользователей]{\hyperref[sec:user_manager]{Менеджер пользователей}} 

Для управления пользователями базы данных (создания, изменения и удаления) через графический интерфейс.
	
	\begin{figure}[H]
		\centering
		\includegraphics[width = 0.8\linewidth]{img/user_manager.png}
	\end{figure}

\anonsubsubsection[Менеджер прав]{\hyperref[sec:priv_manager]{Менеджер прав}} 

Графический интерфейс для назначения и отзыва привилегий.
	\begin{figure}[H]
		\centering
		\includegraphics[width = 0.8\linewidth]{img/priv_manager.png}
	\end{figure}

\anonsubsubsection[Трейс менеджер]{\hyperref[sec:trace_manager]{Трейс менеджер}}  

Позволяют серверу отслеживать и записывать в лог-файлы все события, происходящие с базой данных, такие как: соединение и отсоединение от БД (создание и удаление БД), операции DML и DDL, выполнение хранимых процедур и т.д.
	\begin{figure}[H]
		\centering
		\includegraphics[width = 0.8\linewidth]{img/trace_manager.png}
	\end{figure}

\anonsubsubsection[Генератор тестовых данных]{\hyperref[sec:test_data_generator]{Генератор тестовых данных}}

Инструмент для быстрой генерации тестовых данных в базах.  С помощью него вы сможете быстро заполнять любую таблицу Ред Базы Данных сгенерированными данными с помощью следующих методов:
\begin{itemize}
	\item случайно;
	\item данными из полей другой таблицы;
	\item по спискам значений, заданных вручную;
	\item по спискам значений из файла;
	\item автоинкремент (для числовых типов данных).
\end{itemize}
В методах можно настраивать различные параметры генерации для каждого типа поля. 

\begin{figure}[H]
	\centering
	\includegraphics[width = 1\linewidth]{img/test_data_generator_1.png}
\end{figure}

\anonsubsubsection[Обозреватель системного журнала состояний вывода]{\hyperref[sec:logout]{Обозреватель системного журнала состояний вывода}}

Отображение всей информации, которая выводится в стандартный поток вывода, ошибки и предупреждения.
	\begin{figure}[H]
		\centering
		\includegraphics[width = 0.8\linewidth]{img/system_log.png}
	\end{figure}

\anonsubsubsection[Построитель ER-диаграмм]{\hyperref[sec:erd]{Построитель ER-диаграмм}} 

Графическая визуализация схемы базы данных.
	\begin{figure}[H]
		\centering
		\includegraphics[width = 0.8\linewidth]{img/er_diagram.png}
	\end{figure}

\anonsubsubsection[Модуль сравнения баз данных]{\hyperref[sec:compareDB_module]{Модуль сравнения баз данных}}

Инструмент сравнивает две базы данных и генерирует \ttt{SQL-скрипт}, позволяющий сделать две базы идентичными. 

\begin{figure}[H]
	\centering
	\includegraphics[width = 0.8\linewidth]{img/compareDB_module.png}
\end{figure}

\anonsubsubsection[Импорт данных из файла]{\hyperref[sec:import_data_from_file]{Импорт данных из файла}}

Инструмент для импорта данных из файла в таблицу. Поддерживается импорт из XLSX, XML и CSV файлов.

\begin{figure}[H]
	\centering
	\includegraphics[width = 0.95\linewidth]{img/ImportData.png}
\end{figure}

\anonsubsubsection[Компонент для сравнения типов данных]{\hyperref[sec:compare_datatype]{Компонент для сравнения типов данных}}

Для поиска аналогичных типов данных двух разных СУБД, которые наиболее точно соответствуют друг другу.
	\begin{figure}[H]
		\centering
		\includegraphics[width = 0.8\linewidth]{img/compare_types2.png}
	\end{figure}

\anonsubsubsection[Компонент для выполнения хранимых процедур и функций]{\hyperref[sec:execute_proc]{Компонент для выполнения хранимых процедур и функций}}

Выполнение хранимых объектов через графический интерфейс.
	\begin{figure}[H]
		\centering
		\includegraphics[width = 0.7\linewidth]{img/execute_store_proc.png}
	\end{figure}
	
\anonsubsubsection[Экспорт результатов]{\hyperref[sec:export_results]{Экспорт результатов}} 

Данная опция предназначена для экспорта данных таблиц, представлений, а также результатов выборки \ttt{SELECT} в различные форматы. 
	\begin{figure}[H]
		\centering
		\includegraphics[width = 0.6\linewidth]{img/export_to_sql.png}
	\end{figure}

\anonsubsubsection[Экспорт метаданных БД]{\hyperref[sec:export_db_meta]{Экспорт метаданных БД}}

Инструмент \ttt{Экспорт метаданных БД} извлекает метаданные выбранной БД в \ttt{SQL-скрипт}, выполнение которого позволяет создать дубликат этой базы.

\begin{figure}[H]
	\centering
	\includegraphics[width = 0.8\linewidth]{img/export_db_meta.png}
\end{figure}

\anonsubsubsection[Компонент для выполнения скриптов из файла]{\hyperref[sec:execute_scripts]{Компонент для выполнения скриптов из файла}}
Можно заранее написать код из SQL команд во внешнем редакторе и сохранить его в файле операционной системы. Red Expert позволяет работать с такими файлами: считывать их в буфер, запускать на выполнение и сохранять набранные команды в файле.

	\begin{figure}[H]
		\centering
		\includegraphics[width = 0.8\linewidth]{img/execute_script.png}
	\end{figure}

Список всех открытых в настоящее время вкладок можно посмотреть, нажав на кнопку \begin{tikzpicture}
\pgftext{\includegraphics[width=14pt]{img/button_alltabs.png}} at (0pt,0pt)
\end{tikzpicture} в правой верхней части центральной панели.


\textit{Левая, нижняя и правая} части окна приложения предназначены для прикрепляемых панелей. Их можно расположить в любой из трех позиций по желанию. Открыть или скрыть нужную панель можно с помощью меню \textit{Вид}. Открытые панели могут быть свернуты с помощью кнопки в области заголовка панели. Приложение сохраняет последнюю позицию каждого компонента при перезапуске.

В настоящее время доступны следующие прикрепляемые панели:

\anonsubsubsection{Консоль вывода}

Основным журналом вывода системы является файл \ttt{system.log} (\ttt{\$HOME/.redexpert/logs}). Консоль вывода приложения, обычно расположенная в виде прикрепленного компонента в нижней части окна приложения (выберите меню \textit{Вид} $\rightarrow$ \textit{Консоль вывода}, если она не видна) выводит записи из этого файла журнала.

\begin{redremark}
Пожалуйста, используйте содержимое файла журнала, сообщая о любых ошибках приложения.
\end{redremark}

\begin{figure}[H]
	\centering
	\includegraphics[width = 0.9\linewidth]{img/output_console.png}
\end{figure}

 В \textit{Настройках} можно задать уровень вывода для ограничения содержащейся в логе информации.

\anonsubsubsection{Панель Подключения}

На боковой панели \textit{Подключения} отображены имена всех ранее устанавливаемых соединений (см. раздел \hyperref[sec:connections]{«Подключение к базе данных»}). 

\begin{figure}[H]
	\centering
	\includegraphics[width = 0.5\linewidth]{img/connections.png}
\end{figure}

После подключения к базе данных на этой панели появляется структура дерева, узлы которого представляют собой объекты базы данных. В скобках указано количество объектов каждого вида.

Щелчок правой кнопкой мыши на любом узле вызовет всплывающее меню с доступными действиями. Оно включает в себя перезагрузку выбранного узла, создание, редактирование и удаление объекта и многое другое.

Двойной клик по узлу или нажатие на Enter  откроет вкладку на главной панели с детальным описанием каждого объекта (см. раздел \hyperref[sec:objects]{«Объекты»}). Например, двойной клик (или Enter) по узлу с названием конкретной процедуры отобразит таблицу с ее входными и выходными параметрами, их типами, тело процедуры и SQL команду создания этой процедуры.

\begin{figure}[H]
	\centering
	\includegraphics[width = 0.7\linewidth]{img/procedure_view.png}
\end{figure}

\anonsubsubsection{Ключевые слова}

Компонент \textit{Ключевые слова} отображает список ключевых слов SQL, которые в настоящее время поддерживаются системой. Он включает ключевые слова стандарта SQL92 в дополнение к ключевым словам, определенным пользователями и специфичным для базы данных, которые извлекаются из источника данных при подключении.

Тип ключевого слова обозначается своим значком, а подсказка мыши указывает откуда оно было получено - из SQL92, определено пользователем или зарегистрировано в базе данных.

Ключевые слова можно искать с помощью текстового поля в верхней части компонента.

\begin{figure}[H]
	\centering
	\includegraphics[width = 0.3\linewidth]{img/keywords.png}
\end{figure}

\anonsubsubsection{Драйверы}

Панель \textit{Драйверы} перечисляет все установленные в настоящее время драйверы JDBC в приложении. Дополнительную информацию см. в разделе \hyperref[sec:drivers]{«Драйверы»}.

\begin{figure}[H]
	\centering
	\includegraphics[width = 0.3\linewidth]{img/drivers.png}
\end{figure}

\anonsubsubsection{Строка состояния}

\textit{Строка состояния} это полоса в нижней части приложения с информацией о количестве подключений, доступных обновлениях, версии JDK, состоянии памяти.

\begin{figure}[H]
	\centering
	\includegraphics[width = 1\linewidth]{img/status_bar.png}
\end{figure}

\hyperref[sec:toolbar]{Панель инструментов} в верхней части окна приложения также настраивается через меню \textit{Вид} и в \hyperref[sec:pref_toolbar]{настройках} приложения.

\begin{figure}[H]
	\centering
	\includegraphics[width = 1\linewidth]{img/toolbar.png}
\end{figure}

\anonsubsubsection[Статистика БД]{\hyperref[sec:gstat]{Статистика БД}}

Отображает статистику базы данных, а также позволяет сравнить результаты анализа двух баз.

\begin{figure}[H]
	\centering
	\includegraphics[width = 0.7\linewidth]{img/gstat_table.png}
	\caption{Статистика по таблицам}
\end{figure}

\newpage
\section{Панель инструментов}\label{sec:toolbar}

Red Expert включает 6 основных групп панелей инструментов в верхней части окна приложения. Отдельные группы могут перемещаться горизонтально или вертикально вдоль верхней части окна приложения. Каждую группу инструментов можно скрыть или сделать видимой через меню \textit{Вид} или в настройках приложения.

Далее следует описание каждой панели инструментов и связанных с ней кнопок и действий. 

\begin{longtable}[c]{|m{5mm}|m{9cm}|>{\ttfamily}m{4cm}|}
	\caption{Панель инструментов для работы с файлами}\\
	\hline
	&
	\centering\bfseries Описание &
	\centering\arraybslash\normalfont\bfseries Горячие клавиши\\\hline
	\endfirsthead
	\hline
	&
	\centering\bfseries Описание &
	\centering\arraybslash\normalfont\bfseries Горячие клавиши\\\hline
	\endhead
	\begin{tikzpicture}
	\pgftext{\includegraphics[width=11pt]{img/new16.png}} at (0pt,0pt)
	\end{tikzpicture} & Открыть редактор запросов & Shift+F12\\
	\hline
	\begin{tikzpicture}
	\pgftext{\includegraphics[width=14pt]{img/open16.png}} at (0pt,0pt)
	\end{tikzpicture} & Открыть файл & Ctrl+L\\
	\hline
	\begin{tikzpicture}
	\pgftext{\includegraphics[width=11pt]{img/save16.png}} at (0pt,0pt)
	\end{tikzpicture} & Сохранить файл & Ctrl+S\\
	\hline
	\begin{tikzpicture}
	\pgftext{\includegraphics[width=12pt]{img/print16.png}} at (0pt,0pt)
	\end{tikzpicture} & Печать & Ctrl+P\\
	\hline
	\begin{tikzpicture}
	\pgftext{\includegraphics[width=12pt]{img/PrintPreview16.png}} at (0pt,0pt)
	\end{tikzpicture} & Предварительный просмотр перед печатью & \\
	\hline
	\begin{tikzpicture}
	\pgftext{\includegraphics[width=11pt]{img/SaveAs16.png}} at (0pt,0pt)
	\end{tikzpicture} & Сохранить как & Ctrl+Shift+S\\
	\hline
	\begin{tikzpicture}
	\pgftext{\includegraphics[width=12pt]{img/PageSetup16.png}} at (0pt,0pt)
	\end{tikzpicture} & Свойства печати &\\
	\hline
\end{longtable}

\begin{longtable}[c]{|m{5mm}|m{9cm}|>{\ttfamily}m{4cm}|}
	\caption{Панель инструментов для редактирования}\\
	\hline
	&
	\centering\bfseries Описание &
	\centering\arraybslash\normalfont\bfseries Горячие клавиши\\\hline
	\endfirsthead
	\hline
	&
	\centering\bfseries Описание &
	\centering\arraybslash\normalfont\bfseries Горячие клавиши\\\hline
	\endhead
	\begin{tikzpicture}
	\pgftext{\includegraphics[width=12pt]{img/cut16.png}} at (0pt,0pt)
	\end{tikzpicture} & Вырезать & Ctrl+X\\
	\hline
	\begin{tikzpicture}
	\pgftext{\includegraphics[width=12pt]{img/Copy16.png}} at (0pt,0pt)
	\end{tikzpicture} & Копировать & Ctrl+C\\
	\hline
	\begin{tikzpicture}
	\pgftext{\includegraphics[width=12pt]{img/Paste16.png}} at (0pt,0pt)
	\end{tikzpicture} & Вставить & Ctrl+V\\
	\hline
	\begin{tikzpicture}
	\pgftext{\includegraphics[width=12pt]{img/Undo16.png}} at (0pt,0pt)
	\end{tikzpicture} & Отмена действий & Ctrl+Z\\
	\hline
	\begin{tikzpicture}
	\pgftext{\includegraphics[width=12pt]{img/Redo16.png}} at (0pt,0pt)
	\end{tikzpicture} & Возврат отмененного действия & Ctrl+Y\\
	\hline
\end{longtable}	

\begin{longtable}[c]{|m{5mm}|m{9cm}|>{\ttfamily}m{4cm}|}
	\caption{Панель инструментов для работы с базой данных}\\
	\hline
	&
	\centering\bfseries Описание &
	\centering\arraybslash\normalfont\bfseries Горячие клавиши\\\hline
	\endfirsthead
	\hline
	&
	\centering\bfseries Описание &
	\centering\arraybslash\normalfont\bfseries Горячие клавиши\\\hline
	\endhead
	\begin{tikzpicture}
	\pgftext{\includegraphics[width=12pt]{img/DBImage16.png}} at (0pt,0pt)
	\end{tikzpicture} & Создать подключение & Ctrl+Shift+N\\
	\hline
	\begin{tikzpicture}
	\pgftext{\includegraphics[width=12pt]{img/Edit16.png}} at (0pt,0pt)
	\end{tikzpicture} & Открыть редактор запросов & \\
	\hline
	\begin{tikzpicture}
	\pgftext{\includegraphics[width=12pt]{img/create_database16.png}} at (0pt,0pt)
	\end{tikzpicture} & Создать базу данных & \\
	\hline
	\begin{tikzpicture}
	\pgftext{\includegraphics[width=12pt]{img/CreateScripts16.png}} at (0pt,0pt)
	\end{tikzpicture} & Извлечь метаданные БД в SQL скрипт & \\
	\hline
	\begin{tikzpicture}
	\pgftext{\includegraphics[width=10pt]{img/Procedure16.png}} at (0pt,0pt)
	\end{tikzpicture} & Выполнить хранимые объекты & \\
	\hline
	\begin{tikzpicture}
	\pgftext{\includegraphics[width=10pt]{img/user_manager_16.png}} at (0pt,0pt)
	\end{tikzpicture} & Менеджер пользователей & \\
	\hline
	\begin{tikzpicture}
	\pgftext{\includegraphics[width=12pt]{img/grant_manager_16.png}} at (0pt,0pt)
	\end{tikzpicture} & Менеджер прав & \\
	\hline
	\begin{tikzpicture}
	\pgftext{\includegraphics[width=12pt]{img/CompareDataTypes16.png}} at (0pt,0pt)
	\end{tikzpicture} & Сравнить типы данных & \\
	\hline
	\begin{tikzpicture}
	\pgftext{\includegraphics[width=12pt]{img/Keywords16.png}} at (0pt,0pt)
	\end{tikzpicture} & Пользовательские ключевые слова & \\
	\hline
	\begin{tikzpicture}
	\pgftext{\includegraphics[width=12pt]{img/SystemConsole16.png}} at (0pt,0pt)
	\end{tikzpicture} & Окно системной консоли & \\
	\hline
\end{longtable}	

\begin{longtable}[c]{|m{5mm}|m{9cm}|>{\ttfamily}m{4cm}|}
	\caption{Панель инструментов для поиска}\\
	\hline
	&
	\centering\bfseries Описание &
	\centering\arraybslash\normalfont\bfseries Горячие клавиши\\\hline
	\endfirsthead
	\hline
	&
	\centering\bfseries Описание &
	\centering\arraybslash\normalfont\bfseries Горячие клавиши\\\hline
	\endhead
	\begin{tikzpicture}
	\pgftext{\includegraphics[width=12pt]{img/Find16.png}} at (0pt,0pt)
	\end{tikzpicture} & Найти & Ctrl+F\\
	\hline
	\begin{tikzpicture}
	\pgftext{\includegraphics[width=12pt]{img/FindAgain16.png}} at (0pt,0pt)
	\end{tikzpicture} & Найти далее & F3\\
	\hline
	\begin{tikzpicture}
	\pgftext{\includegraphics[width=12pt]{img/Replace16.png}} at (0pt,0pt)
	\end{tikzpicture} & Заменить & Ctrl+R\\
	\hline
	\begin{tikzpicture}
	\pgftext{\includegraphics[width=12pt]{img/GoTo16.png}} at (0pt,0pt)
	\end{tikzpicture} & Перейти к ... & Ctrl+G\\
	\hline
	\begin{tikzpicture}
	\pgftext{\includegraphics[width=12pt]{img/FindInFiles16.png}} at (0pt,0pt)
	\end{tikzpicture} & Найти файлы & \\
	\hline
\end{longtable}	

\begin{longtable}[c]{|m{5mm}|m{9cm}|>{\ttfamily}m{4cm}|}
	\caption{Системные настройки}\\
	\hline
	&
	\centering\bfseries Описание &
	\centering\arraybslash\normalfont\bfseries Горячие клавиши\\\hline
	\endfirsthead
	\hline
	&
	\centering\bfseries Описание &
	\centering\arraybslash\normalfont\bfseries Горячие клавиши\\\hline
	\endhead
	\begin{tikzpicture}
	\pgftext{\includegraphics[width=12pt]{img/Preferences16.png}} at (0pt,0pt)
	\end{tikzpicture} & Настройки & \\
	\hline
	\begin{tikzpicture}
	\pgftext{\includegraphics[width=12pt]{img/Help16.png}} at (0pt,0pt)
	\end{tikzpicture} & Разделы справки & F1\\
	\hline
	\begin{tikzpicture}
	\pgftext{\includegraphics[width=12pt]{img/About16.png}} at (0pt,0pt)
	\end{tikzpicture} & О программе & \\
	\hline
	\begin{tikzpicture}
	\pgftext{\includegraphics[width=12pt]{img/WebComponent16.png}} at (0pt,0pt)
	\end{tikzpicture} & Посетить red-soft.ru & \\
	\hline
	\begin{tikzpicture}
	\pgftext{\includegraphics[width=12pt]{img/ResourceHeap16.png}} at (0pt,0pt)
	\end{tikzpicture} & Состояние памяти & \\
	\hline
\end{longtable}
\newpage

\section{Драйверы}\label{sec:drivers}

Red Expert использует информацию, доступную из установленных JDBC драйверов базы данных, для генерации большей части ее содержимого.

По умолчанию в Red Expert уже установлена библиотека драйвера Jaybird 3 Driver, которая позволяет работать с базами данных Firebird и Ред Базой Данных.

\anonsubsection{Установка нового драйвера}\label{sec:new_driver}

\begin{redremark}
	Для получения информации о загрузке и использовании конкретных JDBC драйверов обратитесь к документации поставщиков баз данных.
\end{redremark}

\begin{enumerate}[leftmargin=15pt]
	\item Откройте панель \textit{Драйверы} в пункте меню \textit{Инструменты}.
	\item Откройте вкладку «Драйверы» на центральной панели (если вкладка не видна, кликните мышью по узлу «JDBC Drivers» панели  \textit{Драйверы}).
	\item Нажмите кнопку «New driver» на вкладке «Драйверы» или на кнопку 
	\begin{tikzpicture}
	\pgftext{\includegraphics[width=14pt]{img/NewJDBCDriver16.png}} at (0pt,0pt)
	\end{tikzpicture}
	на панели инструментов компонента \textit{Драйверы}.
	\item Заполните соответствующие поля, как показано ниже:
	\begin{longtable}[r]{|>{\bfseries}m{3.5cm}|m{11.4cm}|}
		\hline
		Имя драйвера & Имя драйвера для идентификации\\\hline
		Описание & Краткое описание этого драйвера \\\hline
		База данных & Выберите СУБД, для которой этот драйвер предназначен\\\hline
		JDBC URL & Шаблон URL-адреса для этого JDBC драйвера. Например:
		\ttt{jdbc:firebirdsql://[host]:[port]/[source]} \\\hline
		Путь & Путь к jar-файлу JDBC-драйвера \\\hline
		Имя класса & Имя класса JDBC-драйвера. Выберите кнопку поиска, если имя неизвестно, и система сканирует файл jar, введенный в поле пути, для поиска имени класса драйвера\\\hline
	\end{longtable}
	\item Теперь драйвер готов к использованию и появится в списке доступных драйверов при подключении к базе данных в качестве имени, указанного выше.
\end{enumerate}

\newpage

\section{Создание базы данных}\label{sec:create_db}
\begin{enumerate}[leftmargin=15pt]
	\item Откройте вкладку «Create Database» на центральной панели. Для этого выберите пункт меню \textit{База данных} $\rightarrow$ \textit{Создать базу данных}. Либо нажмите кнопку 
	\begin{tikzpicture}
	\pgftext{\includegraphics[width=14pt]{img/create_database16.png}} at (0pt,0pt)
	\end{tikzpicture}
	на главной панели инструментов.
	\item Заполните соответствующие поля, как показано ниже:
	\begin{longtable}[r]{|>{\bfseries}m{3.9cm}|m{11cm}|}
		\hline
		JDBC драйвер & Выберите JDBC драйвер из выпадающего списка для создания новой базы данных. Для Ред Базы Данных и Firebird рекомендуется JDBC драйвер Jaybird 3 \\\hline
		Имя подключения & Имя подключения к этой базе данных \\\hline
		Имя сервера & Имя хоста сервера базы данных или IP-адрес\\\hline
		Порт & Номер порта подключения к базе данных \\\hline
		Файл базы данных & Путь к файлу базы данных или ее алиас\\\hline
		Имя пользователя & Логин пользователя, который в последствии станет владельцем базы данных\\\hline
		Пароль & Пароль пользователя \\\hline
		Сохранить пароль & Следует ли сохранять значение пароля с другой информацией о соединении\\\hline
		Зашифровать 
		
		пароль 	& Следует ли шифровать пароль при хранении, чтобы избежать открытого просмотра сохраненных паролей\\\hline
		Кодировка & Задает набор символов по умолчанию для строковых
		(символьных) данных для всей базы данных.  \\\hline
		Размер страницы & Размер страницы базы данных в байтах. Допустимыми значениями являются 4096, 8192 и 16384. \\\hline
	\end{longtable}
	
	\item Нажмите кнопку «Создать» для создания новой базы данных с введенные параметрами. 
\end{enumerate}	

\begin{redremark}
При попытке создания базы данных может возникнуть ошибка шифрования сетевого соединения (wire encryption). Для поддержки шифрования необходима версия Java не ниже 1.8.0\_161 или установленное JCE дополнение. В противном случае следует изменить параметр \ttt{firebird.conf}: \ttt{WireCrypt = Disabled}.
\end{redremark}

\anonsubsection{Импорт данных из файла}\label{sec:import_data_from_file}

\begin{enumerate}
	\item Откройте вкладку "Импорт данных" на центральной панели. Для этого выберите пунк меню \textit{База данных} $\rightarrow$ \textit{Импорт данных}. 
	\begin{figure}[H]
		\centering
		\includegraphics[width = 0.6\linewidth]{img/OpenImportData.png}
	\end{figure}
	Откроется вкладка:
	\begin{figure}[H]
		\centering
		\includegraphics[width = 0.95\linewidth]{img/ImportData.png}
	\end{figure}
	\item Заполните соответствующие поля, как показано ниже:
	\begin{longtable}[r]{|>{\bfseries}m{3.9cm}|m{11cm}|}
		\hline
		Импортировать из соединения & Позволяет выбрать базу данных в качестве источника.\\\hline
		Файл данных & Путь к файлу-источнику данных. Поддерживается импорт из XLSX, XML и CSV файлов.\\\hline
		Файл блобов & Путь к файлу-источнику с блобами (с расширением \ttt{.lob}).\\\hline
		Целевая БД & База данных, в которую будет производиться импорт.\\\hline
		Целевая таблица & Таблица, в которую будет производиться импорт.\\\hline
		Разделитель & Тип разделителя данных в CSV файлах.\\\hline
		Номер страницы & Для \ttt{XLSX} файлов можно выбрасть страницу excel-файла, из которой нужно импортировать данные.\\\hline
		Первая строка заголовки столбцов источника & Следует ли сделать значения в первой строке выбранного файла заголовками столбцов в табличном представлении. В таком случае при импорте первая строка не будет учитываться.\\\hline
		Первая импортируемая строка & Строка, начиная с которой будут импортированы данные.\\\hline
		Последняя импортируемая строка & Строка, поcле которой импорт данных будет прекращён. Строки, непопадающие в диапазон между первой и последней импортируемой строкой, будут проигнорированы.\\\hline
		Шаг фиксации & Количество записей, после которого будет произведено подтверждение транзакии и сохранение данных в таблице.\\\hline
		Очистить таблицу перед вставкой & Следует ли очистить целевую таблицу от данных, которые в ней были до импорта.\\\hline
		Столбец источника & Имя столбца, из которого нужно импортировать данные.  Кнопка \ttt{Сопоставить} автоматически распределяет импортируемые столбцы к целевым столбцам по их именам.\\\hline
		Параметры & Импортировать \ttt{BLOB} как файл или как текст.\\\hline
	\end{longtable}
	Пример заполнения полей:
	\begin{figure}[H]
		\centering
		\includegraphics[width = 1\linewidth]{img/ExampleImportData.png}
	\end{figure}
	\item Нажмите на кнопку "Импортировать Данные".
\end{enumerate}

В процессе импортирования будет выведено диалоговое окно. Нажатие на кнопку "Отмена" остановит процесс импорта данных. Все изменения, произошедшие до остановки импортирования, сохранятся в таблице.

\anonsubsection{Дополнительные параметры}

Дополнительные параметры соединения можно ввести, выбрав вкладку «Расширенные». На ней располагается таблица с двумя столбцами. Левый соответствует параметру подключения, а правый -- его значению. Обратитесь к соответствующей документации JDBC драйвера, чтобы узнать, какие (если есть) дополнительные параметры подключения к базе данных могут быть установлены.

Для драйвера Jaybird 3 в таблице \hyperlink{tab:01}{ниже} перечислены основные параметры подключения.

Под таблицей предлагается выбрать уровень изоляции транзакций. Более подробно о них рассказано \hyperlink{tab:isollevel}{далее}.

\newpage
\section{Подключение к базе данных}\label{sec:connections}

Red Expert позволяет одновременно использовать несколько подключений к базе данных. Информация о подключении может быть получена через вкладку «Браузер базы данных».

\anonsubsection{Создание нового подключения}

\begin{enumerate}[leftmargin=15pt]
	\item Убедитесь, что открыта панель \textit{Подключения} (если панель не открыта, установите галочку в пункте меню \textit{Вид} $\rightarrow$ \textit{Подключения}).
	\item Откройте вкладку «Браузер баз данных» на центральной панели (если вкладка не видна, кликните мышью по узлу «Подключения к базе данных» панели \textit{Подключения}).
	\item Нажмите кнопку «Новое подключение» на вкладке «Браузер баз данных» или кнопку 
	\begin{tikzpicture}
	\pgftext{\includegraphics[width=14pt]{img/NewConnection16.png}} at (0pt,0pt)
	\end{tikzpicture}
	на панели инструментов компонента \textit{Подключения}.
	\item Заполните соответствующие поля, как показано ниже:
	\begin{longtable}[r]{|>{\bfseries}m{3.9cm}|m{11cm}|}
		\hline
		Состояние & Индикация, было ли установлено выбранное соединение - подключено или не подключено\\\hline
		JDBC драйвер & Выберите JDBC драйвер из выпадающего списка для этого подключения. Для Ред Базы Данных и Firebird рекомендуется JDBC драйвер Jaybird 3\\\hline
		Имя подключения & Имя подключения \\\hline
		Параметры
		
		подключения & Выпадающий список предлагает выбрать два способа подключения: 
		\begin{itemize}[leftmargin=10pt]
			\item Стандартно: отдельно указываются параметры подключения -- имя сервера, порт, файл БД, способ аутентификации, имя пользователя, пароль, кодировка и роль пользователя. 
			\item JDBC: по URL-адресу выбранного JDBC драйвера, например: \ttt{jdbc:firebirdsql://localhost:3050/employee?user=SYSDBA\& password=masterkey}
		\end{itemize}\\\hline
		Имя сервера & Имя хоста сервера базы данных или IP-адрес\\\hline
		Порт & Номер порта подключения к базе данных \\\hline
		Файл базы данных & Путь к файлу базы данных или ее алиас\\\hline
		Аутентификация & Способ аутентификации клиента:
		\begin{itemize}[leftmargin=10pt]
			\item Базовая: это традиционная (native для Ред Базы Данных 2.6; Legacy\_Auth для Ред Базы Данных 3.0) или безопасная парольная (Srp для Ред Базы Данных 3.0)
			\item Gss: доверенная (для Ред Базы Данных 3.0)
			\item Multifactor: многофакторная (для Ред Базы Данных)
		\end{itemize} 
		Если нужная аутентификация на сервере не включена или ее не существует, то будет выдано сообщение об ошибке о неверном логине-пароле. \\\hline
		Имя пользователя & Логин пользователя для подключения\\\hline
		Пароль & Пароль пользователя для подключения \\\hline
		Сохранить пароль & Следует ли сохранять значение пароля с другой информацией о соединении\\\hline
		Зашифровать 
		
		пароль 	& Следует ли шифровать пароль при хранении, чтобы избежать открытого просмотра сохраненных паролей\\\hline
		Кодировка & Набор символов для подключения \\\hline
		Роль & Роль подключающегося пользователя в текущей базе данных\\\hline
		JDBC URL & URL-адрес выбранного JDBC драйвера\\\hline
	\end{longtable}
	\item Нажмите кнопку «Подключиться», чтобы приложение попыталось установить соединение, используя введенные значения.
\end{enumerate}	

\begin{redremark}
	При попытке подключения может возникнуть ошибка шифрования сетевого соединения (wire encryption). Для поддержки шифрования необходима версия Java не ниже 1.8.0\_161 или установленное JCE дополнение. В противном случае следует изменить параметр \ttt{firebird.conf}: \ttt{WireCrypt = Disabled}.
\end{redremark}


\anonsubsubsection{Расширенные параметры подключения}

Дополнительные свойства соединения можно ввести, выбрав вкладку «Расширенные». На ней можно увидеть таблицу с двумя столбцами. Левый соответствует параметру подключения, а правый -- его значению. Обратитесь к соответствующей документации JDBC драйвера, чтобы узнать, какие (если есть) дополнительные параметры подключения к базе данных могут быть установлены.

В таблице перечислены некоторый параметры для драйвера Jaybird 3: \hypertarget{tab:01}{}
\begin{longtable}[c]{|>{\ttfamily}m{4.2cm}|>{\ttfamily\centering}m{1.7cm}|m{9.5cm}|}
	\hline
	\centering\normalfont\bfseries Параметр &
	\centering\normalfont\bfseries Тип значения&
	\centering\arraybslash\bfseries Описание\\\hline
	\endfirsthead
	\hline
	\centering\normalfont\bfseries Параметр &
	\centering\normalfont\bfseries Тип значения &
	\centering\arraybslash\bfseries Описание\\\hline
	\endhead
	isc\_dpb\_user\_name & string & Имя подключающегося пользователя\\\hline
	isc\_dpb\_password & string & Пароль пользователя\\\hline
	isc\_dpb\_sql\_role\_name & string & Имя роли\\\hline
	isc\_dpb\_sql\_dialect & byte & SQL диалект \\\hline
	isc\_dpb\_process\_id & int & ID процесса \\\hline
	isc\_dpb\_process\_name & string & Имя процесса \\\hline
	isc\_dpb\_lc\_ctype & string & Кодировка символов соединения. Этот параметр сообщает серверу базы данных в какой кодировке клиент ждет строки и текст.\\\hline
	isc\_dpb\_connect\_timeout & int & Время ожидания подключения (в секундах)\\\hline
	isc\_dpb\_gss & -- & Использовать доверительную аутентификацию Gss\\\hline
	isc\_dpb\_num\_buffers & int & Количество страниц базы данных, которые будут кэшироваться\\\hline
	isc\_dpb\_set\_db\_readonly & boolean & Установить базу данных в режим только для чтения\\\hline
	isc\_dpb\_set\_db\_charset & string & Установить набор символов для базы данных \\\hline
	isc\_dpb\_trusted\_auth & -- & Информирование, что будет использоваться доверительная аутентификация\\\hline
	isc\_dpb\_multi\_factor\_ auth & boolean &  Информирование, что будет использоваться многофакторная аутентификация\\\hline
	isc\_dpb\_utf8\_filename  & -- & Информирование, что имя файла передается в формате UTF-8\\\hline
	isc\_dpb\_certificate & string & Алиас сертификата клиента при многофакторном подключении\\\hline
	isc\_dpb\_verify\_server & -- & Флаг для проверки сертификата сервера (многофакторная аутентификация) \\\hline
	isc\_dpb\_repository\_pin & string & Пин код для контейнера с сертификатом (многофакторная аутентификация)\\\hline
	isc\_dpb\_wire\_crypt\_
	level & string & Параметр устанавливает, следует ли шифровать сетевое соединение. Он может принимать три	возможных значения: \ttt{REQUIRED, ENABLED, DISABLED} и \ttt{DEFAULT}. По умолчанию шифрование включено (\ttt{Enabled}).\\\hline
\end{longtable} 

Ниже под таблицей предлагается выбрать уровень изоляции транзакций. Различные уровни изоляции транзакций определяют поведение данного клиентского приложения, запустившего эту транзакцию, по отношению к другим параллельным процессам, выполняющимся на любом компьютере локальной сети, одновременно выполняющих чтение и/или изменение в той же базе данных, что и текущий процесс. 

Для транзакций существует несколько типов чтения:
\begin{itemize}
	\item \textit{Грязное чтение} (\ttt{dirty reads}) происходит, когда транзакциям разрешено видеть несохраненные изменения данных. Иными словами, изменения, сделанные в одной транзакции, видны вне ее до того, как она была сохранена. Если изменения не будут сохранены, то, вероятно, другие транзакции выполняли работу на основе некорректных данных;
	\item \textit{Неповторяющееся чтение} (\ttt{non-repeatable reads}) происходит, когда транзакция А читает строку, транзакция Б изменяет эту строку, транзакция А читает ту же строку и получает обновленные данные;
	\item \textit{Фантомное чтение} (\ttt{phantom reads}) происходит, когда транзакция А считывает все строки, удовлетворяющие \ttt{WHERE}-условию, транзакция Б вставляет новую или удаляет одну из строк, которая удовлетворяет этому условию, транзакция А еще раз считывает все строки, удовлетворяющие \ttt{WHERE}-условию, уже вместе с новой строкой или недосчитавшись старой.
\end{itemize}

Уровни изоляции транзакций определены в виде констант по возрастанию уровня ограничения:\hypertarget{tab:isollevel}{}
\begin{longtable}[c]{|>{\ttfamily}m{5.2cm}|m{10.2cm}|}
	\hline
	\centering\normalfont\bfseries Уровень изоляции &
	\centering\arraybslash\bfseries Описание\\\hline
	\endfirsthead
	\hline
	\centering\normalfont\bfseries Уровень изоляции &
	\centering\arraybslash\bfseries Описание\\\hline
	\endhead
	Стандартная база данных & Уровень изоляции по умолчанию, который установлен для данной базы данных. Для Ред Базы Данных и Firebird это \ttt{READ\_COMMITED}\\\hline
	TRANSACTION\_NONE & Информирует о том, что драйвер не поддерживает транзакции\\\hline
	TRANSACTION\_READ\_UNCOMMITED & Позволяет транзакциям видеть несохраненные изменения данных, что разрешает грязное, неповторяющееся и фантомное чтения\\\hline
	TRANSACTION\_READ\_COMMITED & Означает, что любое изменение, сделанное в транзакции, не видно вне неё, пока она не сохранена. Это предотвращает грязное чтение, но разрешает неповторяющееся и фантомное\\\hline
	TRANSACTION\_REPEATABLE\_READ & Запрещает грязное и неповторяющееся, но фантомное чтение разрешено\\\hline
	TRANSACTION\_SERIALIZABLE & Определяет, что грязное, неповторяющееся и фантомное чтения запрещены\\\hline
\end{longtable} 	


\anonsubsubsection{SSH туннель}

Есть удобная возможность подключиться к базе данных через SSH-туннель. Как и при любом SSH-соединении, весь трафик между вами и БД будет шифроваться.

Для этого введите параметры подключения на вкладке «Базовые», переключитесь на вкладку «SSH Тоннель» и заполните параметры для SSH-соединения. Имя хоста переносится с вкладки «Базовые».

\anonsubsection{Импорт подключений}
Для тех, у кого уже есть настроенные в IBExpert подключения, разработана удобная опция импорта подключений в Red Expert.

\begin{enumerate}[leftmargin=15pt]
	\item Запустите IBExpert и откройте пункт меню \textit{Options} $\rightarrow$ \textit{Enviroment Options...}. Откроется дополнительное окно:
	\begin{figure}[H]
		\centering
		\includegraphics[width = 0.55\linewidth]{img/IBExpert_env_opt.png}
	\end{figure}
	\item Выберите пункт \textit{IBExpert User Database}.
	\item В поле «User Database Connection String» прописан путь к пользовательской базе данных, которая хранит регистрационную информацию.
	\item Далее запустите Red Expert и подключитесь к той пользовательской базе данных.
	\item Выберите пункт меню \textit{База данных} $\rightarrow$ \textit{Подключения} $\rightarrow$ \textit{Импорт подключений} $\rightarrow$ \textit{Импорт подключений из базы данных}. Откроется окно:
	\begin{figure}[H]
		\centering
		\includegraphics[width = 0.35\linewidth]{img/import_connections.png}
	\end{figure}
	\item Выберите нужное подключение и нажмите \textit{OK}.
	\item На панели «Подключения» Вы найдете импортированные подключения.
\end{enumerate}	

\anonsubsection{Закрытие соединения}

Подключение можно разорвать в любой момент, нажав на кнопку «Отключиться» на вкладке «Браузер баз данных». Есть и другой способ. На боковой панели \textit{Подключения} щелкните правой кнопкой мыши на нужное соединение и выберите пункт всплывающего меню \textit{Отключиться от базы данных <имя подключения>}.

\anonsubsection{Свойства базы данных}

После подключения к базе данных становится активной вкладка «Свойства базы данных». На ней располагается таблица с информацией о возможностях конкретной СУБД в сочетании с конкретным JDBC драйвером. Описание отдельных свойств см.~\hyperref[app:1]{Приложение~\ref{app:1}}

\anonsubsection{Ключевые слова SQL}

После подключения к базе данных становится активной вкладка «Ключевые слова SQL». На ней в одной из таблиц перечислены ключевые слова, определенные в стандарте SQL92, в другой -- определённые для конкретной СУБД.

Ключевые слова нельзя использовать в качестве имен объектов базы данных, переменных и параметров. 

\newpage

\section{Редактор запросов}\label{sec:query_editor}

Чтобы открыть редактор запросов в центральной панели, выберите пункт меню \textit{Инструменты} $\rightarrow$ \textit{Редактор запросов} или нажмите кнопку \begin{tikzpicture}
\pgftext{\includegraphics[width=12pt]{img/Edit16.png}} at (0pt,0pt)
\end{tikzpicture} 
на главной панели инструментов. Он открывается в виде вкладки с названием \textit{Query Editor -- script<N>.sql}. Файл с запросами к базе данных \ttt{script<N>.sql} можно сохранить (по умолчанию файл сохраняется в \ttt{\$HOME\textbackslash .redexpert\textbackslash<buildNumber>\textbackslash QueryEditor}). 

Редактор запросов представляет собой настраиваемый инструмент просмотра и выполнения операторов SQL. В любой момент времени может быть открыто любое количество редакторов.

Редактор поддерживает следующие функции:
\begin{itemize}
	\item Настраиваемая подсветка синтаксиса SQL;
	\item Подсказки ключевых слов и имен объектов базы данных;
	\item Выполнение нескольких запросов;
	\item Выполнение и отображение нескольких запросов с множеством результатов (\hyperref[sec:resultset]{Result Set});
	\item Поддержка \hyperref[sec:param_queries]{параметризованных запросов};
	\item Полная поддержка печати;
	\item Управление транзакциями;
	\item Функции текстового редактора стиля IDE - поиск, замена, вставка и т.д.;
	\item Экспорт результатов;
	\item Поддержка нескольких открытых соединений;
	\item Поисковая исполняемая \hyperref[sec:query_history]{история} SQL-запросов;
	\item Быстрый переход из редактора к просмотру объекта базы данных двойным кликом или нажатием \ttt{CTRL + Левая клавиша мыши} по имени объекта;
	\item Выбор уровня изоляции транзакции.
\end{itemize}

\anonsubsection{Панель инструментов}

Следующие команды доступны на панели инструментов редактора запросов:

\begin{longtable}[c]{|m{5mm}|m{10.6cm}|>{\ttfamily}m{4cm}|}
	\hline
	&
	\centering\bfseries Описание &
	\centering\arraybslash\normalfont\bfseries Горячие клавиши\\\hline
	\endfirsthead
	\hline
	&
	\centering\bfseries Описание &
	\centering\arraybslash\normalfont\bfseries Горячие клавиши\\\hline
	\endhead
	\begin{tikzpicture}
	\pgftext{\includegraphics[width=12pt]{img/Execute16.png}} at (0pt,0pt)
	\end{tikzpicture} & Выполнить все операторы в редакторе & Ctrl+Shift+Enter\\\hline
	\begin{tikzpicture}
	\pgftext{\includegraphics[width=12pt]{img/ExecuteSelection16.png}} at (0pt,0pt)
	\end{tikzpicture} & Выполнить любой выделенный текст в редакторе & F9\\\hline
 	\begin{tikzpicture}
 	\pgftext{\includegraphics[width=12pt]{img/ExecuteSqlScript16.png}} at (0pt,0pt)
 	\end{tikzpicture}& Выполнить SQL скрипт & F10\\\hline
 	\begin{tikzpicture}
 	\pgftext{\includegraphics[width=12pt]{img/ExecuteInProfiler.png}} at (0pt,0pt)
 	\end{tikzpicture} & Выполнить в профайлере & \\\hline
 	\begin{tikzpicture}
 	\pgftext{\includegraphics[width=12pt]{img/PrintPlan.png}} at (0pt,0pt)
 	\end{tikzpicture}& Печать плана запроса & Ctrl+Shift+P\\\hline
	\begin{tikzpicture}
	\pgftext{\includegraphics[width=12pt]{img/plan.png}} at (0pt,0pt)
	\end{tikzpicture} & Печать детального плана запроса & Ctrl+Alt+P\\\hline
	\begin{tikzpicture}
	\pgftext{\includegraphics[width=12pt]{img/Stop16.png}} at (0pt,0pt)
	\end{tikzpicture} & Отмена выполнения текущего оператора & \\\hline
	\begin{tikzpicture}
	\pgftext{\includegraphics[width=12pt]{img/ClearOutput16.png}} at (0pt,0pt)
	\end{tikzpicture} & Очистить окно вывода (окно Output) & Ctrl+Shift+Backspace\\\hline
	\begin{tikzpicture}
	\pgftext{\includegraphics[width=12pt]{img/History16.png}} at (0pt,0pt)
	\end{tikzpicture} & Показать диалоговое окно истории выполненных запросов & Ctrl+Shift+H\\\hline
	\begin{tikzpicture}
	\pgftext{\includegraphics[width=12pt]{img/Bookmarks16.png}} at (0pt,0pt)
	\end{tikzpicture} & Закладки для хранения запросов. Для удобного поиска важных запросов  & Ctrl+B\\\hline
	\begin{tikzpicture}
	\pgftext{\includegraphics[width=12pt]{img/Shortcut16.png}} at (0pt,0pt)
	\end{tikzpicture} & Сокращения запросов. Появляется диалоговое окно, в котором можно найти или создать новое сокращение для определенных запросов к базе данных & \\\hline
	\begin{tikzpicture}
	\pgftext{\includegraphics[width=12pt]{img/Previous16.png}} at (0pt,0pt)
	\end{tikzpicture} & Показать в редакторе предыдущий выполненный запрос  & Ctrl+Shift+Down\\\hline
	\begin{tikzpicture}
	\pgftext{\includegraphics[width=12pt]{img/Forward16.png}} at (0pt,0pt)
	\end{tikzpicture} & Показать в редакторе следующий выполненный запрос  & Ctrl+Shift+Up\\\hline
	\begin{tikzpicture}
	\pgftext{\includegraphics[width=12pt]{img/Commit16.png}} at (0pt,0pt)
	\end{tikzpicture} & Коммит всех изменений с последнего коммита/отката  & Ctrl+Shift+C\\\hline
	\begin{tikzpicture}
	\pgftext{\includegraphics[width=12pt]{img/Rollback16.png}} at (0pt,0pt)
	\end{tikzpicture} & Отменить все изменения с последнего коммита/отката  & Ctrl+Shift+R\\\hline
	\begin{tikzpicture}
	\pgftext{\includegraphics[width=12pt]{img/AutoCommit16.png}} at (0pt,0pt)
	\end{tikzpicture} & Переключатель режима ручной/авто -- коммит. Если режим автокоммита выключен, то для подтверждения или отката транзакций следует нажимать кнопки \begin{tikzpicture}
	\pgftext{\includegraphics[width=12pt]{img/Commit16.png}} at (0pt,0pt)
	\end{tikzpicture} или  \begin{tikzpicture}
	\pgftext{\includegraphics[width=12pt]{img/Rollback16.png}} at (0pt,0pt)
	\end{tikzpicture} соответственно или прописывать команды \ttt{commit} или \ttt{rollback}. Если же режим автокоммита включен, то каждый оператор будет закоммичен по мере его выполнения. & \\\hline
	\begin{tikzpicture}
	\pgftext{\includegraphics[width=12pt]{img/RecycleConnection16.png}} at (0pt,0pt)
	\end{tikzpicture} & Закрытие соединения, выделенного для редактора запросов, и его повторное открытие & \\\hline
	\begin{tikzpicture}
	\pgftext{\includegraphics[width=12pt]{img/TableRefresh16.png}} at (0pt,0pt)
	\end{tikzpicture} & Обновить список ключевых слов & Alt+Shift+R\\\hline
	\begin{tikzpicture}
	\pgftext{\includegraphics[width=12pt]{img/RSMetaData16.png}} at (0pt,0pt)
	\end{tikzpicture} & Отображает метаданные выбранного множества результатов (Result Set) & \\\hline
	\begin{tikzpicture}
	\pgftext{\includegraphics[width=12pt]{img/TableColumn16.png}} at (0pt,0pt)
	\end{tikzpicture} & Показывает диалоговое окно, в котором можно выбрать какие столбцы выборки скрыть & \\\hline
	\begin{tikzpicture}
	\pgftext{\includegraphics[width=12pt]{img/ExportDelimited16.png}} at (0pt,0pt)
	\end{tikzpicture} & Экспорт выборки в файл & \\\hline
	\begin{tikzpicture}
	\pgftext{\includegraphics[width=12pt]{img/ToggleEditorOutput16.png}} at (0pt,0pt)
	\end{tikzpicture} & Показать/скрыть окно вывода & \\\hline
	\begin{tikzpicture}
	\pgftext{\includegraphics[width=12pt]{img/ShiftTextLeft16.png}} at (0pt,0pt)
	\end{tikzpicture} & Сдвиг текущей строки влево & Ctrl+Shift+D\\\hline
	\begin{tikzpicture}
	\pgftext{\includegraphics[width=12pt]{img/ShiftTextRight16.png}} at (0pt,0pt)
	\end{tikzpicture} & Сдвиг текущей строки вправо & Ctrl+Shift+T\\\hline
	\begin{tikzpicture}
	\pgftext{\includegraphics[width=12pt]{img/AddComment16.png}} at (0pt,0pt)
	\end{tikzpicture} & Закомментировать текущую строку или выделение & Ctrl+/\\\hline
	\begin{tikzpicture}
	\pgftext{\includegraphics[width=12pt]{img/FormatSql16.png}} at (0pt,0pt)
	\end{tikzpicture} & Форматирование SQL кода & Ctrl+Shift+F\\\hline
	\begin{tikzpicture}
	\pgftext{\includegraphics[width=12pt]{img/split.png}} at (0pt,0pt)
	\end{tikzpicture}& Изменить ориентацию разделителя & Ctrl+Alt+Q\\\hline
\end{longtable}	

Горячие клавиши для действий, которых нет в панели инструментов редактора запроса:
\begin{longtable}[r]{|>{\ttfamily}m{3.5cm}|m{11.5cm}|}
	\hline
	\centering\normalfont\bfseries Горячие клавиши &
	\centering\arraybslash\bfseries Описание\\\hline
	\endfirsthead
	\hline
	\centering\normalfont\bfseries Горячие клавиши &
	\centering\arraybslash\bfseries Описание\\\hline
	\endhead
	Ctrl+F & Поиск \\\hline
	Ctrl+R & Найти и заменить\\\hline
	Ctrl+K & Найти слово в позиции курсора, выделить все совпадения и перейти к следующему совпадению\\\hline
	Ctrl + Shift + К & Найти слово в позиции курсора, выделить все совпадения и перейти к предыдущему совпадению\\\hline
	Ctrl + D & Удалить текущую строку\\\hline
	Ctrl + J & Объединить текущую строку со следующей\\\hline
\end{longtable}

\begin{longtable}[c]{|>{\ttfamily}m{2.8cm}|m{12.5cm}|}
	\hline
	Connection & Подключение к базе данных может быть изменено, используя выпадающий список всех доступных соединений \\\hline
	Transaction Isolation Level & Уровень изоляции транзакций можно поменять в текущем окне редактора запросов \\\hline
	Filter & Фильтрация записей в Result Set. Выполните запрос, который вернет множество результатов. Введите в фильтр искомое значение и нажмите \ttt{Enter} --- в новой вкладке откроются совпадающие с фильтром записи. \\\hline
	Max Rows & Урезает выборку до заданного количества строк \\\hline
\end{longtable}	

Для настройки параметров транзакций включите параметр \ttt{Показать параметры транзакции}. В открывшейся панели выберите необходимые вам параметры:

\begin{figure}[H]
	\centering
	\includegraphics[width = 1\linewidth]{img/transaction_param.png}
	\caption{Выбор параметров транзакции}
\end{figure}

Описание параметров:
\begin{itemize}
	\item \ttt{READ ONLY} -- Если включено, то транзакция работает в режиме \ttt{READ ONLY}. Если флаг не установлен, то допустимы как чтение, так и запись (режим \ttt{READ WRITE}).
	\item \ttt{NO AUTO UNDO} -- При использовании опции \ttt{NO AUTO UNDO} оператор \ttt{ROLLBACK} только помечает транзакцию как отменённую без удаления созданных в этой транзакции версий, которые будут удалены позднее в соответствии с выбранной политикой сборки мусора.
	\item \ttt{IGNORE LIMBO} -- При указании опции \ttt{IGNORE LIMBO} игнорируются записи, создаваемые "потерянными" (т.е. незавершёнными) транзакциями (\ttt{limbo transaction}). Транзакции считается "потерянной" , если не завершён второй этап двухфазного подтверждения (\ttt{two-phase commit}).
	\item \ttt{Уровень изоляции} -- Выбор уровня изоляции транзакции.
	\item \ttt{RECORD VERSION} -- При задании \ttt{RECORD\_VERSION} транзакция в процессе работы читает последнюю подтвержденную версию записей в таблицах независимо от того, существуют ли другие неподтвержденные версии этих записей. В этом случае режим разрешения блокировок (\ttt{WAIT} или \ttt{NO WAIT}) никак не влияет на поведение транзакции.
	\item \ttt{WAIT} -- Если включено, то используется режим разрешения блокировок \ttt{WAIT}, при котором в случае возникновения конфликта клиентское приложение будет переведено в режим ожидания до момента разрешения конфликта. Если выключено, то используется режим \ttt{NO WAIT}.
	\item \ttt{Время ожидания} -- Если для режима \ttt{WAIT} задать этот параметр, то ожидание будет продолжаться только указанное количество секунд.
	\item \ttt{RESERVING} -- резервирует указанные таблицы, то есть запрещает другим транзакциям вносить в эти таблицы изменения в то время как выполняется данная транзакция. 
\end{itemize}

\anonsubsection{Панель Вывода}

Панель вывода редактора запросов представляет собой интерфейс с вкладками, отображающими сообщения о выполнении SQL запросов (информация, предупреждения, ошибки и т.д.) и результирующие наборы (Result Set).

При выполнении любого запроса в окне \textit{Output} будет логгироваться информация: о текущем подключении, содержимое SQL запроса, сообщения об успешном выполнении и ошибках, планы запросов и время выполнения.

Панель может быть очищена от ее содержимого в любое время, используя кнопку \begin{tikzpicture}
\pgftext{\includegraphics[width=12pt]{img/ClearOutput16.png}} at (0pt,0pt)
\end{tikzpicture} на панели инструментов редактора или через всплывающее меню по нажатию правой кнопки мыши.

Содержимое окна вывода редактора может сохраняться в файл. Для это откройте настройки редактора (\textit{Инструменты}$\rightarrow$\textit{Настройки}$\rightarrow$\textit{Editor}) и поставьте галочку напротив параметра «\ttt{Log output to file}». А в качестве значения параметра  «\ttt{Output log file path}» укажите путь к лог-файлу, в который будет записываться вывод.

\anonsubsubsection{Result Set}\label{sec:resultset}

Если результатом запроса является Result Set, то он отображается в отдельной вкладке на панели вывода. Если таких запросов несколько, то для каждого создается своя вкладка в порядке выполнения.

При наведении курсором мыши на открытую вкладку \textit{Result Set <N>} вылезает всплывающее окно с запросом, который выдал данный результат. Можно перейти к строке с этим запросом в редакторе, нажав на кнопку 
\begin{tikzpicture}
\pgftext{\includegraphics[width=12pt]{img/GoToResultSetQuery16.png}} at (0pt,0pt)
\end{tikzpicture} или скопировать его в буфер обмена. 

Над самой таблицей с выборкой также можно производить некоторые действия:
\begin{itemize}
	\item сортировать по полю, щелкнув мышью по заголовку нужного столбца
	\item перемещать в разном порядке столбцы таблицы
	\item скопировать выделенные строки или столбцы (см. всплывающее меню)
	\item экспортировать выделенное/таблицу в файл (см. всплывающее меню)
	\item скрыть некоторые поля выборки (см. всплывающее меню)
	\item печать выделенного или всей таблицы (см. всплывающее меню)
	\item посмотреть метаданные результирующей таблицы, нажав на кнопку \begin{tikzpicture}
	\pgftext{\includegraphics[width=12pt]{img/RSMetaData16.png}} at (0pt,0pt)
	\end{tikzpicture} на панели инструментов.
	
	\begin{redremark}
		Эта функция недоступна по умолчанию и должна быть включена в настройках приложения.
	\end{redremark}
\end{itemize}

\anonsubsection{Параметризованные запросы}\label{sec:param_queries}

В некоторых случаях нужно создать запрос, который можно использовать многократно, но каждый раз с разными «входными» значениями. Например, можно написать несколько запросов, чтобы найти данные о сотруднике с определенной фамилией. А можно написать один запрос, меняя только фамилию сотрудника.

Чтобы создать запрос, который в разное время может иметь разные «входные» данные, используются параметры запроса. Параметры могут быть именованными или неименованными. Неименованный параметр — это вопросительный знак («?»), который можно подставлять в любое место запроса, вместо литерального значения. Например:

\begin{redexample}\ttfamily
SELECT * FROM employees WHERE (surname = ?)  
\end{redexample}

Запустив на выполнение такой запрос в редакторе запросов, на экране появится диалоговое окно для ввода значения параметра (фамилии сотрудника):

\begin{figure}[H]
	\centering
	\includegraphics[width = 0.4\linewidth]{img/input_param.png}
\end{figure}

Именованные параметры -- это комбинация двоеточия и имени параметра (\ttt{:<paramname>}), которую также можно подставлять вместо литерального значения. Именованные параметры особенно полезны, если в запросе имеется несколько параметров. Например:

\begin{redexample}\ttfamily
	SELECT * FROM employees WHERE (surname = :surname AND name =:name)  
\end{redexample}

Запустив на выполнение такой запрос в редакторе запросов, на экране появится диалоговое окно для ввода значений параметров (фамилии и имени сотрудника):

\begin{figure}[H]
	\centering
	\includegraphics[width = 0.4\linewidth]{img/input_param2.png}
\end{figure}

\anonsubsection{История запросов}\label{sec:query_history}

После успешного выполнения любого запроса, он сохраняется в кеше журнала редактора для последующего удобного поиска поиска. Количество хранящихся в истории запросов указывается в настройках редактора. Сохраненные запросы не теряются после перезапуска приложения или редактора запросов

Чтобы посмотреть все сохраненные запросы, нажмите на кнопку 
\begin{tikzpicture}
\pgftext{\includegraphics[width=12pt]{img/History16.png}} at (0pt,0pt)
\end{tikzpicture} 
на панели инструментов редактора. Появится диалоговое окно, в котором можно воспользоваться строкой поиска нужного запроса, вставить его в открытое или новое окно редактора (весь текст открытого редактора будет перезаписан выбранным текстом запроса), вставить сохраненный запрос в место позиции курсора, копировать в буфер обмена текст запроса, очистить всю историю. 

Кэш истории можно прокручивать с помощью кнопок 
\begin{tikzpicture}
\pgftext{\includegraphics[width=12pt]{img/Previous16.png}} at (0pt,0pt)
\end{tikzpicture} 
и 
\begin{tikzpicture}
\pgftext{\includegraphics[width=12pt]{img/Forward16.png}} at (0pt,0pt)
\end{tikzpicture} на панели инструментов редактора. 

\newpage


\section{Объекты}\label{sec:objects}

\anonsubsection{ИНДЕКСЫ}
\anonsubsubsection{Создание индексов}\label{sec:crind}

Диалоговое окно \textit{Создать индекс} позволяет создавать новые индексы таблиц базы данных в выбранном соединении с базой данных.

\begin{enumerate}[leftmargin=26pt]
	\item В дереве объектов панели \textit{Подключения} нажмите правой кнопкой мыши по узлу \textit{ИНДЕКСЫ} и выберите пункт \textit{Создать индекс}. Откроется диалоговое окно.
	\begin{figure}[H]
		\centering
		\includegraphics[width = 0.75\linewidth]{img/create_index.png}
	\end{figure}	
	\item Выберите открытое соединение с базой данных.
	\item Введите уникальное имя индекса.
	\item Выберите таблицу базы данных, для которой создается индекс.
	\item Выберите направление индекса (по возрастанию или по убыванию значений ключей индекса).
	\item Если поддерживаются табличные пространства выберите нужное табличное пространство
	\item Для уникальности значений индекса поставьте галочку на параметре \textit{Уникальный}.
	\item Чтобы перевести индекс в активное состояние, поставьте галочку на параметре \textit{Активный}. Иначе индекс будет переведен в неактивное состояние.
	\item Выберите столбцы выбранной таблицы для построения индекса или, если вы хотите построить вычисляемый индекс, поставьте галочку на поле \textit{Вычисляемое как} и впишите выражение.
	\item Напишите комментарий к индексу, если это необходимо.
	\item Нажмите кнопку \textit{OK}. Появится новое диалоговое окно для подтверждения или отката операции создания индекса (см. \hyperref[ris:01]{рис.~\ref{ris:01}}).
	\item Проверьте сгенерированный SQL код.
	\item Сделайте \ttt{commit} или \ttt{rollback} транзакции, нажав на соответствующие кнопки.
\end{enumerate}

\begin{figure}[H]
	\caption{Диалоговое окно подтверждения или отката DDL операций}\label{ris:01}
	\centering
	\includegraphics[width = 0.48\linewidth]{img/create_index2.png}
\end{figure}

\anonsubsubsection{Изменение индексов}
На боковой панели \textit{Подключения} щелкните правой кнопкой мыши по имени индекса и выберите пункт всплывающего меню \textit{Редактировать <имя индекса>} или откройте двойным кликом по нужному индексу. Откроется диалоговое окно \textit{Alter Index}. Алгоритм редактирования индекса полностью аналогичен его \hyperref[sec:crind]{созданию}.

\begin{figure}[H]
	\caption{Диалоговое окно подтверждения или отката DDL операций}\label{ris:01}
	\centering
	\includegraphics[width = 0.48\linewidth]{img/create_index2.png}
\end{figure}

\anonsubsubsection{Удаление индексов}
На боковой панели \textit{Подключения} щелкните правой кнопкой мыши по имени индекса и выберите пункт всплывающего меню \textit{Удалить <имя индекса>}. Откроется  диалоговое окно для подтверждения или отката операции удаления индекса (см. \hyperref[ris:01]{рис.~\ref{ris:01}}). Для подтверждения операции нажмите кнопку «Подтвердить».

\anonsubsubsection{Пересчет селективности индексов}\label{sec:crind}

Диалоговое окно \textit{Пересчет селективности} позволяет обновлять статистику существующих индексов.

\begin{enumerate}[leftmargin=26pt]
	\item В дереве объектов панели \textit{Подключения} нажмите правой кнопкой мыши по узлу \textit{ИНДЕКСЫ} и выберите пункт \textit{Пересчитать статистику всех индексов}. Откроется диалоговое окно.
	\begin{figure}[H]
		\centering
		\includegraphics[width = 0.75\linewidth]{img/selectivity.png}
	\end{figure}	
	\item Нажмите кнопку \textit{Да}. После этого вы увидите список \ttt{SQL-запросов}, выполняющих пересчет статистики у всех индексов.  
	\begin{figure}[H]
		\centering
		\includegraphics[width = 0.75\linewidth]{img/selectivity2.png}
	\end{figure}	
	\item Нажмите кнопку \textit{ОК}.
\end{enumerate}

\anonsubsection{ПРОЦЕДУРЫ}

\anonsubsubsection{Создание процедур}\label{sec:crproc}
Диалоговое окно \textit{Создать процедуру} позволяет создавать новые процедуры в выбранном соединении с базой данных.
\begin{enumerate}[leftmargin=26pt]
	\item На боковой панели \textit{Подключения} щелкните правой кнопкой мыши по объекту \textit{ПРОЦЕДУРЫ} и выберите пункт \textit{Создать процедуру}. Откроется диалоговое окно.
	\begin{figure}[H]
		\centering
		\includegraphics[width = 0.7\linewidth]{img/create_proc.png}
	\end{figure}
	\item Выберите открытое соединение с базой данных.
	\item Введите уникальное имя процедуры.
	\item На вкладке \textit{Edit} можно заполнить таблицы с входными, выходными параметрами и локальными переменными процедуры. \begin{longtable}[r]{|>{\ttfamily}m{3cm}|m{11.5cm}|}
		\hline
		\centering\normalfont\bfseries Имя столбца &
		\centering\arraybslash\bfseries Описание\\\hline
		\endfirsthead
		\hline
		\centering\normalfont\bfseries Имя столбца &
		\centering\arraybslash\bfseries Описание\\\hline
		\endhead
		Name & Имя параметра, может содержать до 31 байта \\\hline
		Datatype & Тип данных SQL параметра\\\hline
		Type of & Используется только с доменом. Если стоит галочка, значит используется только тип данных домена — не проверяется его ограничение (если оно есть в домене) на \ttt{NOT NULL, CHECK} ограничения и/или значения по умолчанию. \\\hline
		Domain & Имя домена в качестве типа параметра\\\hline
		Table & Имя таблицы, если в качестве типа параметра используется тип данных столбцов существующих таблиц.\\\hline
		Column & Имя столбца, если в качестве типа параметра используется тип данных столбцов существующих таблиц.\\\hline
		Size	 & Максимальный размер строки в символах (для строковых типов) Или точность от 1 до 18 для типов данных \ttt{DECIMAL, NUMERIC}. Или размер сегмента (не может превышать 65535) для типа данных \ttt{BLOB}.\\\hline
		Scale	 & Масштаб  для типов данных \ttt{DECIMAL, NUMERIC}. От 0 до 18, должно быть меньше или равно \ttt{Size} \\\hline
		Subtype	 &  Номер подтипа \ttt{BLOB}. \\\hline
		Description	 &  Добавление комментария для параметра процедуры\\\hline
		Default Value	 & Значение по умолчанию для параметра процедуры \\\hline
		Encoding & Кодировка для строковых типов \\\hline	
		Required & Запретить передавать в параметр значение \ttt{NULL}.\\\hline	
	\end{longtable}	
	\item На вкладке \textit{Редактирование} в поле ввода напишите тело процедуры.
	\item На вкладке \textit{Описание} напишите комментарий к процедуре, если это необходимо.
	\item На вкладке \textit{DDL} проверьте сгенерированный SQL код и в случае необходимости исправьте его. 
	
	\item Нажмите кнопку \textit{OK}. Появится новое диалоговое окно для подтверждения или отката операции создания процедуры (см. \hyperref[ris:01]{рис.~\ref{ris:01}}).
	\item Сделайте \ttt{commit} или \ttt{rollback} транзакции, нажав на соответствующие кнопки.
\end{enumerate}

\anonsubsubsection{Изменение процедур}
На боковой панели \textit{Подключения} щелкните правой кнопкой мыши по имени процедуры и выберите пункт всплывающего меню \textit{Редактировать <имя процедуры>}. Откроется диалоговое окно \textit{Редактировать процедуру}. Алгоритм редактирования процедуры полностью аналогичен его \hyperref[sec:crproc]{созданию}.
Также можно открыть вкладку редактирования процедуры двойным кликом по нужной процедуре. В режиме редактирования можно выполнить процедуру, нажав соответствующую кнопку.

\anonsubsubsection{Удаление процедур}
На боковой панели \textit{Подключения} щелкните правой кнопкой мыши по имени процедуры и выберите пункт всплывающего меню \textit{Удалить <имя процедуры>}. Откроется  диалоговое окно для подтверждения или отката операции удаления процедуры (см. \hyperref[ris:01]{рис.~\ref{ris:01}}). Для подтверждения операции нажмите кнопку «Подтвердить».

\anonsubsection{ГЕНЕРАТОРЫ}
\anonsubsubsection{Создание генераторов}\label{sec:crseq}

Диалоговое окно \textit{Создать генератор} позволяет создавать новые генераторы в выбранном соединении с базой данных.

\begin{enumerate}[leftmargin=26pt]
	\item  В дереве объектов панели \textit{Подключения} нажмите правой кнопкой мыши по узлу \textit{ГЕНЕРАТОРЫ} и выберите пункт \textit{Создать генератор}. Откроется диалоговое окно.
	\begin{figure}[H]
		\centering
		\includegraphics[width = 0.65\linewidth]{img/create_seq.png}
	\end{figure}
	\item Выберите открытое соединение с базой данных.
	\item Введите уникальное имя генератора, начальное значение и инкремент.
	\item Напишите комментарий к генератору, если это необходимо.
	\item Нажмите кнопку \textit{OK}. Появится новое диалоговое окно для подтверждения или отката операции создания генератора (см. \hyperref[ris:01]{рис.~\ref{ris:01}}).
	\item Проверьте сгенерированный SQL код.
	\item Сделайте \ttt{commit} или \ttt{rollback} транзакции, нажав на соответствующие кнопки.
\end{enumerate}

\anonsubsubsection{Изменение генераторов}
На боковой панели \textit{Подключения} щелкните правой кнопкой мыши по имени генератора и выберите пункт всплывающего меню \textit{Редактировать <имя генератора>}. Откроется диалоговое окно \textit{Редактировать генератор}. Алгоритм редактирования генератора полностью аналогичен его \hyperref[sec:crseq]{созданию}.
Также можно открыть вкладку редактирования двойным кликом по нужному объекту. 

\anonsubsubsection{Удаление генераторов}
На боковой панели \textit{Подключения} щелкните правой кнопкой мыши по имени генератора и выберите пункт всплывающего меню \textit{Удалить <имя генератора>}. Откроется  диалоговое окно для подтверждения или отката операции удаления генератора (см. \hyperref[ris:01]{рис.~\ref{ris:01}}). Для подтверждения операции нажмите кнопку «Подтвердить».

\anonsubsection{ТАБЛИЦЫ}

\anonsubsubsection{Создание таблиц}\label{sec:111}

Диалоговое окно \textit{Создать таблицу} позволяет создавать новые таблицы базы данных в выбранном соединении с базой данных.

\begin{enumerate}[leftmargin=26pt]
	\item В дереве объектов панели \textit{Подключения} нажмите правой кнопкой мыши по узлу \textit{ТАБЛИЦЫ} и выберите пункт \textit{Создать таблцу}. Откроется диалоговое окно.
	\begin{figure}[H]
		\centering
		\includegraphics[width = 0.65\linewidth]{img/create_table.png}
	\end{figure}
	\item Выберите открытое соединение с базой данных
	\item Введите имя таблицы
	\item Создайте поля таблицы, заполнив таблицу: \label{tab:2}
	\begin{longtable}[r]{|>{\ttfamily}m{3cm}|m{11.5cm}|}
		\hline
		\centering\normalfont\bfseries Имя столбца &
		\centering\arraybslash\bfseries Описание\\\hline
		\endfirsthead
		\hline
		\centering\normalfont\bfseries Имя столбца &
		\centering\arraybslash\bfseries Описание\\\hline
		\endhead
		PK & Двойным кликом по ячейке можно создать ограничение столбца \ttt{PRIMARY KEY}\\\hline
		Name & Имя столбца таблицы, может содержать до 31 байта \\\hline
		Datatype & Тип данных SQL \\\hline
		Domain & Имя домена \\\hline
		Size	 & Максимальный размер строки в символах (для строковых типов) Или точность от 1 до 18 для типов данных \ttt{DECIMAL, NUMERIC}. Или размер сегмента (не может превышать 65535) для типа данных \ttt{BLOB}.\\\hline
		Scale	 & Масштаб  для типов данных \ttt{DECIMAL, NUMERIC}. От 0 до 18, должно быть меньше или равно \ttt{Size} \\\hline
		Subtype	 &  Номер подтипа \ttt{BLOB}. \\\hline
		Required & Указывает, что столбцу не может быть присвоено значение \ttt{NULL}.\\\hline
		Check	 & Ограничение \ttt{CHECK} задаёт условие, которому должны удовлетворять значения, помещаемые 	в данный столбец. Условие — это логическое выражение, называемое также предикат, которое может возвращать значения \ttt{TRUE} (истина), \ttt{FALSE} (ложь) и \ttt{UNKNOWN} (неизвестно) \\\hline
		Description	 &  Добавление комментария для поля таблицы\\\hline
		Computed By  &  Выражение для вычисляемого поля\\\hline
		Autoincrement  & Автоинкремент. Двойной щелчок по ячейке вызовет диалоговое окно, в котором предложат создать или использовать существующий генератор, а также создать триггер, увеличивающий значение генератора и вставляющий в поле его значение \\\hline
		Default Value	 & Значение по умолчанию для столбца таблицы.  \\\hline
		Encoding & Кодировка для строковых типов \\\hline		
	\end{longtable}	
	Новое поле таблицы можно создавать, нажав на одну из кнопок \begin{tikzpicture}
	\pgftext{\includegraphics[width=11pt]{img/ColumnInsertBefore16.png}} at (0pt,0pt)
	\end{tikzpicture} или
	\begin{tikzpicture}
	\pgftext{\includegraphics[width=11pt]{img/ColumnInsertAfter16.png}} at (0pt,0pt)
	\end{tikzpicture}. 
	Чтобы удалить поле, нажмите на кнопку 
	\begin{tikzpicture}
	\pgftext{\includegraphics[width=11pt]{img/ColumnDelete16.png}} at (0pt,0pt)
	\end{tikzpicture}. Также есть возможность изменять последовательность полей с помощью стрелок \begin{tikzpicture}
	\pgftext{\includegraphics[width=11pt]{img/Up16.png}} at (0pt,0pt)
	\end{tikzpicture}, \begin{tikzpicture}
	\pgftext{\includegraphics[width=11pt]{img/Down16.png}} at (0pt,0pt)
	\end{tikzpicture}.
	\item Заполните таблицу с ограничениями: \label{tab:3}
	\begin{longtable}[r]{|>{\ttfamily}m{3cm}|m{11.5cm}|}
		\hline
		\centering\normalfont\bfseries Имя столбца &
		\centering\arraybslash\bfseries Описание\\\hline
		\endfirsthead
		\hline
		\centering\normalfont\bfseries Имя столбца &
		\centering\arraybslash\bfseries Описание\\\hline
		\endhead
		Name &  Имя ограничения, может содержать до 31 байта\\\hline
		Type &  Вид ограничения: \ttt{UNIQUE, PRIMARY KEY, FOREIGN KEY}\\\hline
		Table Column & Имя поля \\\hline
		Reference Schema & Имя схемы, на которую ссылается внешний ключ\\\hline
		Reference Table &  Имя таблицы, на которую ссылается внешний ключ\\\hline
		Reference Column & Столбец таблицы, на которую ссылается внешний ключ\\\hline	
	\end{longtable}	
	\item Во время вставки новых записей в таблицы диалогового окна автоматически генерируется SQL код.
	
	Проверьте сгенерированный SQL код и внесите необходимые изменения для своей базы данных.
	
	Изменения, внесенные в текст SQL, НЕ отражаются в полях таблиц диалогового окна.
	\item Нажмите кнопку «Create».
\end{enumerate}

\anonsubsubsection{Изменение таблиц}\label{sec:alttable}
На боковой панели \textit{Подключения} щелкните правой кнопкой мыши по имени таблицы и выберите пункт всплывающего меню \textit{Редактировать <имя таблицы>}. На центральной панели приложения откроется новая вкладка с подробным описанием данной таблицы. 

\begin{figure}[H]
	\centering
	\includegraphics[width = 1\linewidth]{img/alter_table.png}
\end{figure}

На вкладках \textit{Описание}, \textit{Ограничения}, \textit{Индексы}, \textit{Триггеры} доступны следующие операции изменения таблицы:
\begin{itemize}
	\item добавление нового элемента(столбца, ограничения, индекса или триггера) (нажмите кнопку \begin{tikzpicture}
	\pgftext{\includegraphics[width=11pt]{img/ColumnInsertBefore16.png}} at (0pt,0pt)
	\end{tikzpicture} на панели инструментов)
	\item удаление выбранного элемента (нажмите кнопку \begin{tikzpicture}
	\pgftext{\includegraphics[width=11pt]{img/ColumnDelete16.png}} at (0pt,0pt)
	\end{tikzpicture} на панели инструментов)
	\item изменение характеристик существующих элементов (диалоговое окно редактирования элемента вызовется по двойному клику на ячейку таблицы)
\end{itemize}

\anonsubsubsection{Отображение зависимостей}

С помощью обозревателя зависимостей можно выяснить любые зависимости от объекта и зависимости объекта от других объектов. Этот инструмент имеет большую вспомогательную роль перед изменением типов данных или удалением объектов.

Посмотреть зависимости можно открыв обозреватель конкретного объекта на вкладке \textit{Зависимости}.

\begin{figure}[H]
	\centering
	\includegraphics[width = 0.8\linewidth]{img/dependence.png}
\end{figure}

Колонка \textit{<<Зависимые от объекта>>} отображает какие объекты ссылаются на выбранный объект. Например, из примера видно, что на таблицу \ttt{EMPLOYEE} ссылаются подчиненные объекты: \ttt{EMPLOYEE} (ссылается на себя), \ttt{DEPARTMENT, EMPLOYEE\_PROJECT, PROJECT, SALARY\_HISTORY, SALES}.


Колонка \textit{<<Объекты БД от которых зависит объект>>} показывает, какие объекты используются выбранным объектом. В приведенном выше примере это таблицы \ttt{DEPARTMENT, JOB} и \ttt{EMPLOYEE}.

Двойной щелчок по любому из объектов в средстве просмотра открывает соответствующий редактор объектов.

\anonsubsubsection{Удаление таблиц}
На боковой панели \textit{Подключения} щелкните правой кнопкой мыши по имени таблицы и выберите пункт всплывающего меню \textit{Удалить <имя таблицы>}. Откроется  диалоговое окно для подтверждения или отката операции удаления таблицы (см. \hyperref[ris:01]{рис.~\ref{ris:01}}). Для подтверждения операции нажмите кнопку «Подтвердить».

\anonsubsubsection{Вставка, удаление и обновление данных}
Выполнение запросов на добавление, удаление и обновление записей в таблице с помощью операторов \ttt{INSERT, DELETE} и \ttt{UPDATE} соответственно возможно в Редакторе запросов. Но в Red Expert есть возможность выполнить тоже самое не прибегая к написанию запросов.

Для этого войдите в режим редактирования таблицы, откройте вкладку \textit{Данные}. На центральной панели вы увидите графическое представление таблицы с данными. Нажмите на кнопку \begin{tikzpicture}
\pgftext{\includegraphics[width=11pt]{img/add_16.png}} at (0pt,0pt)
\end{tikzpicture} и добавится новая строка таблицы. Выделите строку таблицы и нажмите кнопку \begin{tikzpicture}
\pgftext{\includegraphics[width=11pt]{img/delete_16.png}} at (0pt,0pt)
\end{tikzpicture}, чтобы удалить эту строку. Двойной клик по ячейке таблицы позволяет редактировать ее содержимое. 

После выполнения всех изменений необходимо подтвердить действия, нажав на кнопку \begin{tikzpicture}
\pgftext{\includegraphics[width=11pt]{img/Commit16.png}} at (0pt,0pt)
\end{tikzpicture} или же откатить их (\begin{tikzpicture}
\pgftext{\includegraphics[width=11pt]{img/Rollback16.png}} at (0pt,0pt)
\end{tikzpicture}).

Ячейки с измененными, но не подтвержденными данными подсвечиваются. Цвет подсветки можно задать в настройках приложения.


\anonsubsection{ТРИГГЕРЫ}

\anonsubsubsection{Создание табличных триггеров}\label{sec:crtrig}
Диалоговое окно \textit{Создать триггер} позволяет создавать новые табличные триггеры в выбранном соединении с базой данных.

\begin{enumerate}[leftmargin=26pt]
	\item  На панели \textit{Подключения} щелкните правой кнопкой мыши по объекту ТРИГГЕРЫ и выберите пункт \textit{Создать TRIGGER} всплывающего меню. Откроется диалоговое окно.
	\begin{figure}[H]
		\centering
		\includegraphics[width = 0.75\linewidth]{img/create_trigger.png}
	\end{figure}
	\item Выберите открытое соединение с базой данных.
	\item Введите уникальное имя триггера.
	\item Задайте порядок, в котором будут выполняться триггеры с одинаковой фазой и событием (или группы событий).
	\item Выберите состояние триггера: активное или неактивное. 
	\item Выберите фазу события срабатывания триггера: \ttt{BEFORE} или \ttt{AFTER}.
	\item Выберите таблицу, на которою создается триггер.
	\item Укажите одно или несколько событий (\ttt{INSERT, UPDATE, DELETE}), при которых вызывается триггер.  
	\item Напишите SQL код тела триггера на вкладке \textit{SQL} .
	\item Напишите комментарий к триггеру, если это необходимо.
	\item Нажмите кнопку \textit{OK}. Появится новое диалоговое окно для подтверждения или отката операции создания триггера (см. \hyperref[ris:01]{рис.~\ref{ris:01}}).
	\item Проверьте сгенерированный SQL код .
	\item Сделайте \ttt{commit} или \ttt{rollback} транзакции, нажав на соответствующие кнопки.
\end{enumerate}	

\anonsubsubsection{Создание триггеров базы данных}
Диалоговое окно \textit{Create Trigger} позволяет создавать новые триггеры на события базы данных в выбранном соединении с базой данных.

\begin{enumerate}[leftmargin=26pt]
	\item  На панели \textit{Подключения} щелкните правой кнопкой мыши по объекту ТРИГГЕРЫ БАЗЫ ДАННЫХ и выберите пункт \textit{Создать DATABASE TRIGGER} всплывающего меню. Откроется диалоговое окно.
	\begin{figure}[H]
		\centering
		\includegraphics[width = 0.75\linewidth]{img/create_dbtrigger.png}
	\end{figure}
	\item Выберите открытое соединение с базой данных
	\item Введите уникальное имя триггера
	\item Выберите тип триггера: \textit{Триггер базы данных}
	\item Задайте порядок, в котором будут выполняться триггеры с одинаковой фазой и событием (или группы событий)
	\item Выберите состояние триггера: активное или неактивное 
	\item Выберите событие, на которое создается триггер 
	\item Напишите SQL код тела триггера на вкладке \textit{SQL} 
	\item Напишите комментарий к триггеру, если это необходимо
	\item Нажмите кнопку \textit{OK}. Появится новое диалоговое окно для подтверждения или отката операции создания триггера (см. \hyperref[ris:01]{рис.~\ref{ris:01}}).
	\item Проверьте сгенерированный SQL код 
	\item Сделайте \ttt{commit} или \ttt{rollback} транзакции, нажав на соответствующие кнопки
\end{enumerate}	

\anonsubsubsection{Создание DDL триггеров}
Диалоговое окно \textit{Create Trigger} позволяет создавать новые триггеры на события изменения метаданных в выбранном соединении с базой данных.

\begin{enumerate}[leftmargin=26pt]
	\item  На панели \textit{Подключения} щелкните правой кнопкой мыши по объекту ТРИГГЕРЫ БАЗЫ ДАННЫХ и выберите пункт \textit{Создать DATABASE TRIGGER} всплывающего меню. Откроется диалоговое окно.
	\begin{figure}[H]
		\centering
		\includegraphics[width = 0.75\linewidth]{img/create_ddltrigger.png}
	\end{figure}
	\item Выберите открытое соединение с базой данных
	\item Введите уникальное имя триггера
	\item Выберите тип триггера: \textit{DDL триггер}
	\item Задайте порядок, в котором будут выполняться триггеры с одинаковой фазой и событием (или группы событий)
	\item Выберите состояние триггера: активное или неактивное 
	\item Выберите фазу события срабатывания триггера: \ttt{BEFORE} или \ttt{AFTER}
	\item Укажите одно или несколько событий, при которых вызывается триггер.  
	\item Напишите SQL код тела триггера на вкладке \textit{SQL} 
	\item Напишите комментарий к триггеру, если это необходимо
	\item Нажмите кнопку \textit{OK}. Появится новое диалоговое окно для подтверждения или отката операции создания триггера (см. \hyperref[ris:01]{рис.~\ref{ris:01}}).
	\item Проверьте сгенерированный SQL код 
	\item Сделайте \ttt{commit} или \ttt{rollback} транзакции, нажав на соответствующие кнопки
\end{enumerate}	

\anonsubsubsection{Изменение триггеров}
На боковой панели \textit{Подключения} щелкните правой кнопкой мыши по имени триггера и выберите пункт всплывающего меню \textit{Редактировать <имя триггера>}. Откроется диалоговое окно \textit{Редактировать триггер}. Алгоритм редактирования триггера полностью аналогичен его \hyperref[sec:crtrig]{созданию}.

\anonsubsubsection{Удаление триггеров}
На боковой панели \textit{Подключения} щелкните правой кнопкой мыши по имени триггера и выберите пункт всплывающего меню \textit{Удалить <имя триггера>}. Откроется  диалоговое окно для подтверждения или отката операции удаления триггера (см. \hyperref[ris:01]{рис.~\ref{ris:01}}). Для подтверждения операции нажмите кнопку «Подтвердить».

\anonsubsection{ПРЕДСТАВЛЕНИЯ}

\anonsubsubsection{Создание представлений}\label{sec:crview}

Диалоговое окно \textit{Создать представление} позволяет создавать новые представления в выбранном соединении с базой данных.

\begin{enumerate}[leftmargin=26pt]
	\item  На панели \textit{Подключения} щелкните правой кнопкой мыши по объекту \textit{ПРЕДСТАВЛЕНИЯ} и выберите пункт \textit{Создать VIEW} всплывающего меню. Откроется диалоговое окно.
	\begin{figure}[H]
		\centering
		\includegraphics[width = 0.7\linewidth]{img/create_view.png}
	\end{figure}
	\item Выберите открытое соединение с базой данных.
	\item Введите уникальное имя представления.
	\item Напишите SQL оператор создания представления, воспользовавшись предоставленным шаблоном.
	\item Напишите комментарий к представлению, если это необходимо.
	\item Нажмите кнопку \textit{OK}. Появится новое диалоговое окно для подтверждения или отката операции создания представления (см. \hyperref[ris:01]{рис.~\ref{ris:01}}).
	\item Проверьте сгенерированный SQL код.
	\item Сделайте \ttt{commit} или \ttt{rollback} транзакции, нажав на соответствующие кнопки.
\end{enumerate}	

\anonsubsubsection{Изменение представлений}
На боковой панели \textit{Подключения} щелкните правой кнопкой мыши по имени представления и выберите пункт всплывающего меню \textit{Редактировать <имя представления>}. Откроется диалоговое окно \textit{Alter View}. Алгоритм редактирования представления полностью аналогичен его \hyperref[sec:crview]{созданию}.

\anonsubsubsection{Удаление представлений}
На боковой панели \textit{Подключения} щелкните правой кнопкой мыши по имени представления и выберите пункт всплывающего меню \textit{Удалить <имя представления>:VIEW}. Откроется  диалоговое окно для подтверждения или отката операции удаления представления (см. \hyperref[ris:01]{рис.~\ref{ris:01}}). Для подтверждения операции нажмите кнопку «Подтвердить».

\anonsubsection{ДОМЕНЫ}

\anonsubsubsection{Создание доменов}\label{sec:crdomain}

Диалоговое окно \textit{Создать домен} позволяет создавать новые представления в выбранном соединении с базой данных.

\begin{enumerate}[leftmargin=26pt]
	\item На панели \textit{Подключения} щелкните правой кнопкой мыши по объекту ДОМЕНЫ и выберите пункт \textit{Создать домен} всплывающего меню. Откроется диалоговое окно.
	\begin{figure}[H]
		\centering
		\includegraphics[width = 0.75\linewidth]{img/create_domain.png}
	\end{figure}
	\item Выберите открытое соединение с базой данных.
	\item Введите уникальное имя домена.
	\item \ttt{NOT NULL} запрещает столбцам и переменным, основанным на домене,
	присваивать значение \ttt{NULL}.
	\item На вкладке \textit{Тип} заполните соответствующие параметры:
	\begin{longtable}[r]{|m{3cm}|m{11cm}|}
		\hline
		\centering\normalfont\bfseries Имя параметра &
		\centering\arraybslash\bfseries Описание\\\hline
		\endfirsthead
		\hline
		\centering\normalfont\bfseries Имя параметра &
		\centering\arraybslash\bfseries Описание\\\hline
		\endhead
		Тип данных & Тип данных SQL \\\hline
		Размер	 & Максимальный размер строки в символах (для строковых типов) Точность от 1 до 18 для типов данных \ttt{DECIMAL, NUMERIC}. Размер сегмента (не может превышать 65535) для типа данных \ttt{BLOB}.\\\hline
		Точность	 & Масштаб  для типов данных \ttt{DECIMAL, NUMERIC}. Должен быть меньше или равно \textit{Размер}\\\hline
		Подтип	 &  Номер подтипа \ttt{BLOB}. \\\hline
		Кодирока & Кодировка для строковых типов \\\hline		
	\end{longtable}	 
	\item Укажите значение по умолчанию для домена (не обязательно).
	\item Задайте ограничение домена (не обязательно).
	\item Напишите комментарий к домену, если это необходимо.
	\item Проверьте сгенерированный SQL код на вкладке \textit{SQL}.
	\item Нажмите кнопку \textit{OK}. Появится новое диалоговое окно для подтверждения или отката операции создания домена (см. \hyperref[ris:01]{рис.~\ref{ris:01}}).
	\item Сделайте \ttt{commit} или \ttt{rollback} транзакции, нажав на соответствующие кнопки.
\end{enumerate}

\anonsubsubsection{Изменение доменов}
На боковой панели \textit{Подключения} щелкните правой кнопкой мыши по имени домена и выберите пункт всплывающего меню \textit{Редактировать <имя домена>}. Откроется диалоговое окно \textit{Редактировать домен}. Алгоритм редактирования домена полностью аналогичен его \hyperref[sec:crdomain]{созданию}.
Также можно открыть вкладку редактирования двойным кликом по нужному объекту. 

\anonsubsubsection{Удаление доменов}
На боковой панели \textit{Подключения} щелкните правой кнопкой мыши по имени домена и выберите пункт всплывающего меню \textit{Удалить <имя домена>}. Откроется  диалоговое окно для подтверждения или отката операции удаления домена (см. \hyperref[ris:01]{рис.~\ref{ris:01}}). Для подтверждения операции нажмите кнопку «Подтвердить».



\anonsubsection{ИСКЛЮЧЕНИЯ}

\anonsubsubsection{Создание исключений}\label{sec:crexc}

Диалоговое окно \textit{Создать исключение} позволяет создавать новые исключения в выбранном соединении с базой данных.

\begin{enumerate}[leftmargin=26pt]
	\item На панели \textit{Подключения} щелкните правой кнопкой мыши по объекту ИСКЛЮЧЕНИЯ и выберите пункт \textit{Создать исключение} всплывающего меню. Откроется диалоговое окно.
	\begin{figure}[H]
		\centering
		\includegraphics[width = 0.75\linewidth]{img/create_exception.png}
	\end{figure}
	\item Выберите открытое соединение с базой данных.
	\item Введите уникальное имя исключения.
	\item Введите сообщение об ошибке.
	\item Напишите комментарий к исключению, если это необходимо.
	\item Нажмите кнопку \textit{OK}. Появится новое диалоговое окно для подтверждения или отката операции создания исключения (см. \hyperref[ris:01]{рис.~\ref{ris:01}}).
	\item Проверьте сгенерированный SQL код.
	\item Сделайте  \ttt{commit} или \ttt{rollback} транзакции, нажав на соответствующие кнопки.
\end{enumerate}

\anonsubsubsection{Изменение исключений}
На боковой панели \textit{Подключения} щелкните правой кнопкой мыши по имени исключения и выберите пункт всплывающего меню \textit{Редактировать <имя исключения>}. Откроется диалоговое окно \textit{Редактировать исключение}. Алгоритм редактирования исключения полностью аналогичен его \hyperref[sec:crexc]{созданию}.
Также можно открыть вкладку редактирования двойным кликом по нужному объекту. 

\anonsubsubsection{Удаление исключений}
На боковой панели \textit{Подключения} щелкните правой кнопкой мыши по имени исключения и выберите пункт всплывающего меню \textit{Удалить <имя исключения>}. Откроется  диалоговое окно для подтверждения или отката операции удаления исключения (см. \hyperref[ris:01]{рис.~\ref{ris:01}}). Для подтверждения операции нажмите кнопку «Подтвердить».

\anonsubsection{ВНЕШНИЕ ФУНКЦИИ}

\anonsubsubsection{Создание внешних функций}\label{sec:crexfunc}

Диалоговое окно \textit{Создать внешнюю функцию} позволяет создавать новые функции, определенные пользователем, в выбранном соединении с базой данных.

\begin{enumerate}[leftmargin=26pt]
	\item На панели \textit{Подключения} щелкните правой кнопкой мыши по объекту ВНЕШНИЕ ФУНКЦИИ и выберите пункт \textit{Создать внешнюю функцию} всплывающего меню. Откроется диалоговое окно.
	\begin{figure}[H]
		\centering
		\includegraphics[width = 0.75\linewidth]{img/create_exfunc.png}
	\end{figure}
	\item Выберите открытое соединение с базой данных.
	\item Введите уникальное имя внешней функции.
	\item Введите имя модуля, в котором находится экспортируемая функция.
	\item Введите имя точки входа (имя экспортируемой функции) в модуле.
	\item Выберите будет ли выходной параметр SQL типом, нуль терминальной строкой или одним из входных параметров.
	\item Если выходной параметр будет SQL типом, выберите как он будет передан: по ссылке, по дескриптору или по значению. На вкладке Returns Type заполните значения для выходного параметра:
	\begin{longtable}[r]{|m{3cm}|m{11cm}|}
		\hline
		\centering\normalfont\bfseries Имя параметра &
		\centering\arraybslash\bfseries Описание\\\hline
		\endfirsthead
		\hline
		\centering\normalfont\bfseries Имя параметра &
		\centering\arraybslash\bfseries Описание\\\hline
		\endhead
		Тип данных & Тип данных SQL \\\hline
		Размер	 & Максимальный размер строки в символах (для строковых типов) Точность от 1 до 18 для типов данных \ttt{DECIMAL, NUMERIC}. Размер сегмента (не может превышать 65535) для типа данных \ttt{BLOB}.\\\hline
		Точность	 & Масштаб  для типов данных \ttt{DECIMAL, NUMERIC}. Должен быть меньше или равно \textit{Размер}\\\hline
		Подтип	 &  Номер подтипа \ttt{BLOB}. \\\hline
		Кодирока & Кодировка для строковых типов \\\hline		
	\end{longtable}
	
	\item Если выходной параметр будет нуль терминальной строкой, то задайте ее максимальную длину в байтах.
	\item Если выходной параметр будет одним из входных параметров, то укажите номер входного параметра, который будет возвращён функцией.
	\item Параметр \ttt{FREE\_IT} означает, что память, выделенная для хранения возвращаемого значения, будет освобождена после завершения выполнения функции. Применяется только 	в том случае, если эта память в UDF выделялась динамически. 
	\item Введите данные входных параметров функции в таблицу вкладки Input Parameters:
	\begin{longtable}[r]{|m{3cm}|m{11cm}|}
		\hline
		\centering\normalfont\bfseries Имя параметра &
		\centering\arraybslash\bfseries Описание\\\hline
		\endfirsthead
		\hline
		\centering\normalfont\bfseries Имя параметра &
		\centering\arraybslash\bfseries Описание\\\hline
		\endhead
		Datatype & Тип данных SQL \\\hline
		Size	 & Максимальный размер строки в символах (для строковых типов) Точность от 1 до 18 для типов данных \ttt{DECIMAL, NUMERIC}. Размер сегмента (не может превышать 65535) для типа данных \ttt{BLOB}.\\\hline
		Scale	 & Масштаб  для типов данных \ttt{DECIMAL, NUMERIC}. Должен быть меньше или равно \textit{Размер}\\\hline
		Subtype	 &  Номер подтипа \ttt{BLOB}. \\\hline
		Encoding & Кодировка для строковых типов \\\hline	
		By Descriptor &  Если указано, то входной параметр передаётся по дескриптору. Иначе по ссылке\\\hline
		NULL &  Если указано, то при передаче значение \ttt{NULL} оно попадёт в функцию в виде нулевого указателя (\ttt{null}).\\\hline
		CSTRING &  Нуль терминальная строка вместо SQL типа \\\hline	
	\end{longtable}
	\item Напишите комментарий к функции, если это необходимо.
	\item Нажмите кнопку \textit{OK}. Появится новое диалоговое окно для подтверждения или отката операции создания внешней функции (см. \hyperref[ris:01]{рис.~\ref{ris:01}}).
	\item Проверьте сгенерированный SQL код.
	\item Сделайте \ttt{commit} или \ttt{rollback} транзакции, нажав на соответствующие кнопки.
\end{enumerate}

\anonsubsubsection{Изменение внешних функций}
На боковой панели \textit{Подключения} щелкните правой кнопкой мыши по имени внешней функции и выберите пункт всплывающего меню \textit{Редактировать <имя функции>}. Откроется диалоговое окно \textit{Редактировать внешнюю функцию}. Алгоритм редактирования внешних функций полностью аналогичен его \hyperref[sec:crexfunc]{созданию}.
Также можно открыть вкладку редактирования двойным кликом по нужному объекту. 

\anonsubsubsection{Удаление внешних функций}
На боковой панели \textit{Подключения} щелкните правой кнопкой мыши по имени внешней функции и выберите пункт всплывающего меню \textit{Удалить <имя функции>}. Откроется  диалоговое окно для подтверждения или отката операции удаления UDF (см. \hyperref[ris:01]{рис.~\ref{ris:01}}). Для подтверждения операции нажмите кнопку «Подтвердить».

\anonsubsection{ФУНКЦИИ}

\anonsubsubsection{Создание функций}\label{sec:crfunc}

Диалоговое окно \textit{Создать функцию} позволяет создавать новые хранимые функции в выбранном соединении с базой данных.

\begin{enumerate}[leftmargin=26pt]
	\item На панели \textit{Подключения} щелкните правой кнопкой мыши по объекту ФУНКЦИИ и выберите пункт \textit{Создать FUNCTION} всплывающего меню. Откроется диалоговое окно.
	\begin{figure}[H]
		\centering
		\includegraphics[width = 0.75\linewidth]{img/create_func.png}
	\end{figure}
	\item Выберите открытое соединение с базой данных.
	\item Введите уникальное имя функции.
	\item На вкладке Arguments заполните таблицу, описывающую входные параметры функции.
	\item На вкладке Variables заполните таблицу, описывающую локальные переменные.
	\item На вкладке Returns type задается тип возвращаемого значения. В качестве типа выходного значения можно указать: тип SQL, имя домена, ссылку на тип столбца таблицы (с помощью параметра Type Of).
	\item На вкладке Edit напишите тело функции.
	\item Напишите комментарий к функции, если это необходимо.
	\item Проверьте сгенерированный SQL код на вкладке DDL.
	\item Нажмите кнопку \textit{OK}. Появится новое диалоговое окно для подтверждения или отката операции создания функции (см. \hyperref[ris:01]{рис.~\ref{ris:01}}).
	\item Проверьте сгенерированный SQL код.
	\item Сделайте \ttt{commit} или \ttt{rollback} транзакции, нажав на соответствующие кнопки.
\end{enumerate}

\anonsubsubsection{Изменение функций}
На боковой панели \textit{Подключения} щелкните правой кнопкой мыши по имени хранимой функции и выберите пункт всплывающего меню \textit{Редактировать <имя функции>}. Откроется диалоговое окно \textit{Редактировать функцию}. Алгоритм редактирования функций полностью аналогичен его \hyperref[sec:crfunc]{созданию}.
В режиме редактирования можно выполнить функцию, нажав соответствующую кнопку.

\anonsubsubsection{Удаление функций}
На боковой панели \textit{Подключения} щелкните правой кнопкой мыши по имени функции и выберите пункт всплывающего меню \textit{Удалить <имя функции>}. Откроется  диалоговое окно для подтверждения или отката операции удаления функции (см. \hyperref[ris:01]{рис.~\ref{ris:01}}). Для подтверждения операции нажмите кнопку «Подтвердить».

\anonsubsection{ГЛОБАЛЬНЫЕ ВРЕМЕННЫЕ ТАБЛИЦЫ}

\anonsubsubsection{Создание глобальных временных таблиц}

Диалоговое окно \textit{Создать таблицу} позволяет создавать новые глобальные временные таблицы в выбранном соединении с базой данных.

\begin{enumerate}[leftmargin=26pt]
	\item На панели \textit{Подключения} щелкните правой кнопкой мыши по объекту ГЛОБАЛЬНЫЕ ВРЕМЕННЫЕ ТАБЛИЦЫ и выберите пункт \textit{Создать глобальную временную таблицу} всплывающего меню. Откроется диалоговое окно.
	\begin{figure}[H]
		\centering
		\includegraphics[width = 0.75\linewidth]{img/create_gtt.png}
	\end{figure}
	\item Выберите открытое соединение с базой данных.
	\item Введите уникальное имя таблицы.
	\item Выберите будет ли создана GTT транзакционного уровня (DELETE ROWS). или уровня соединения с базой данных (PRESERVE ROWS ).
	\item Создайте поля таблицы (см. \hyperref[sec:111]{Создание таблиц}).
	\item Заполните таблицу с ограничениями (см. \hyperref[sec:111]{Создание таблиц}).
	\item Нажмите кнопку \textit{Create}. Появится новое диалоговое окно для подтверждения или отката операции создания таблицы (см. \hyperref[ris:01]{рис.~\ref{ris:01}}).
	\item Проверьте сгенерированный SQL код.
	\item Сделайте \ttt{commit} или \ttt{rollback} транзакции, нажав на соответствующие кнопки.
\end{enumerate}

\anonsubsubsection{Изменение GTT}
На боковой панели \textit{Подключения} щелкните правой кнопкой мыши по имени глобальной временной таблицы и выберите пункт всплывающего меню \textit{Редактировать <имя таблицы>}. На центральной панели откроется новая вкладка с описанием данной таблицы. О выполнении операций изменения таблиц уже рассказано в разделе \hyperref[sec:alttable]{Изменение таблиц}.

\anonsubsubsection{Удаление GTT}
На боковой панели \textit{Подключения} щелкните правой кнопкой мыши по имени глобальной временной таблицы и выберите пункт всплывающего меню \textit{Удалить <имя таблицы>}. Откроется  диалоговое окно для подтверждения или отката операции удаления GTT (см. \hyperref[ris:01]{рис.~\ref{ris:01}}). Для подтверждения операции нажмите кнопку «Подтвердить».

\anonsubsection{ПАКЕТЫ}

\anonsubsubsection{Создание пакетов}\label{sec:crpack}

Диалоговое окно \textit{Создать пакет} позволяет создавать новые пакеты в выбранном соединении с базой данных.

\begin{enumerate}[leftmargin=26pt]
	\item На панели \textit{Подключения} щелкните правой кнопкой мыши по объекту ПАКЕТЫ и выберите пункт \textit{Создать пакет} всплывающего меню. Откроется диалоговое окно.
	\begin{figure}[H]
		\centering
		\includegraphics[width = 0.75\linewidth]{img/create_package.png}
	\end{figure}
	\item Выберите открытое соединение с базой данных.
	\item Введите уникальное имя пакета.
	\item Создайте заголовок пакета.
	\item Создайте тело пакета.
	\item Напишите комментарий к пакету, если это необходимо.
	\item Нажмите кнопку \textit{OK}. Появится новое диалоговое окно для подтверждения или отката операции создания пакета (см. \hyperref[ris:01]{рис.~\ref{ris:01}}).
	\item Проверьте сгенерированный SQL код.
	\item Сделайте \ttt{commit} или \ttt{rollback} транзакции, нажав на соответствующие кнопки.
\end{enumerate}

\anonsubsubsection{Изменение пакетов}
На боковой панели \textit{Подключения} щелкните правой кнопкой мыши по имени пакета и выберите пункт всплывающего меню \textit{Редактировать <имя пакета>}. Откроется диалоговое окно \textit{Редактировать пакет}. Алгоритм редактирования пакета полностью аналогичен его \hyperref[sec:crpack]{созданию}.

\anonsubsubsection{Удаление пакетов}
На боковой панели \textit{Подключения} щелкните правой кнопкой мыши по имени пакета и выберите пункт всплывающего меню \textit{Удалить <имя пакета>}. Откроется  диалоговое окно для подтверждения или отката операции удаления пакета (см. \hyperref[ris:01]{рис.~\ref{ris:01}}). Для подтверждения операции нажмите кнопку «Подтвердить».

\anonsubsection{РОЛИ}

\anonsubsubsection{Создание ролей}

Диалоговое окно \textit{Create role} позволяет создавать новые роли в выбранном соединении с базой данных.

\begin{enumerate}[leftmargin=26pt]
	\item На панели \textit{Подключения} щелкните правой кнопкой мыши по объекту РОЛИ и выберите пункт \textit{Создать роль} всплывающего меню. Откроется диалоговое окно.
	\begin{figure}[H]
		\centering
		\includegraphics[width = 0.35\linewidth]{img/create_role.png}
	\end{figure}
	\item Введите уникальное имя роли.
	\item Нажмите кнопку \textit{OK}. Появится новое диалоговое окно для подтверждения или отката операции создания роли (см. \hyperref[ris:01]{рис.~\ref{ris:01}}).
	\item Проверьте сгенерированный SQL код.
	\item Сделайте \ttt{commit} или \ttt{rollback} транзакции, нажав на соответствующие кнопки.
\end{enumerate}

\anonsubsubsection{Изменение ролей}
На боковой панели \textit{Подключения} щелкните правой кнопкой мыши по имени роли и выберите пункт всплывающего меню \textit{Редактировать <имя роли>}. На центральной панели откроется новая вкладка для просмотра и редактирования привилегий роли. 	

\begin{figure}[H]
	\centering
	\includegraphics[width = 1\linewidth]{img/alter_role.png}
\end{figure}

На ней представлена таблица, отражающая привилегии данной роли со всеми объектами базы данных. Так точка \begin{tikzpicture}
\pgftext{\includegraphics[width=12pt]{img/grant.png}} at (0pt,0pt)
\end{tikzpicture} говорит о том, что роли предоставлена привилегия (\ttt{SELECT, UPDATE, DELETE, INSERT, EXECUTE} или \ttt{REFERENCES}) для объектов базы
данных. Точка  \begin{tikzpicture}
\pgftext{\includegraphics[width=12pt]{img/no_grant.png}} at (0pt,0pt)
\end{tikzpicture} означает, что у роли нет таких прав на объект. Точка \begin{tikzpicture}
\pgftext{\includegraphics[width=12pt]{img/admin_option.png}} at (0pt,0pt)
\end{tikzpicture} говорит о возможности роли передавать привилегию другому пользователю или роли. 

Изменить привилегии роли можно с помощью оператора \ttt{GRANT} или двойным кликом по нужной ячейке таблицы. 

\anonsubsubsection{Удаление ролей}
На боковой панели \textit{Подключения} щелкните правой кнопкой мыши по имени роли и выберите пункт всплывающего меню \textit{Удалить <имя роли>}. Откроется  диалоговое окно для подтверждения или отката операции удаления роли (см. \hyperref[ris:01]{рис.~\ref{ris:01}}). Для подтверждения операции нажмите кнопку «Подтвердить».

\anonsubsection{ПОЛЬЗОВАТЕЛИ}

\anonsubsubsection{Создание пользователей}

Диалоговое окно \textit{Создать пользователя} позволяет создавать новых пользователей в выбранном соединении с базой данных.

\begin{enumerate}[leftmargin=26pt]
	\item На панели \textit{Подключения} щелкните правой кнопкой мыши по объекту ПОЛЬЗОВАТЕЛИ и выберите пункт \textit{Создать пользователя} всплывающего меню. Откроется диалоговое окно.
	\begin{figure}[H]
		\centering
		\includegraphics[width = 0.6\linewidth]{img/create_user.png}
	\end{figure}
	\item Введите запрашиваемые данные в данном окне.
	\item Нажмите кнопку \textit{OK}. Появится новое диалоговое окно для подтверждения или отката операции создания пользователя (см. \hyperref[ris:01]{рис.~\ref{ris:01}}).
	\item Проверьте сгенерированный SQL код.
	\item Сделайте \ttt{commit} или \ttt{rollback} транзакции, нажав на соответствующие кнопки.
\end{enumerate}

\anonsubsubsection{Изменение пользователей}
На боковой панели \textit{Подключения} щелкните правой кнопкой мыши по имени пользователя и выберите пункт всплывающего меню \textit{Редактировать <имя пользователя>}. откроется окно редактирования пользователя

\anonsubsubsection{Удаление пользователей}
На боковой панели \textit{Подключения} щелкните правой кнопкой мыши по имени пользователя и выберите пункт всплывающего меню \textit{Удалить <имя пользователя>}. Откроется  диалоговое окно для подтверждения или отката операции удаления пользователя (см. \hyperref[ris:01]{рис.~\ref{ris:01}}). Для подтверждения операции нажмите кнопку «Подтвердить».

\anonsubsection{ТАБЛИЧНЫЕ ПРОСТРАНСТВА}

\anonsubsubsection{Создание табличного пространства}

Диалоговое окно \textit{Создать табличное пространство} позволяет создавать новые табличные пространства в выбранном соединении с базой данных.

\begin{enumerate}[leftmargin=26pt]
	\item На панели \textit{Подключения} щелкните правой кнопкой мыши по объекту ТАБЛИЧНЫЕ ПРОСТРАНСТВА и выберите пункт \textit{Создать табличное пространство} всплывающего меню. Откроется диалоговое окно.
	\begin{figure}[H]
		\centering
		\includegraphics[width = 0.6\linewidth]{img/create_tablespace.png}
	\end{figure}
	\item Введите запрашиваемые данные в данном окне.
	\item Нажмите кнопку \textit{OK}. Появится новое диалоговое окно для подтверждения или отката операции создания роли (см. \hyperref[ris:01]{рис.~\ref{ris:01}}).
	\item Проверьте сгенерированный SQL код.
	\item Сделайте \ttt{commit} или \ttt{rollback} транзакции, нажав на соответствующие кнопки.
\end{enumerate}

\anonsubsubsection{Изменение табличных пространств}
На боковой панели \textit{Подключения} кликните 2 раза мышкой по нужному табличному пространству. откроется вкладка редактирования табличного пространства. Для добавления таблиц и индексов в табличное пространство перетащите их из дерева объектов. 

\anonsubsubsection{Удаление табличных пространств}
На боковой панели \textit{Подключения} щелкните правой кнопкой мыши по имени табличного пространства и выберите пункт всплывающего меню \textit{Удалить <имя табличного пространства>}. Откроется  диалоговое окно для подтверждения или отката операции удаления пространства (см. \hyperref[ris:01]{рис.~\ref{ris:01}}). Для подтверждения операции нажмите кнопку «Подтвердить».

\anonsubsection{ЗАДАЧИ}

\anonsubsubsection{Создание задачи}

Диалоговое окно \textit{Создать задачу} позволяет создавать новые зададания для планировщика в выбранном соединении с базой данных.

\begin{enumerate}[leftmargin=26pt]
	\item На панели \textit{Подключения} щелкните правой кнопкой мыши по объекту ЗАДАЧИ и выберите пункт \textit{Создать задачу} всплывающего меню. Откроется диалоговое окно.
	\begin{figure}[H]
		\centering
		\includegraphics[width = 0.6\linewidth]{img/create_job.png}
	\end{figure}
	\item Выберите открытое соединение с базой данных.
	\item Введите уникальное имя задания.
	\item Необязательные параметры \textit{Время запуска} и \textit{Время завершения} задают промежуток времени, в котором будет запускаться задание. Значение \ttt{NULL} используется, чтобы убрать интервал. 
	\item Выберите тип задачи. Тип \ttt{BASH} дает возможность запланировать выполнение команды ОС. Такой тип заданий может создавать только \ttt{SYSDBA}. Поле \textit{Активный} указывает, что задание будет запускаться по расписанию. 
	\item Введите содержание задания. 
	\item Заполните расписание, в соответствии с которым будет выполняться задание. Если поставить галочку \textit{Расширенные свойства}, то откроются расширенные параметры расписания.
	\begin{figure}[H]
		\centering
		\includegraphics[width = 0.6\linewidth]{img/create_job_cron.png}
	\end{figure}
\end{enumerate}

\anonsubsubsection{Изменение задачи}

На боковой панели \textit{Подключения} щелкните правой кнопкой мыши по задаче и выберите пункт всплывающего меню \textit{Редактировать <задача>}. Также открыть окно редактирования можно двойным кликом по задаче.

Для просмотра событий выполнения задачи откройте вкладку \textit{Журнал}. Параметры \textit{Начало периода} и \textit{Конец периода} позволяют указать период, за который нужно отобразить лог. Чтобы посмотреть события только определённого типа, поставьте нужные галочки для полей \textit{RUN\_START}, \textit{RUN\_FINISH} и \textit{RUN\_ERROR}.
\begin{figure}[H]
	\centering
	\includegraphics[width = 0.6\linewidth]{img/create_job_log.png}
\end{figure}

\anonsubsubsection{Удаление задачи}

На боковой панели \textit{Подключения} щелкните правой кнопкой мыши по задаче и выберите пункт всплывающего меню \textit{Удалить <задача>}. Откроется  диалоговое окно для подтверждения или отката операции удаления задачи. Для подтверждения нажмите кнопку «Подтвердить».

\newpage
\section{Перекомпиляция объектов}\label{sec:recompile_objects}
\begin{enumerate}[leftmargin=20pt]
	\item Перекомпиляция доступна для процедур, пакетов, функций и триггеров. Здесь показан пример с процедурами. Для остальных объектов алгоритм действий аналогичный. Есть возможность перекомпилировать только невалидные объекты. При полной перекомпиляции невалидные объекты будут пересозданы в первую очередь.
	\item Нажмите правой кнопкой мыши на узел с процедурами (триггерами, функциями и пакетами). В контекстном меню выберите пункт \textit{Перекомпилировать все Процедуры} или \textit{Перекомпилировать все невалидные Процедуры}.
	\item Подтвердите действие в открывшемся диалоговом окне. После этого вы увидите список \ttt{SQL-запросов} пересоздания каждой процедуры. 
	\begin{figure}[H]
		\centering
		\includegraphics[width = 0.7\linewidth]{img/recompile2.png}
	\end{figure}
	\item Нажмите кнопку \textit{ОК}. 
\end{enumerate}
\newpage
\section{Выполнение хранимых объектов}\label{sec:execute_proc}

\begin{enumerate}[leftmargin=20pt]
	\item Нажмите на кнопку \begin{tikzpicture}
	\pgftext{\includegraphics[width=10pt]{img/Procedure16.png}} at (0pt,0pt)
	\end{tikzpicture} на главной панели инструментов. Откроется вкладка на центральной панели.
	\begin{figure}[H]
		\centering
		\includegraphics[width = 0.9\linewidth]{img/execute_store_proc.png}
	\end{figure}
	
	\item Выберите соединение, которое будет использоваться для выполнения объектов.
	\item Выберите тип объекта: функция или процедура
	\item Выберите доступные объекты для отображения соответствующих входных и выходных параметров
	\item Нажмите кнопку \textit{Execute}, чтобы выполнить процедуру или функцию
	\item Введите значения входных параметров в новом диалоговом окне
	\item Результаты или ошибки будут отображаться на панели вывода в нижней части вкладки.
\end{enumerate}

\newpage

\section{Сравнение типов данных}\label{sec:compare_datatype}

В Red Expert есть возможность сравнить типы данных разных СУБД. Задача состоит в том, чтобы найти аналогичные типы данных, которые наиболее точно соответствуют друг другу. Это часто требуется при переносе базы данных из одной СУБД в другую.

Требуется два открытых подключения к базам данных от разных СУБД.

\begin{enumerate}[leftmargin=39pt]
	\item Выберите пункт меню  \textit{База данных}$\rightarrow$\textit{Сравнить типы данных} или кнопку \begin{tikzpicture}
	\pgftext{\includegraphics[width=10pt]{img/CompareDataTypes16.png}} at (0pt,0pt)
	\end{tikzpicture} на главной панели инструментов. Откроется вкладка на центральной панели.
	\begin{figure}[H]
		\flushright
		\includegraphics[width = 0.9\linewidth]{img/compare_types.png}
	\end{figure}
	\item Выберите два соединения для сравнения. Первое соединение используется как основное, против которого выполняется сравнение.
	\item Выберите какой-либо тип данных из первого подключения для отображения эквивалентных типов данных из второго подключения
	\item Изучите различия типов данных в таблице метаданных в нижней части вкладки. Первая строка относится к основному подключению и она всегда выделена зеленым цветом. Любые различия между типами данных будут выделены желтым цветом.
\end{enumerate}

\newpage

\section{Экспорт результатов}\label{sec:export_results}
Данная опция предназначена для экспорта данных таблиц, представлений, а также результатов выборки \ttt{SELECT} в различные форматы. 

\anonsubsection{Экспорт данных таблиц и представлений}

Чтобы экспортировать данные таблицы или представления выполните следующее:
\begin{enumerate}
	\item Откройте интересующую таблицу/представление и перейдите во вкладку \ttt{Данные};
	\item Нажмите на таблицу правой кнопкой мыши и выберите в открывшемся списке пункт \ttt{Экспортировать}:
	\begin{figure}[H]
		\centering
		\includegraphics[width = 0.7\linewidth]{img/export_set.png}
	\end{figure}
	\item В открывшемся окне выберите нужные параметры экспорта и нажмите кнопку \ttt{Экспортировать}:
	\begin{figure}[H]
		\centering
		\includegraphics[width = 0.6\linewidth]{img/export_table_sql.png}
	\end{figure}
	\item Данные будут экспортированы в указанный файл.
\end{enumerate}

\anonsubsection{Экспорт результатов запроса}

Экспортировать результаты запроса можно двумя способами:
\begin{itemize}
	\item Нажав правой кнопкой мыши на полученный результирующий набор и выбрав в открывшемся списке пункт \ttt{Экспортировать}:
	\begin{figure}[H]
		\centering
		\includegraphics[width = 1\linewidth]{img/export_result_set1.png}
	\end{figure}
	\item Выбрав пункт \ttt{Выполнить в файл} перед выполнением запроса:
	\begin{figure}[H]
		\centering
		\includegraphics[width = 1\linewidth]{img/export_result_set2.png}
	\end{figure}
\end{itemize}
 В открывшемся окне выберите нужные параметры экспорта и нажмите кнопку \ttt{Экспортировать}. Данные будут экспортированы в указанный файл.
 
\anonsubsection{Форматы экспорта}

Поддерживается экспорт в \ttt{CSV}, \ttt{XLSX}, \ttt{XML} и \ttt{SQL} файлы.

Общие параметры экспорта:
\begin{itemize}
	\item Формат файла -- Фомат файла, в который нужно экспортировать данные;
	\item Путь к файлу -- Путь к файлу, в который будут экспортированы данные;
	\item Путь к BLOB файлам -- Определяет путь к папке, в которую будут экспортированы блоб файлы, если указана галочка \ttt{Сохранять блобы по-отдельности}. Если не указана, то параметр определяет файл, куда будут экспортированы все блобы;
	\item Вкладка \ttt{Поля} -- Показывает поля, которые будут экспортированы. Экспорт блобов и вычисляемых полей по умолчанию отключен. 
\end{itemize}

\anonsubsubsection{Экспорт в CSV файл}
Панель экспорта в \ttt{CSV-файл}:

\begin{figure}[H]
	\begin{center}
		\begin{minipage}[h]{0.49\linewidth}
			\includegraphics[width=1\linewidth]{img/export_to_csv.png}
			\caption{Экспорт в CSV файл} 
		\end{minipage}
		\hfill 
		\begin{minipage}[h]{0.49\linewidth}
			\includegraphics[width=1\linewidth]{img/export_to_csv_params.png}
			\caption{Параметры экспорта}
		\end{minipage}
	\end{center}
\end{figure}

Описание параметров экспорта:
\begin{itemize}
	\item Включить имена столбцов в качестве первой строки -- Если включено, то первой строке будут добавлены заголовки столбцов;
	\item Используйте двойные кавычки для столбцов char/varchar/longvarchar -- Если включео, то строковые значения будут обёрнуты в двойные кавычки;
	\item Сохранять BLOB файлы по отдельности -- Если включено, то блобы будут сохраняться в указанную папку отдельными файлами;
	\item Заменить \rr{NULL} на --- Если включено, то определяет, на что заменить значения NULL при экспорте;
	\item Заменить \rr{\textbackslash n} -- Если включено, то определяет на что заменить перенос строки при экспорте;
	\item Разделитель -- Опредеяет, что будет использоваться в качестве разделителя.
\end{itemize}

\anonsubsubsection{Экспорт в XLSX файл}
Панель экспорта в \ttt{XLSX-файл}:

\begin{figure}[H]
	\centering
	\includegraphics[width = 0.5\linewidth]{img/export_to_xlsx.png}
	\caption{Экспорт в XLSX файл} 
\end{figure}

Описание параметров экспорта:
\begin{itemize}
	\item Включить имена столбцов в качестве первой строки -- Если включено, то первой строке будут добавлены заголовки столбцов;
	\item Сохранять BLOB файлы по отдельности -- Если включено, то блобы будут сохраняться в указанную папку отдельными файлами;
	\item Заменить \rr{NULL} на --- Если включено, то определяет, на что заменить значения NULL при экспорте;
\end{itemize}

\anonsubsubsection{Экспорт в XML файл}
Панель экспорта в \ttt{XML-файл}:

\begin{figure}[H]
	\centering
	\includegraphics[width = 0.5\linewidth]{img/export_to_xml.png}
	\caption{Экспорт в XML файл} 
\end{figure}

Описание параметров экспорта:
\begin{itemize}
	\item Сохранять BLOB файлы по отдельности -- Если включено, то блобы будут сохраняться в указанную папку отдельными файлами;
	\item Заменить \rr{NULL} на --- Если включено, то определяет, на что заменить значения NULL при экспорте;
\end{itemize}

\anonsubsubsection{Экспорт в SQL файл}
Панель экспорта в \ttt{SQL-файл}:

\begin{figure}[H]
	\centering
	\includegraphics[width = 0.5\linewidth]{img/export_to_sql.png}
	\caption{Экспорт в SQL файл} 
\end{figure}

Описание параметров экспорта:
\begin{itemize}
	\item Открыть сгенерированный SQL-скрипт в Редакторе Запросов -- Если включено, то полученный после экспорта sql-скрипт будет открыт в редакторе запросов;
	\item Добавить скрипт создания таблицы -- Если включено, то в указанный файл будет добавлен скрипт создания таблицы, в которую будут занесены данные при выполнении скрипта;
	\item Сохранять BLOB файлы по отдельности -- Если включено, то блобы будут сохраняться в указанную папку отдельными файлами;
	\item Имя таблицы -- Имя таблицы, которое будет использовано при генерации скрипта.
\end{itemize}

\newpage

\section{Экспорт метаданных БД}\label{sec:export_db_meta}

Инструмент \ttt{Экспорт метаданных БД} извлекает метаданные выбранной БД в \ttt{SQL-скрипт}, выполнение которого позволяет создать дубликат этой базы.

Открыть панель экспорта метаданных можно двумя способами:
\begin{itemize}
	\item На панели инструменов нажпите на кнопку \ttt{Извлечь метаданные БД в SQL скрипт}:
		\begin{figure}[H]
			\centering
			\includegraphics[width = 0.5\linewidth]{img/export_db_meta_button.png}
		\end{figure}
	\item На панели \ttt{Подключения} кникните правой кнопкой мыши по названию подключения и выберите пункт \ttt{Экспортировать метаданные}:
		\begin{figure}[H]
			\centering
			\includegraphics[width = 0.5\linewidth]{img/export_db_meta1.png}
		\end{figure}
\end{itemize}

Панель экспорта метаданных выглядит следующим образом:
\begin{figure}[H]
	\centering
	\includegraphics[width = 1\linewidth]{img/export_db_meta.png}
	\caption{Панель экспорта метаданных}
\end{figure}
 
Блок \textit{Соединения}:
 \begin{itemize}
 	\item База-образец -- база данных, метаданные которой нужно извлечь;
 	\item Извлечь -- кнопка, которая запускает экспорт метаданных.
 \end{itemize}	
 
Блок \textit{Атрибуты} представляет собой список элементов базы данных, которые нужно/не нужно учитывать при экспорте.

Блок \textit{Параметры} представляет собой список условий, влияющих на извлечение метаданных и формирование \ttt{SQL-скрипта}, создающего выбранную базу данных:
 \begin{itemize}
 	\item Игнорировать комментарии - учитывать ли комментарии при экспорте;
 	\item Игнорировать вычисляемые столбцы -- учитывать ли вычисляемые столбцы при экспорте; 
 	\item Игнорировать позиции столбцов -- учитывать ли позиции столбцов при экспорте;
 	\item Игнорировать \ttt{PK/FK/UK/CK} -- учитывать ли ограничения при экспорте.
 \end{itemize}
\newpage

Посмотреть результаты экспорта метаданных можно во вкладках \ttt{Вывод}, \ttt{Обзор} и \ttt{SQL}.

Во вкладке \ttt{Вывод} перечислены элементы, метаданные которых были извлечены:
\begin{figure}[H]
	\centering
	\includegraphics[width = 1\linewidth]{img/export_db_meta_output.png}
	\caption{Вкладка Вывод}
\end{figure}

Во вкладке \ttt{Обзор} отображены извлечённые элементы, которые будут созданы при выполнении сгенерированного скрипта:
\begin{figure}[H]
	\centering
	\includegraphics[width = 1\linewidth]{img/export_db_meta_view.png}
	\caption{Вкладка обзор}
\end{figure}

Во вкладке \ttt{SQL} размещён сгенерированный \ttt{SQL-скрипт}:
\begin{figure}[H]
	\centering
	\includegraphics[width = 1\linewidth]{img/export_db_meta_sql.png}
	\caption{Вкладка SQL}
\end{figure}

Кнопка \ttt{Сохранить скрипт} позволяет сохранить сгенерированный скрипт в файл.

Кнопка \ttt{Выполнить скрипт} открывает сгенерированный скрипт в редакторе запросов.

\section{Выполнение SQL скриптов}\label{sec:execute_scripts}

Данная опция позволяет задать файл с SQL запросами для выполнения.

\begin{enumerate}[leftmargin=39pt]
	\item Выберите пункт меню \textit{База данных}$\rightarrow$\textit{Выполнить SQL скрипт}. На центральной панели приложения откроется новая вкладка.
	\begin{figure}[H]
		\flushright
		\includegraphics[width = 0.92\linewidth]{img/execute_input.png}
	\end{figure}
	\item Выберите открытое соединение.
	\item Выберите действие в случае возникновения ошибки: остановиться или продолжить.
	\item Укажите путь к файлу с SQL запросами, если вы хотите выполнить скрипт из файла.
	\item Поставьте галочку, если вы хотите видеть подробный вывод о ходе выполнения.
	\item Нажмите кнопку \textit{Start}.
	\item Сделайте \ttt{commit} или \ttt{rollback} транзакции, нажав на соответствующие кнопки.
\end{enumerate}

\newpage

\section{Модуль сравнения баз данных}\label{sec:compareDB_module}

Инструмент \textit{Сравнить базы данных} сравнивает две базы данных и генерирует \ttt{SQL-скрипт}, позволяющий сделать две базы идентичными. 

\begin{redimportant}
	В результате работы модуля у сравниваемых баз будет идентичная структура, но не содержание.
\end{redimportant}	

В меню \textit{Инструменты} выберите \textit{Сравнить базы данных}. Откроется вкладка \textit{Сравнение баз данных}.

\begin{figure}[H]
	\centering
	\includegraphics[width = 1\linewidth]{img/compareDB_module.png}
	\caption{Модуль сравнения баз данных}
\end{figure}

Блок \textit{Соединения}:
\begin{itemize}
	\item База-образец -- база данных, к состоянию которой нужно привести целевую базу;
	\item Целевая база -- база данных, к которой будут применяться изменения.
\end{itemize}	

Блок \textit{Атрибуты} представляет собой список элементов базы данных, которые нужно/не нужно учитывать при сравнении.

Блок \textit{Параметры} представляет собой список условий, влияющих на сравнение баз и формирование \ttt{SQL-скрипта}, приводящего структуру целевой базы к состоянию базы-образца:
\begin{itemize}
	\item Поиск \ttt{CREATE} элементов -- если установлена галочка, то в \ttt{SQL-скрипт} будут добавлены \ttt{SQL-запросы} для создания объектов, которые отсутствуют в целевой базе, но есть в базе-образце;
	\item Поиск \ttt{ALTER} элементов -- если установлена галочка, то в \ttt{SQL-скрипт} будут добавлены \ttt{SQL-запросы} для изменения объектов с одинаковыми именами и типами данных;
	\item Поиск \ttt{DROP} элементов -- если установлена галочка, то в \ttt{SQL-скрипт} будут добавлены \ttt{SQL-запросы} для удаления объектов, которые есть в целевой базе, но отсутствуют в базе-образце;
	\item Игнорировать комментарии - учитывать ли комментарии при сравнении баз;
	\item Игнорировать вычисляемые столбцы -- учитывать ли вычисляемые столбцы при сравнении баз; 
	\item Игнорировать \ttt{PK/FK/UK/CK} -- учитывать ли ограничения при сравнении баз.
\end{itemize}

Найденные различия будут записаны во вкладку \textit{Вывод}. Во вкладке \textit{SQL} будет сформирован скрипт, приводящий структуру целевой базы к структуре базы-образца.

Пример сравнения баз данных:
\begin{figure}[H]
	\centering
	\includegraphics[width = 1\linewidth]{img/compareDB_module_example.png}
	\caption{Найденные различия}
\end{figure}

\begin{figure}[H]
	\centering
	\includegraphics[width = 1\linewidth]{img/compareDB_module_exampleSQL.png}
	\caption{Сформированный SQL-скрипт}
\end{figure}

\newpage

\section{Профайлер}\label{sec:profiler}

Инструмент \textit{Профайлер} позволяет пользователям измерять затраты на производительность \ttt{SQL} и \ttt{PSQL} кода. 

Чтобы открыть панель профайлера, в меню \textit{Инструменты} выберите \textit{Профайлер}. 

\begin{figure}[H]
	\centering
	\includegraphics[width = 0.5\linewidth]{img/open_profiler_module.png}
	\caption{Открытие панели профайлера}
\end{figure}

После этого откроется панель профайлера, состоящая из выпадающих списков для выбора базы данных и соединения, набора кнопок для управления сессией профайлера, результирующей таблицы и панели для выбора формата отображения дерева вызовов.

\begin{figure}[H]
	\centering
	\includegraphics[width = 1\linewidth]{img/profiler_expanded_data_display.png}
	\caption{Панель профайлера}
\end{figure}

Кнопки, управляющие сессией профайлера:
\begin{itemize}
\item \ttt{Старт} -- запускает новый сеанс профайлера;
\item \ttt{Пауза} -- приостанавливает текущий сеанс профайлера;
\item \ttt{Возобновить} -- возобновляет приостановленный сеанс профайлера;
\item \ttt{Стоп} -- завершает текущий сеанс профайлера;
\item \ttt{Отменить} -- удаляет все данные сеанса, присутствующие в плагине профайлера (данные, которые были обновлены, не удаляются);
\item \ttt{Очистить} -- удаляет все сеансы профайлера из памяти, не cбрасывая их данных. Если есть активный сеанс, он завершается. Также очищает дерево вызовов, которое было построено профайлером.
\end{itemize}

Форматы отображения результата: 
\begin{itemize}
\item \ttt{Компактное отображение} -- Повторяющиеся процессы внутри общего родительского будут объединены в один. Является значением по умолчанию.
 \begin{figure}[H]
	\centering
	\includegraphics[width = 0.8\linewidth]{img/profiler_compact_data_display.png}
	\caption{Компактное отображение данных}
\end{figure}
\item \ttt{Стандартное отображение} -- Каждый процесс отображается в отдельном узле. Повторяющиеся процессы не объединяются.
 \begin{figure}[H]
	\centering
	\includegraphics[width = 0.8\linewidth]{img/profiler_standard_data_display.png}
	\caption{Стандартное отображение данных}
\end{figure}
\item \ttt{Расширенное отображение} -- В дополнение к стандартному отображению добавлена PSQL статистика для внутренних объектов процедур и функций.
 \begin{figure}[H]
	\centering
	\includegraphics[width = 0.8\linewidth]{img/profiler_expanded_data_display.png}
	\caption{Расширенное отображение данных}
\end{figure}
\item \ttt{Округлять значения} - Если общее или среднее время больше 1000000ns, то оно будет переводиться в большую единицу измерения, пока значение не станет меньше 1000000. По умолчанию включено.
\end{itemize}
В компактном и стандартном отображении для каждого не последнего узла (за исключением \ttt{ROOT NODE} - корневого узла) есть узел \ttt{Собственное время}, показывающий затраченное время без учета дочерних процессов.

В результирующей таблице отображается информация, собранная профайлером, а именно:
\begin{itemize}
	\item Имя процесса или \ttt{SQL-код};
	\item Затраченное на процесс время в наносекундах (с учетом дочерних процессов) и процент времени от родительского процесса;
	\item Среднее затраченное на процесс время в наносекундах (с учетом дочерних процессов) для повторяющихся процессов, объединенных в один узел; 
	\item Количество вызовов повторяющихся процессов.
\end{itemize}
При выборе узла будет открыто всплывающее окно.
\begin{figure}[H]
	\centering
	\includegraphics[width = 0.4\linewidth]{img/profiler_nodes.png}
\end{figure}
\begin{itemize}
	\item \ttt{Копия} -- Скопировать содержимое выбранного узла в буфер обмена.
	\item \ttt{Просмотреть} -- Открыть содержимое узла в отдельном окне. Если это процедура или функция, будет отображен ее код.
	\begin{figure}[H]
		\centering
		\includegraphics[width = 0.5\linewidth]{img/profiler_node_view.png}
	\end{figure}
	\item \ttt{Информация} -- Открыть в отдельном окне всю информацию, которую собрал профайлер для выбранного узла.
	\begin{figure}[H]
		\centering
		\includegraphics[width = 0.6\linewidth]{img/profiler_node_info.png}
	\end{figure}
\end{itemize}
Сессию профайлера можно запустить для одного запроса из \ttt{редактора запросов} кликом по кнопке "Выполнить в профайлере". В этом случае будет выполнено следующее:
\begin{enumerate}
	\item Запустится сессия профайлера.
	\item Выполнится находящийся в редакторе запрос.
	\item Завершится сессия профайлера.
	\item Отобразится панель профайлера с собранной информацией.
\end{enumerate}	

\newpage

\section{Генератор тестовых данных}\label{sec:test_data_generator}

Инструмент \textit{Генератор тестовых данных} предназначен для быстрого и удобного заполнения полей таблицы большим объемом данных. Его можно открыть через пункт меню \textit{Инструменты}.
На главной панели откроется вкладка с названием \textit{Генератор тестовых данных}. 

Прежде чем воспользоваться генератором, убедитесь, что Вы подключены к базе данных и в ней существует хотя бы одна таблица.  

\begin{enumerate}[leftmargin=39pt]
	\item В поле \textit{Подключение} выберите базу данных, для которой должны быть созданы тестовые данные (если подключено более одной базы данных). 

	\item Далее выберите таблицу. Поля и их типы вы увидите перед собой. Напротив поля поставьте галочку, если Вы хотите, чтобы для него были сгенерированы данные. Иначе это поле будет заполнено значениями \ttt{NULL}.

	\item Напишите количество записей, которые будут сгенерированы. Генерация больших объемов данных может занять некоторое время. 
	
	\item В поле \textit{Фиксация после} введите число строк, после вставки которых будет совершаться коммит.
	
	\item При генерации могут возникать ошибки. По умолчанию они не пишутся в лог-файл, чтобы не замедлять работу генератора. Поставьте соответствующую галочку, если вы хотите, чтобы все ошибки фиксировались в логе.
	
	\item Если в процессе генерации возникает ошибка для некоторых записей (например, из-за ограничений столбца), то по умолчанию генератор не остановится и продолжит свою работу. Вы можете изменить это поведение, поставив напротив поля \textit{Остановить при ошибке} галочку.

	\item С правой стороны окна можно выбрать метод генерации отдельно для каждого поля таблицы. 
Предлагается выбрать один из следующих методов:
\begin{itemize}[leftmargin=10pt]
	\item \textbf{\textit{Случайно}}
	
	В зависимости от типа поля настраиваются различные параметры генерации:
	\begin{itemize}[leftmargin=5pt]
		\item[] \textit{Целочисленные типы}: 
		\begin{itemize}
			\item минимальное и максимальное значение диапазона;
		\end{itemize}	
		\item[] \textit{Числа с плавающей и фиксированной точкой}:  
		\begin{itemize}
			\item  минимальное и максимальное значение диапазона; 
			\item  число знаков после запятой;
		\end{itemize}		
		\item[] \textit{Тип даты и времени}: 
		\begin{itemize}
			\item  минимальное и максимальное значение диапазона даты и времени;
		\end{itemize}		
		\item[] \textit{Строковые типы}:
		\begin{itemize}
			\item минимальная и максимальная длина строки; 
			\item набор символов;
		\end{itemize}		
		\item[] \textit{Тип BLOB}:  
		\begin{itemize}
			\item  минимальная и максимальная длина строки; 
			\item  минимальное и максимальное значение, которое могут принимать байты.
		\end{itemize}	
	\end{itemize}	
	\item \textbf{\textit{Получить из другой таблицы}}
	
	В этом методе нужно выбрать таблицу, столбец и количество записей. Записи из таблицы выбираются случайно в указанном количестве, далее из этого списка значений заполняется основная таблица. 
	\item \textbf{\textit{Получить из списка}}
	
	Значения списка формируют содержимое поля.
	
	Список должен состоять из элементов соответствующего типа данных, разделенных специальным символом. Разделителем может выступать любой одиночный символ или escape-последовательность, начинающаяся с обратной косой черты (\rr{\textbackslash}). Сам список можно ввести вручную в предназначенном для этого поле, либо загрузить из файла. 

	\item \textbf{\textit{Автоинкремент}}
	
	Для автоинкремента настраивается начальное значение, шаг и направление движения шага (в сторону увеличения или уменьшения).
	
	\end{itemize}
	\item Наконец, нажмите на кнопку \textit{Начать генерацию}. 
\end{enumerate}

Сгенерированные данные можно посмотреть в редакторе соответствующей таблицы на странице \textit{Данные}:
\begin{figure}[H]
	\centering
	\includegraphics[width = 1\linewidth]{img/test_data_generator_2.png}
\end{figure}

\newpage

\section{ER-диаграмма}\label{sec:erd}

Диаграмма «сущность-связь» (англ. entity-relationship diagram, ERD, ER-диаграмма) представляет собой графическую визуализацию схемы базы данных. 
Ее используют при высокоуровневом проектировании баз данных. 

ER-диаграмма визуализирует табличные связи, обеспечивает полный доступ к структурам таблиц и ограничениям и позволяет создавать сценарии \ttt{CREATE TABLE} из самой диаграммы.
\begin{redremark}
	Изменения, внесенные в структуры таблиц и связанные с ними ограничения, не распространяются на базу данных. 
\end{redremark}

В языке ER-диаграммы таблица базы данных есть сущность, т.е. класс однотипных объектов, информация о которых должна быть учтена в модели.

Каждая таблица на диаграмме изображается в виде прямоугольника с названием таблицы:

\begin{figure}[H]
	\centering
	\includegraphics[width = 0.18\linewidth]{img/erd1.png}
\end{figure}

Между таблицами может существовать связь посредством внешний ключей. Связь --- это некоторая ассоциация между двумя сущностями. Одна сущность может быть связана с другой сущностью или сама с собою. Связи позволяют по одной сущности находить другие сущности, связанные с нею.

Связи внешними ключами между таблицами изображены в виде стрелок:

\begin{figure}[H]
	\centering
	\includegraphics[width = 0.45\linewidth]{img/erd2.png}
\end{figure}

\anonsubsection{Панель инструментов}\label{sec:81}

Следующие команды доступны на панели инструментов ERD:

\begin{longtable}[c]{|m{5mm}|m{9cm}|>{\ttfamily}m{4cm}|}
	\hline
	&
	\centering\bfseries Описание &
	\centering\arraybslash\normalfont\bfseries Горячие клавиши\\\hline
	\endfirsthead
	\hline
	&
	\centering\bfseries Описание &
	\centering\arraybslash\normalfont\bfseries Горячие клавиши\\\hline
	\endhead
	\begin{tikzpicture}
	\pgftext{\includegraphics[width=11pt]{img/NewTable16.png}} at (0pt,0pt)
	\end{tikzpicture} & Создать новую таблицу & \\\hline
	\begin{tikzpicture}
	\pgftext{\includegraphics[width=11pt]{img/AddTable16.png}} at (0pt,0pt)
	\end{tikzpicture} & Добавить таблицу из существующей базы данных & \\\hline
	\begin{tikzpicture}
	\pgftext{\includegraphics[width=11pt]{img/TableRelationship16.png}} at (0pt,0pt)
	\end{tikzpicture} & Создать табличную связь посредством внешних ключей & \\\hline
	\begin{tikzpicture}
	\pgftext{\includegraphics[width=11pt]{img/TableRelationshipDelete16.png}} at (0pt,0pt)
	\end{tikzpicture} & Удалить связь между выбранными таблицами & \\\hline
	\begin{tikzpicture}
	\pgftext{\includegraphics[width=11pt]{img/DropTable16.png}} at (0pt,0pt)
	\end{tikzpicture} & Удалить выбранные таблицы с диаграммы & \\\hline
	\begin{tikzpicture}
	\pgftext{\includegraphics[width=11pt]{img/CreateScripts16.png}} at (0pt,0pt)
	\end{tikzpicture} & Извлечь метаданные БД в SQL скрипт & \\\hline
	\begin{tikzpicture}
	\pgftext{\includegraphics[width=11pt]{img/ErdTitle16.png}} at (0pt,0pt)
	\end{tikzpicture} & Добавить описательный заголовок диаграммы & \\\hline
	\begin{tikzpicture}
	\pgftext{\includegraphics[width=11pt]{img/FontStyle16.png}} at (0pt,0pt)
	\end{tikzpicture} & Изменить шрифт всех таблиц & \\\hline
	\begin{tikzpicture}
	\pgftext{\includegraphics[width=11pt]{img/LineStyle16.png}} at (0pt,0pt)
	\end{tikzpicture} & Изменить стиль линий связи & \\\hline
	\begin{tikzpicture}
	\pgftext{\includegraphics[width=11pt]{img/ErdForeground16.png}} at (0pt,0pt)
	\end{tikzpicture} & Изменить цвет фона содержимого выбранных таблиц & \\\hline
	\begin{tikzpicture}
	\pgftext{\includegraphics[width=11pt]{img/ErdBackground16.png}} at (0pt,0pt)
	\end{tikzpicture} & Изменить цвет холста & \\\hline
	\begin{tikzpicture}
	\pgftext{\includegraphics[width=11pt]{img/ZoomOut16.png}} at (0pt,0pt)
	\end{tikzpicture} & Уменьшить размер & \\\hline
	\begin{tikzpicture}
	\pgftext{\includegraphics[width=11pt]{img/ZoomIn16.png}} at (0pt,0pt)
	\end{tikzpicture} & Увеличить размер & \\\hline
\end{longtable}

\anonsubsection{Создание новой диаграммы}

Чтобы создать новую ER-диаграмму, выберите пункт меню \textit{База данных}$\rightarrow$\textit{Построить ER диаграмму}$\rightarrow$\textit{Создать ER-диаграмму}.
На главной панели откроется вкладка с названием \textit{Entity Relationship Diagram -- erd<N>.eqd} на котором будет изображен пустой холст, не содержащий таблиц и связей. 

На диаграмме можно создать новые таблицы (\begin{tikzpicture}
\pgftext{\includegraphics[width=11pt]{img/NewTable16.png}} at (0pt,0pt)
\end{tikzpicture} ) или добавить существующие (\begin{tikzpicture}
\pgftext{\includegraphics[width=11pt]{img/AddTable16.png}} at (0pt,0pt)
\end{tikzpicture}) из открытого подключения к базе данных.
Об этих и других функциях рассказано в подразделе \hyperref[sec:84]{Управление таблицами}.

\anonsubsection{Диаграмма существующей базы данных}
Существует возможность наоборот построить диаграмму по имеющейся базе данных. Для этого воспользуйтесь пунктом меню \textit{База данных}$\rightarrow$\textit{Построить ER диаграмму}$\rightarrow$\textit{ER-диаграмма базы данных}.

Откроется диалоговое окно, в котором:
\begin{enumerate}
	\item Выберите открытое соединение
	\item Выберите таблицы для включения в ERD с помощью кнопок-стрелок
	\item Нажмите кнопку \textit{Generate}
\end{enumerate}
На холсте будут изображены выбранные сущности и имеющиеся между ними связи.

\anonsubsection{Управление таблицами}\label{sec:84}
\anonsubsubsection{Создание новых таблиц}
\begin{enumerate}
	\item Нажмите кнопку \begin{tikzpicture}
	\pgftext{\includegraphics[width=11pt]{img/NewTable16.png}} at (0pt,0pt)
	\end{tikzpicture} на панели инструментов ERD
	\item Выберите открытое соединение, которое будет использовано в этой таблице
	\item Введите имя таблицы, может содержать до 31 байта.
	\item Создайте поля таблицы, заполнив соответствующую таблицу (см. \hyperref[tab:2]{подробное описание})
	\item Заполните табличку с ограничениями таблицы (см. \hyperref[tab:3]{подробное описание})
	\item Нажмите кнопку \textit{Создать}
\end{enumerate}

\anonsubsubsection{Добавление таблицы из существующей базы данных}
\begin{enumerate}
	\item Нажмите на кнопку \begin{tikzpicture}
	\pgftext{\includegraphics[width=11pt]{img/AddTable16.png}} at (0pt,0pt)
	\end{tikzpicture} на панели инструментов ERD
	\item Выберите открытое соединение
	\item В поле \ttt{Available Tables} будет подгружены все таблицы, доступные в выбранном соединении. Необходимые для построения диаграммы таблицы перенесите в окно \ttt{Selected Tables} с помощью кнопок \begin{tikzpicture}
	\pgftext{\includegraphics[width=11pt]{img/Forward16.png}} at (0pt,0pt)
	\end{tikzpicture},
	\begin{tikzpicture}
	\pgftext{\includegraphics[width=11pt]{img/SelectAll16.png}} at (0pt,0pt)
	\end{tikzpicture}
	\item Нажмите на кнопку \textit{Добавить}
\end{enumerate}

\anonsubsubsection{Изменение таблиц}
Структуру таблиц можно просмотреть более подробно, дважды щелкнув по интересующей таблице. Покажется диалоговое окно, которое позволит изменять столбцы и ограничения таблицы. Таблицы могут быть переименованы, связаны с другим открытым соединением, а также столбцы таблицы могут быть переупорядочены.

\anonsubsubsection{Создание новых табличных связей}

Схема создания ограничения внешних ключей между выбранными объектами:
\begin{enumerate}
	\item Нажмите кнопку 
	\begin{tikzpicture}
	\pgftext{\includegraphics[width=11pt]{img/TableRelationship16.png}} at (0pt,0pt)
	\end{tikzpicture}
	на панели инструментов ERD
	\item Введите (уникальное) имя для этого ограничения
	\item Выберите таблицу и столбец, который ссылается внешним ключом
	\item Выберите таблицу и столбец, который на который ссылается внешний ключ
	\item нажмите кнопку «Создать»
\end{enumerate}
Созданная связь будет показана с использованием одной стрелки между таблицами.

\anonsubsubsection{Удаление табличных связей}
Удаление связей между таблицами приведет к удалению ограничения внешнего ключа. Это будет отражено как в диалоговом окне описания таблицы, так и в любом сгенерированном SQL.
\begin{enumerate}
	\item Выделите на холсте две таблицы, которые содержат ограничение внешнего ключа, с помощью левой кнопки мыши и клавиши Ctrl
	\item Нажмите кнопку 
	\begin{tikzpicture}
	\pgftext{\includegraphics[width=11pt]{img/DropTable16.png}} at (0pt,0pt)
	\end{tikzpicture} 
	на панели инструментов ERD. Откроется диалоговое окно «Delete Table Relationship», в таблице которого перечислены все ограничения для таблиц. 
	\item Поставьте галочку на нужном ограничении
	\item Нажмите кнопку «Удалить»
\end{enumerate}
Видимая ранее стрелка, связывающая таблицы, будет удалена.

\anonsubsubsection{Удаление таблиц}
Удаление таблиц(ы) приведет к удалению таблиц(ы) только из ER-диаграммы.
\begin{enumerate}
	\item Выделите на холсте таблицу для удаления при помощи мыши или выделите несколько таблиц, используя клавишу Ctrl
	\item Нажмите кнопку \begin{tikzpicture}
	\pgftext{\includegraphics[width=11pt]{img/DropTable16.png}} at (0pt,0pt)
	\end{tikzpicture}
\end{enumerate}

\anonsubsubsection{Генерирование SQL скриптов}

Для всех таблиц, изображенных на диаграмме, могут быть сгенерированы SQL скрипты создания этих таблиц (оператор \ttt{CREATE TABLE}). При этом вызовы в базу данных не производятся. 
\begin{enumerate}
	\item Нажмите на кнопку 
	\begin{tikzpicture}
	\pgftext{\includegraphics[width=11pt]{img/CreateScripts16.png}} at (0pt,0pt)
	\end{tikzpicture} на панели инструментов ERD. Появляется диалоговое окно «Generate Scripts». 
	\item Выберите нужные таблицы из списка «Available Tables», используя кнопки \begin{tikzpicture}
	\pgftext{\includegraphics[width=11pt]{img/Forward16.png}} at (0pt,0pt)
	\end{tikzpicture},
	\begin{tikzpicture}
	\pgftext{\includegraphics[width=11pt]{img/SelectAll16.png}} at (0pt,0pt)
	\end{tikzpicture}
	\item Укажите путь к файлу, в котором будет сохранен сгенерированный SQL-скрипт
	\item Поставьте галочку, если в скрипт следует включить любые ограничения таблицы (на выбор-или вставить в оператор \ttt{CREATE TABLE}, или сгенерировать отдельно оператор \ttt{ALTER TABLE} с добавлением этих ограничений)
	\item Нажмите кнопку «Generate»
\end{enumerate}

\anonsubsection{Манипулирование диаграммой}
Диаграмма ERD позволяет выполнять несколько настроек, включая позиционирование каждой таблицы на холсте, выбор цвета фона и таблиц, выбор шрифта и стиля линий. Большинство этих функций доступны через панель инструментов ERD, описанную \hyperref[sec:81]{выше}. Остальные доступны через всплывающее меню по нажатию правой кнопки мыши на холсте.

\anonsubsubsection{Перемещение объектов}

Панель ERD дает возможность перемещать выбранные таблицы по холсту для получения более четкого представления структуры базы данных. Чтобы переместить таблицу, щелкните левой кнопкой мыши по заголовку таблицы и перетащите ее в нужное место. Связи между таблицами останутся неизменными.

\anonsubsubsection{Заголовок диаграммы}
\begin{enumerate}
	\item Нажмите кнопку \begin{tikzpicture}
	\pgftext{\includegraphics[width=11pt]{img/ErdTitle16.png}} at (0pt,0pt)
	\end{tikzpicture} на панели инструментов ERD
	\item Заполните соответствующие поля для текущей диаграммы: название, дату, автора, описание, базу данных и пересмотр.
	\item Нажмите кнопку «Добавить»
\end{enumerate}

Заголовок может перемещаться, как любой другой компонент, в нужное место путем перетаскивания. Чтобы отредактировать содержимое поля заголовка после первоначального создания, дважды щелкните в любом месте внутри него, чтобы отобразить диалог заголовка ERD и внести необходимые изменения.

\anonsubsubsection{Отображение полей таблицы}

На диаграмме таблицы могут быть показаны со всеми полями или только теми, которые имеют такие ограничения, как первичный и внешний ключи. Последний вариант может быть более предпочтительным в больших структурах базы данных, где пространство ограничено. Вид можно изменить через всплывающее меню (щелкнуть правой кнопкой мыши в области холста) \textit{View}$\rightarrow$\textit{Display Referenced keys only}.

\anonsubsection{Сохрание и экспорт}

Сгенерированные ERD могут быть сохранены для последующего редактирования или экспортированы в выбранные форматы изображений. Сохранение файла ERD по умолчанию происходит с помощью простой сериализации объектов Java с расширением файла «eqd». Сохраненные ERD могут быть повторно открыты позже, используя команду «Открыть файл» из главного меню или панели инструментов.
Ниже описывается процесс сохранения и экспорта.

\begin{enumerate}
	\item На главной панели меню Red Expert выберите пункт \textit{Файл}$\rightarrow$\textit{Сохранить} или выберите кнопку
	\begin{tikzpicture}
	\pgftext{\includegraphics[width=11pt]{img/save16.png}} at (0pt,0pt)
	\end{tikzpicture} 
	на главной панели инструментов
	\item Выберите файл или формат изображения для сохранения. Доступные форматы изображений включают: JPEG, GIF, PNG, SVG
	\item Выберите любые параметры формата изображения
	\item Выберите путь к файлу для сохранения
	\item Нажмите кнопку \textit{Сохранить}
\end{enumerate}

\newpage

\section{Черновик}\label{sec:draft}

Черновик - это простой компонент текстового редактора, не поддерживающий форматирование, доступный в виде вкладки на центральной панели. Рекомендуется для временного отображения или редактирования текста по мере необходимости.

Открыть черновик можно с помощью пункта меню \textit{Инструменты}$\rightarrow$\textit{Черновик}. 

Данный компонент имеет панель инструментов:
\begin{longtable}[c]{|m{5mm}|m{9cm}|>{\ttfamily}m{4cm}|}
	\hline
	&
	\centering\bfseries Описание &
	\centering\arraybslash\normalfont\bfseries Горячие клавиши\\\hline
	\endfirsthead
	\hline
	&
	\centering\bfseries Описание &
	\centering\arraybslash\normalfont\bfseries Горячие клавиши\\\hline
	\endhead
	\begin{tikzpicture}
	\pgftext{\includegraphics[width=12pt]{img/NewScratchPad16.png}} at (0pt,0pt)
	\end{tikzpicture} & Открыть новый черновик & \\\hline
	\begin{tikzpicture}
	\pgftext{\includegraphics[width=12pt]{img/ScratchToEditor16.png}} at (0pt,0pt)
	\end{tikzpicture} & Открыть редактор запросов с содержимым черновика & \\\hline
	\begin{tikzpicture}
	\pgftext{\includegraphics[width=12pt]{img/Delete16.png}} at (0pt,0pt)
	\end{tikzpicture} & Очистить черновик & \\\hline
\end{longtable}	

\newpage

\section{Пользовательские ключевые слова}

Ключевые слова, добавленные пользователем Red Expert, распознаются в редакторе запросов, подсвечиваются и присутствуют в всплывающей подсказке (по нажатию \ttt{Ctrl+Space}).

\begin{enumerate}[leftmargin=39pt]
	\item Выберите пункт меню \textit{Инструменты}$\rightarrow$\textit{Пользовательские ключевые слова} или нажмите \begin{tikzpicture}
	\pgftext{\includegraphics[width=12pt]{img/Keywords16.png}} at (0pt,0pt)
	\end{tikzpicture} на главной панели инструментов. Откроется диалоговое окно
	\begin{figure}[H]
		\centering
		\includegraphics[width = 0.7\linewidth]{img/user_keywords.png}
	\end{figure}
	В таблице слева перечислены добавленные пользователем ранее ключевые слова. В таблице справа перечислены ключевые слова стандарта SQL-92.
	\item Введите новое ключевое слово в поле \textit{New Keyword}
	\item Нажмите кнопку \textit{Add}, чтобы включить новое ключевое слово в таблицу слева
	\item Чтобы удалить пользовательское ключевое слово из списка, выделите ячейку таблицы и нажмите кнопку \begin{tikzpicture}
	\pgftext{\includegraphics[width=12pt]{img/Delete16.png}} at (0pt,0pt)
	\end{tikzpicture}
	\item Нажмите кнопку \textit{Сохранить}, чтобы изменения вступили в силу.
\end{enumerate}

\begin{redremark}
	Добавить новые ключевые слова можно и из Редактора Запросов просто поставив курсор на нужном слове и выбрав пункт всплывающего меню \textit{Добавить пользовательские ключевые слова}.
\end{redremark}	

\newpage

\section{Менеджер пользователей}\label{sec:user_manager}
С помощью Менеджера пользователей можно управлять пользователями базы данных: добавлять, редактировать и удалять. 
\begin{enumerate}[leftmargin=39pt]
	\item Выберите пункт меню \textit{Инструменты}$\rightarrow$\textit{Менеджер пользователей} или нажмите \begin{tikzpicture}
	\pgftext{\includegraphics[width=12pt]{img/user_manager_16.png}} at (0pt,0pt)
	\end{tikzpicture} на главной панели инструментов. На центральной панели откроется новая вкладка:
	\begin{figure}[H]
		\flushright
		\includegraphics[width = 0.92\linewidth]{img/user_manager.png}
	\end{figure}
	\item Выберите подключение к базе данных (для получения адреса сервера) или укажите свой адрес сервера.
	\item В таблице будут отображены все пользователи, зарегистрированные на сервере.
\end{enumerate}	

\anonsubsection{Добавление пользователя}

\begin{enumerate}[leftmargin=39pt]
	\item Нажмите кнопку \textit{Добавить} на вкладке Менеджера пользователей
	\item На центральной панели откроется новая вкладка. Введите информацию о пользователе:
	\label{tab:adduser}
	\begin{longtable}[r]{|>{\itshape}m{3.2cm}|m{11cm}|}
		\hline
		\centering\normalfont\bfseries Параметр &
		\centering\arraybslash\bfseries Описание\\\hline
		\endfirsthead
		\hline
		\centering\normalfont\bfseries Параметр &
		\centering\arraybslash\bfseries Описание\\\hline
		\endhead
		Имя пользователя & Логин, с которым пользователь будет зарегистрирован. Обязательный параметр. \\\hline
		Пароль & Пароль пользователя. Чувствительный к регистру. Обязательный параметр. \\\hline
		Фамилия & Вспомогательная информация: фамилия пользователя.\\\hline
		Имя & Вспомогательная информация: имя пользователя. \\\hline
		Отчество & Вспомогательная информация: отчество пользователя. \\\hline
		Плагин & Имя плагина управления пользователями из списка параметра UserManager в файле конфигурации firebird.conf (для Firebird, Ред Базы Данных). По умолчанию используется тот плагин, который был указан первым в списке параметра UserManager. \\\hline
		Активный & Включите этот атрибут и пользователь будет создан в активном состоянии. Иначе пользователь будет	создан в "неактивном состоянии"\ , т.е. подключиться с его учётной записью будет невозможно.\\\hline
		Администратор & Если опция включена, новая учётная запись пользователя создаётся с правами роли \ttt{RDB\$ADMIN} в текущей базе данных безопасности. Это позволяет вновь созданному пользователю управлять учётными записями пользователей, но не даёт ему специальных полномочий в обычных базах данных.\\\hline
		Комментарий & Примечание к новой учетной записи. Необязательно.\\\hline
		Метки & Можно добавить неограниченное количество пользовательских атрибутов (тегов). \\\hline
	\end{longtable} 	
	\item Нажмите кнопку \textit{OK} и создастся новая учетная запись.
\end{enumerate}	


\anonsubsection{Редактирование пользователя}

\begin{enumerate}[leftmargin=39pt]
	\item В таблице Менеджера пользователей выделите строку таблицы, соответствующую пользователю, информацию о котором нужно изменить.
	\item Нажмите кнопку \textit{Редактировать}.
	\item На центральной панели откроется новая вкладка. Отредактируйте необходимую информацию о пользователе (см. \hyperref[tab:adduser]{таблицу выше}). Логин пользователя и плагин, с которым он создан, изменять нельзя. 	
	\item Нажмите кнопку \textit{OK} и учетная запись изменится.
\end{enumerate}	

\anonsubsection{Назначение ролей}
Менеджер пользователей позволяет просматривать не только учетные записи, но и роли, хранящиеся в базе данных. Они перечислены на вкладке \textit{Роли}.

Кроме того, имеющиеся роли можно назначить другим пользователям и ролям:

\begin{enumerate}[leftmargin=39pt]
	\item Откройте вкладку \textit{Членство}. 
	\item Вы увидите таблицу, верхние ячейки которой соответствуют всем ролям базы данных, а левые -- ролям и пользователям.
	\begin{figure}[H]
		\flushright
		\includegraphics[width = 0.92\linewidth]{img/grant_role.png}
	\end{figure}
	\item Зеленая точка \begin{tikzpicture}
	\pgftext{\includegraphics[width=12pt]{img/grant.png}} at (0pt,0pt)
	\end{tikzpicture} говорит о том, что пользователю или роли (перечисленные слева таблицы) назначена одна из ролей. Точка \begin{tikzpicture}
	\pgftext{\includegraphics[width=12pt]{img/admin_option.png}} at (0pt,0pt)
	\end{tikzpicture} указывает на то, что пользователям или ролям, имеющим данную роль, можно назначить эту роль другим пользователям или ролям.
	
	\item С помощью соответствующих кнопок (или двойным кликом мыши) можно назначать и отбирать роли. Изменения тут же вступят в силу.
\end{enumerate}	

\newpage

\section{Менеджер прав}\label{sec:priv_manager}
Менеджер прав предоставляет графический интерфейс для назначения и отзыва привилегий.

После создания объекта, только пользователь создавший объект (его владелец) и администраторы имеют доступ к нему. Пользователю необходимы привилегии на каждый объект, к которому он должен получить доступ. Как правило, привилегии должны быть предоставлены явно пользователю владельцем объекта или администратором базы данных.

Привилегия включает тип DML доступа (\ttt{SELECT, INSERT, UPDATE, DELETE, EXECUTE, USAGE} и \ttt{REFERENCES}), имя объекта базы данных для которого предоставляется привилегия (таблица, представление, процедура) и имя объекта которому  предоставляется привилегия (пользователь, процедура, триггер, роль, представление). 

Все привилегии по доступу к объектам базы данных хранятся в самой базе, и не могут быть применены к любой другой базе данных.

\begin{enumerate}[leftmargin=39pt]
	\item Выберите пункт меню \textit{Инструменты}$\rightarrow$\textit{Менеджер прав} или нажмите \begin{tikzpicture}
	\pgftext{\includegraphics[width=12pt]{img/grant_manager_16.png}} at (0pt,0pt)
	\end{tikzpicture} на главной панели инструментов. На центральной панели откроется новая вкладка:
	\begin{figure}[H]
		\flushright
		\includegraphics[width = 0.92\linewidth]{img/priv_manager.png}
	\end{figure}
	\item Выберите открытое соединение с базой данных.
	\item На консоли слева выберите получателя привилегий.
	\item На центральной консоли в виде таблицы показаны привилегии получателя для объектов базы данных (таблицы, представления, процедуры) в виде зеленых точек: \begin{tikzpicture}
	\pgftext{\includegraphics[width=12pt]{img/grant.png}} at (0pt,0pt)
	\end{tikzpicture} и \begin{tikzpicture}
	\pgftext{\includegraphics[width=12pt]{img/admin_option.png}} at (0pt,0pt)
	\end{tikzpicture} (\ttt{WITH GRANT OPTION}).
	\item Двойным кликом мыши по ячейке таблицы можно назначить или отозвать ту или иную привилегию.
	\item Для удобного поиска объектов можно воспользоваться строкой фильтра. Фильтр чувствителен к регистру.
	\item Привилегии \ttt{UPDATE} и \ttt{REFERENCES} можно указать для конкретных столбцов таблицы.
	\item Любые изменения прав сразу вступают в силу.
\end{enumerate}	

\newpage

\section{Трейс Менеджер}\label{sec:trace_manager}

Трейс менеджер позволяют отслеживать и анализировать все, что происходит в базе данных в режиме реального времени. Он позволяют серверу отслеживать и записывать в лог-файлы такие события, как: соединения и отсоединения от БД (создания и удаления БД), операции DML и DDL, выполнение хранимых процедур и т.д. 

Открыть менеджер можно через пункт меню \textit{Инструменты}.

\anonsubsection{Включение аудита}

\begin{enumerate}[leftmargin=20pt]
	\item Для включения сервиса аудита необходимо заполнить следующие параметры:
	\begin{longtable}[r]{|>{\ttfamily}m{3cm}|m{11.5cm}|}
		\hline
		\centering\normalfont\bfseries Параметр &
		\centering\arraybslash\bfseries Описание\\\hline
		\endfirsthead
		\hline
		\centering\normalfont\bfseries Параметр &
		\centering\arraybslash\bfseries Описание\\\hline
		\endhead
		Подключения
		& Созданное подключение к базе данных \\\hline
		База данных
		& Путь к базе данных или ее алиас \\\hline
		Имя сессии 
		& Имя сессии аудита \\\hline
		Сервер
		& Имя хоста сервера базы данных или IP-адрес\\\hline
		Порт
		& Номер порта подключения к базе данных \\\hline
		Имя пользователя
		& Имя пользователя для включения трейса. Им может быть любой зарегистрированный пользователь. \\\hline
		Пароль
		& Пароль пользователя \\\hline
		Кодировка
		& Кодировка потока сообщений сервиса аудита \\\hline
	\end{longtable} 
	Вы можете выбрать какое-либо сохраненное подключение к базе данных и поля параметров будут заполнены автоматически на основе этого подключения. 
	\item Выберите файл у параметра \textit{Файл конфигурации}, чтобы указать конфигурационный файл с настройками аудита. Вы можете настроить список отслеживаемых событий, нажав на кнопку \textit{Новый файл конфигурации}.
	\item Поставьте галочку на параметре \textit{Лог в файл} и укажите путь к файлу для размещение логов трассировки. Иначе логгирование будет вестись только в таблицу Трейс Менеджера без сохранения в файл. 
	\item Нажмите кнопку \textit{Старт}, чтобы включить сессию пользовательской трассировки.
\end{enumerate}	

\begin{redremark}
	При закрытии вкладки Трейс Менеджера завершается сеанс текущей трассировки.
\end{redremark}

\anonsubsection{Параметры файла конфигурации}

В зависимости от версии сервера параметры настройки аудита несколько отличаются, поэтому из выпадающего списка выберите сервер, соответствующий базе данных.

Настройте параметры аудита и включите трассировку или сохраните их в файле конфигурации. 

Ниже дается краткое описание каждого параметра аудита.

\begin{longtable}[r]{|>{\ttfamily}m{5cm}|m{10.5cm}|}
	\hline
	\centering\normalfont\bfseries Параметр &
	\centering\arraybslash\bfseries Описание\\\hline
	\endfirsthead
	\hline
	\centering\normalfont\bfseries Параметр &
	\centering\arraybslash\bfseries Описание\\\hline
	\endhead
	\hline
	\rowcolor[gray]{.9}\multicolumn{2}{|m{15.75cm}|}{\bfseries Database}\\\hline 
	log\_security\_incidents &  События, связанные с нарушением безопасности сервера (инциденты безопасности) \\\hline
	log\_initfini/log\_init &  События начала/окончания ведения аудита БД в лог-файл  \\\hline
	log\_connections &  События присоединения/отсоединения к БД  \\\hline
	log\_transactions &  События начала и завершения транзакций  \\\hline
	log\_statement\_prepare &  События подготовки запросов к БД  \\\hline
	log\_statement\_free &  События освобождения запросов к БД \\\hline
	log\_statement\_start &  События начала выполнения запросов к БД  \\\hline
	log\_statement\_finish &  События окончания выполнения запросов к БД \\\hline
	log\_procedure\_start & События начала выполнения хранимых процедур \\\hline
	log\_procedure\_finish & События завершения выполнения хранимых процедур \\\hline
	log\_function\_start & События начала выполнения хранимых функций  \\\hline
	log\_function\_finish & События завершения выполнения хранимых функций  \\\hline
	log\_trigger\_start & События начала выполнения триггеров  \\\hline
	log\_trigger\_finish & События завершения выполнения триггеров \\\hline
	log\_context & События изменений значений контекстных переменных   \\\hline
	log\_errors & Запись об ошибках  \\\hline	
	log\_warnings & Запись о предупреждениях   \\\hline
	print\_plan & Печать планов запросов  \\\hline
	print\_perf & Печать статистики выполнения запросов  \\\hline
	log\_blr\_requests &  События прямого выполнения откомпилированных запросов во внутреннем представлении сервера - BLR \\\hline
	print\_blr &  Содержимое BLR-запросов будет преобразовываться в текстовое представление, иначе останется в двоичном виде  \\\hline
	log\_dyn\_requests & События прямого выполнения откомпилированных запросов на изменение метаданных (DDL) во внутреннем представлении сервера - DYN  \\\hline
	print\_dyn & Содержимое DYN-запросов будет преобразовываться в текстовое представление, иначе останется в двоичном виде   \\\hline
	log\_privilege\_changes & События, связанные с изменением правил разграничения доступа  \\\hline
	log\_changes\_only &  Запись только тех событий, которые изменяли данные в базе \\\hline
	print\_stack\_trace & Печать стека вызовов функций сервера при завершении события с ошибкой \\\hline
	log\_auth\_factors & События проверки предъявленных факторов аутентификации \\\hline
	log\_mandatory\_access & Аудит мандатного доступа \\\hline
	log\_record\_mandatory\_access & События о мандатном доступе к записям \\\hline
	log\_object\_relabeling & События смены метки объектов \\\hline
	log\_record\_relabeling & События смены метки записей \\\hline
	\rowcolor[gray]{.9}\multicolumn{2}{|m{15.75cm}|}{\bfseries Services}\\\hline 
	log\_services & События присоединения/отсоединения и старта сервиса \\\hline
	log\_service\_query & События запросов к сервису \\\hline
	\rowcolor[gray]{.9}\multicolumn{2}{|m{15.75cm}|}{\bfseries Filters}\\\hline
	include\_user\_filter & Регулярное выражение, которому должно соответствовать имя пользователя \\\hline
	exclude\_user\_filter &  Подключения от пользователей, совпавших с этим регулярным выражением не будут регистрироваться \\\hline
	include\_process\_filter &  Регулярное выражение, которому должно соответствовать название пользовательского процесса  \\\hline
	exclude\_process\_filter & Подключения от процессов, совпавших с этим регулярным выражением не будут регистрироваться  \\\hline
	include\_filter &  Регулярное выражение в синтаксисе SQL (\ttt{SIMILAR TO}), которому должен удовлетворять текст SQL-запроса \\\hline
	exclude\_filter & Регулярное выражение в синтаксисе SQL (\ttt{SIMILAR TO}), которому не должен удовлетворять текст SQL-запроса  \\\hline
	connection\_id &  Номер (идентификатор) подключения на сервере, которое
	будет отслеживаться \\\hline
	log\_filename &  Имя файла лога \\\hline
	\rowcolor[gray]{.9}\multicolumn{2}{|m{15.75cm}|}{\bfseries Int Values}\\\hline
	max\_log\_size &   Максимальный размер лог файлов в мегабайтах. Значение 0 означает неограниченный размер \\\hline
	time\_threshold & События, время выполнения которых меньше указанного (в мс), не будут регистрироваться в журнале  \\\hline
	max\_sql\_length & Максимальная длина одной записи SQL-запроса в лог-файле,	в байтах  \\\hline
	max\_blr\_length & Максимальная длина BLR-запроса, сохраняемого в лог, в байтах \\\hline
	max\_dyn\_length & Максимальная длина DYN-запроса, сохраняемого в лог, в байтах \\\hline
	max\_arg\_length & Максимальная длина одного параметра запроса/процедуры в
	лог-файле \\\hline
	max\_arg\_count & Максимальное количество параметров запроса/процедуры, которое заносится в лог-файл  \\\hline
\end{longtable}

\anonsubsection{Таблица событий}

Таблица событий в Трейс Менеджере заполняется либо при включении трассировки, либо при открытии сохраненного лог-файла в поле верхней части вкладки Трес Менеджера. При открытии лог-файла есть возможность выбрать кодировку:

\begin{figure}[H]
	\centering
	\includegraphics[width = 1\linewidth]{img/trace_log_file.png}
\end{figure}

В таблице событий по умолчанию отображаются все возможные столбцы, содержащие информацию о зарегистрированном событии. По желанию некоторые столбцы можно скрыть. Для этого нажмите на кнопку \ttt{Видимые столбцы} и внесите необходимые изменения.

Краткое описание значения каждого столбца таблицы описано ниже:

\begin{longtable}[r]{|>{\ttfamily}m{4.6cm}|m{10.9cm}|}
	\hline
	\centering\normalfont\bfseries Параметр &
	\centering\arraybslash\bfseries Описание\\\hline
	\endfirsthead
	\centering\normalfont\bfseries Параметр &
	\centering\arraybslash\bfseries Описание\\\hline
	\endhead
	\hline
	NUM &  Номер строки таблицы \\\hline
	TSTAMP & Время и дата произошедшего события \\\hline
	ID\_PROCESS &  Идентификатор процесса\\\hline
	ID\_THREAD & Идентификатор потока \\\hline
	EVENT\_TYPE &  Тип события \\\hline
	FAILED &  В случае неуспешной или несанкционированной попытки выполнения запроса фиксируется результат \ttt{FAILED} или \ttt{UNAUTHORIZED}\\\hline
	ID\_SESSION & Идентификатор сессии \\\hline
	NAME\_SESSION & Имя сессии \\\hline
	ID\_SERVICE & Идентификатор сервиса \\\hline
	USERNAME & Имя пользователя \\\hline
	PROTOCOL\_CONNECTION & Протокол соединения \\\hline
	CLIENT\_ADDRESS & IP адрес или имя компьютера \\\hline
	TYPE\_QUERY\_SERVICE & Тип запроса к сервису \\\hline
	OPTIONS\_START\_SERVICE &  Опции, переданные сервис-менеджеру от клиента при запуске \\\hline
	ROLE & Роль пользователя \\\hline
	DATABASE & Имя базы данных \\\hline
	CHARSET & Кодировка \\\hline
	ID\_CONNECTION &  Идентификатор соединения \\\hline
	CLIENT\_PROCESS & Клиентский процесс \\\hline
	ID\_CLIENT\_PROCESS & Идентификатор клиентского процесса \\\hline
	ID\_TRANSACTION &  Идентификатор транзакции \\\hline
	LEVEL\_ISOLATION & Уровень изоляции \\\hline
	MODE\_OF\_BLOCK &  Режим разрешения блокировок \\\hline
	MODE\_OF\_ACCESS & Режим доступа \\\hline
	TIME\_EXECUTION &  Время выполнения (в мс) \\\hline
	COUNT\_READS &  Количество страниц, считанных с диска \\\hline
	COUNT\_WRITES & Количество страниц, записанных на диск \\\hline
	COUNT\_FETCHES & Количество страниц, считанных из страничного кэша  \\\hline
	COUNT\_MARKS & Количество страниц, изменённых в страничном кэше \\\hline
	ID\_STATEMENT & Идентификатор запроса  \\\hline
	RECORDS\_FETCHED & Количество выбранных записей \\\hline
	STATEMENT\_TEXT &  Содержимое запроса\\\hline
	PARAMETERS\_TEXT &  Параметры выполнения запроса> \\\hline
	PLAN\_TEXT &  План запроса \\\hline
	TABLE\_COUNTERS & Табличные счетчики \\\hline
	DECLARE\_CONTEXT\_VARIABLES & Имя и значение контекстной переменной \\\hline
	EXECUTOR & Список пользователей, ролей и объектов базы данных, для которых выданы привилегии \\\hline
	GRANTOR & Пользователь от имени, которого предоставляются привилегии. \\\hline
	PRIVILEGE & Назначенная/отозванная привилегия \\\hline
	PRIVILEGE\_OBJECT & Объект, на который выдана (отозвана) привилегия \\\hline
	PRIVILEGE\_USERNAME & Имя пользователя, кому назначена привилегия \\\hline
	PRIVILEGE\_ATTACHMENT &  Идентификатор соединения \\\hline
	PRIVILEGE\_TRANSACTION & Идентификатор транзакции  \\\hline
	PROCEDURE\_NAME &  Имя процедуры \\\hline
	RETURN\_VALUE & Возвращаемые значения процедуры \\\hline
	TRIGGER\_INFO & \ttt{<имя\_триггера> [FOR <имя\_таблицы>] ({ON <событие БД>} | \{BEFORE | AFTER\} <событие таблицы или DDL-событие>)} \\\hline
	SENT\_DATA & Данные, переданные сервис-менеджеру \\\hline
	RECEIVED\_DATA & Данные, полученные сервис-менеджером \\\hline
	ERROR\_MESSAGE & Текст сообщения об ошибке \\\hline
	OLDEST\_INTERESTING &  Номер старейшей заинтересованной транзакции  \\\hline
	OLDEST\_ACTIVE &  Номер старейшей активной транзакции  \\\hline
	OLDEST\_SNAPSHOT &  Номер транзакции, которая была активной на момент	старта транзакции OAT \\\hline
	NEXT\_TRANSACTION &  Номер следующей транзакции \\\hline
	SORT\_MEMORY\_USAGE\_TOTAL & Суммарный размер кэша (в байтах), выделенного в процессе сортировки\\\hline
	SORT\_MEMORY\_USAGE\_CACHED & Размер RAM кэша (в байтах), выделенного в процессе сортировки\\\hline
	SORT\_MEMORY\_USAGE\_ON\_DISK & Размер дискового кэша\\\hline
\end{longtable}

\anonsubsubsection{Поиск в таблице}

В таблице поиск записей возможен только по одному из столбцов. Для этого служит выпадающий список \textit{Фильтр по столбцу} чуть выше самой таблицы. 

Введите строку для поиска в поле \textit{Текст} и тогда, в зависимости от настройки, строки с совпадениями будут либо подсвечены в таблице:
\begin{figure}[H]
	\centering
	\includegraphics[width = 0.95\linewidth]{img/trace_mag2.png}
\end{figure}
\noindent Либо таблица полностью урезается и остаются только искомые строки.

\anonsubsubsection{Просмотр содержимого}

Для удобного просмотра содержимого выделенной строки предусмотрено поле \textit{Тело события} справа от таблицы. Также при двойном клике на ячейку таблицы, открывается диалог просмотра данной ячейки. Для табличных счетчиков предусмотрена таблица,расположенная под таблицей событий, она заполняется при выделении строки в таблице событий

\anonsubsubsection{Экспорт таблиц}
При нажатии Правой кнопки мыши на любую из таблиц, появляется контекстное меню экспорта таблицы. Оно предусматривает возможность экспорта данных таблицы в следующие форматы:
\begin{enumerate}[leftmargin=20pt]
	\item Файл с разделителем
	\item Электронная таблица Excel
	\item Файл XML
	\item SQL скрипт заполнения таблицы базы данных
\end{enumerate}
	
\anonsubsection{Анализ трейса}	

Анализ трейса отображает информацию о событиях произошедших в указанный период времени. События добавляются либо при включении трассировки, либо при открытии лог-файла.

\begin{figure}[H]
	\centering
	\includegraphics[width = 1\linewidth]{img/trace_analysis.png}
\end{figure}

Параметры анализа трейса:
\begin{itemize}
	\item \ttt{TIME} - Если включен, то в результирующую таблицу будет добавлена информация о времени выполнения запроса;
	\item \ttt{READ} - Если включен, то в результирующую таблицу будет добавлена информация о количестве страниц, считанных с диска;
	\item \ttt{FETCH} -  Если включен, то в результирующую таблицу будет добавлена информация о количестве страниц, считанных из страничного кэша;
	\item \ttt{WRITE} - Если включен, то в результирующую таблицу будет добавлена информация о количестве страниц, записанных на диск;
	\item \ttt{MARK} - Если включен, то в результирующую таблицу будет добавлена информация о количестве страниц, изменённых в страничном кэше;
	\item \ttt{RSORT} - Если включен, то в результирующую таблицу будет добавлена информация об объёме ОЗУ, использованном для сортировки;
	\item \ttt{DSORT} - Если включен, то в результирующую таблицу будет добавлена информация о размере временных файлов, использованных в запросе;
	\item \ttt{Округлить значения} - Если значение больше 10000, то оно будет переводиться в большую единицу измерения, пока не станет меньше 10000;
	\item \ttt{Показать план} - Если включен, то будет показан план выполнения выбранного запроса;
	\item \ttt{Показать дополнительные параметры} - Если включен, то будут показаны дополнительные параметры анализа.
\end{itemize}

Дополнительные параметры:
\begin{figure}[H]
	\centering
	\includegraphics[width = 1\linewidth]{img/trace_analysis_additional_parameters.png}
\end{figure}
\begin{itemize}
	\item \ttt{Период} - Период времени, который нужно проанализировать; после изменения периода нужно нажать кнопку \ttt{Обновить};
	\item \ttt{Сравнить по N} - Считать запросы одинаковыми, если у них совпадают первые N символов;
	\item \ttt{Фильтр событий} - Позволяет выбрать типы событий, которые нужно анализировать.
\end{itemize}

Результирующая таблица:
\begin{figure}[H]
	\centering
	\includegraphics[width = 1\linewidth]{img/trace_analysis_result.png}
\end{figure}

При наведении на ячейку в всплывающей подсказке будет показана сумма значений столбца и среднее значение. Для просмотра полного текста запроса и его плана нужно выбрать ячейку. При двойном клике откроется окно с записью события в текстовом формате.

Описание значений столбцов таблицы:

\begin{longtable}[r]{|>{\ttfamily}m{4.6cm}|m{10.9cm}|}
	\hline
	\centering\normalfont\bfseries Столбец &
	\centering\arraybslash\bfseries Описание\\\hline
	\endfirsthead
	\centering\normalfont\bfseries Столбец &
	\centering\arraybslash\bfseries Описание\\\hline
	\endhead
	\hline
	QUERY & Текст запроса\\\hline
	COUNT & Количество выполнений запроса\\\hline
	PLAN\_COUNT & Количество планов запроса\\\hline
	TOTAL\_TIME & Общее время выполнения в миллисекундах\\\hline
	MIN\_TIME & Мимимальное время выполнения в миллисекундах\\\hline
	AVG\_TIME & Среднее время выполнения в миллисекундах\\\hline
	MAX\_TIME & Максимальное время выполнения в миллисекундах\\\hline
	STD\_DEV\_TIME & Среднеквадратическое отклонение времени выполнения\\\hline
	TOTAL\_READ & Общее количество страниц, считанных с диска\\\hline
	MIN\_READ & Минимальное количество страниц, считанных с диска\\\hline
	AVG\_READ & Среднее количество страниц, считанных с диска\\\hline
	MAX\_READ & Максимальное количество страниц, считанных с диска\\\hline
	STD\_DEV\_READ & Среднеквадратическое отклонение количества страниц, считанных с диска\\\hline
	TOTAL\_FETCH & Общее количество страниц, считанных из страничного кэша\\\hline
	MIN\_FETCH & Минимальное количество страниц, считанных из страничного кэша\\\hline
	AVG\_FETCH & Среднее количество страниц, считанных из страничного кэша\\\hline
	MAX\_FETCH & Максимальное количество страниц, считанных из страничного кэша\\\hline
	STD\_DEV\_FETCH & Среднеквадратическое отклонение количества страниц, считанных из страничного кэша\\\hline
	TOTAL\_WRITE & Общее количество страниц, записанных на диск\\\hline
	MIN\_WRITE & Мминимальное количество страниц, записанных на диск\\\hline
	AVG\_WRITE & Среднее количество страниц, записанных на диск\\\hline
	MAX\_WRITE & Максимальное количество страниц, записанных на диск\\\hline
	STD\_DEV\_WRITE & Среднеквадратическое отклонение количества страниц, записанных на диск\\\hline
	TOTAL\_MARK & Общее количество страниц, изменённых в страничном кэше\\\hline
	MIN\_MARK & Минимальное количество страниц, изменённых в страничном кэше\\\hline
	AVG\_MARK & Среднее количество страниц, изменённых в страничном кэше\\\hline
	MAX\_MARK & Максимальное количество страниц, изменённых в страничном кэше\\\hline
	STD\_DEV\_MARK & Среднеквадратическое отклонение количества страниц, изменённых в страничном кэше\\\hline
	TOTAL\_RSORT & Общий объём оперативной памяти, использованный для сортировки (в байтах)\\\hline
	MIN\_RSORT & Минимальный объём оперативной памяти, использованный для сортировки (в байтах)\\\hline
	AVG\_RSORT & Средний объём оперативной памяти, использованный для сортировки (в байтах)\\\hline
	MAX\_RSORT & Максимальный объём оперативной памяти, использованный для сортировки (в байтах)\\\hline
	STD\_DEV\_RSORT & Среднеквадратическое отклонение объёма оперативной памяти, использованного для сортировки\\\hline
	TOTAL\_DSORT & Общий размер временных файлов, использованный в запросе (в байтах)\\\hline
	MIN\_DSORT & Минимальный размер временных файлов, использованный в запросе (в байтах)\\\hline
	AVG\_DSORT & Средний размер временных файлов, использованный в запросе (в байтах)\\\hline
	MAX\_DSORT & Максимальный размер временных файлов, использованный в запросе (в байтах)\\\hline
	STD\_DEV\_DSORT & Среднеквадратическое отклонение размера временных файлов, использованного в запросе\\\hline
\end{longtable}
	
\anonsubsection{Менеджер сессий}
Менеджер сессий отображает список всех доступных в настоящее время сеансов трассировки. Вкладка менеджера видна только при запущенной текущей сессии аудита. 

\begin{figure}[H]
	\centering
	\includegraphics[width = 0.85\linewidth]{img/session_manager.png}
\end{figure}

В окне слева можно выбрать имя сессии и посмотреть информацию о ней (ID, запустившего пользователя, время запуска), а также завершить сеанс трассировки.

\newpage

\section{Статистика БД}\label{sec:gstat}

Инструмент отображает статистику базы данных, а также позволяет сравнить результаты анализа двух баз.

Открыть \ttt{Статистику БД} можно, выбрав соответствующий пункт в меню \ttt{Инструменты}. 

\begin{figure}[H]
	\centering
	\includegraphics[width = 0.95\linewidth]{img/gstat_con.png}
\end{figure}

Форматы вывода статистики:
\begin{itemize}
	\item По умолчанию -- анализ всей базы данных, вывод аналогичен выполнению \ttt{gstat} без опций;
	\item Таблицы -- статистика страниц данных, вывод аналогичен выполнению \ttt{gstat -data};
	\item Индексы -- анализ индексов, вывод аналогичен \ttt{gstat -index};
	\item Версии записей -- добавляет статистику о средних длинах записей, количестве версий и информацию о \ttt{BLOB};
	\item Системные объекты -- анализ системных таблиц и индексов;
	\item Страница заголовка -- статические данные о базе данных, вывод аналогичен \ttt{gstat -header};
	\item Только выбранные таблицы -- анализ выбранных таблиц, параметр доступен, если установлено соединение с выбранной базой данных.
\end{itemize}

Кнопка \ttt{Старт} запускает сбор статистики. После нажатия будет открыто окно с информацией о базе данных.

\paragraph{Вкладка Текст}
	
Вкладка \ttt{Текст} содержит полный вывод \ttt{gstat}:

\begin{figure}[H]
	\centering
	\includegraphics[width = 1\linewidth]{img/gstat_text.png}
	\caption{Статистика в текстовом формате}
\end{figure}

\paragraph{Вкладка Таблицы}

Вкладка \ttt{Таблицы} отображает статистику по таблицам:

\begin{figure}[H]
	\centering
	\includegraphics[width = 1\linewidth]{img/gstat_table.png}
	\caption{Статистика по таблицам}
\end{figure}

Описание значений столбцов вкладки \ttt{Таблицы}. Вычисляемые столбцы выделены синим.
\begin{longtable}[r]{|>{\ttfamily}m{4.6cm}|m{10.9cm}|}
	\hline
	\centering\normalfont\bfseries Столбец &
	\centering\arraybslash\bfseries Описание\\\hline
	\endfirsthead
	\centering\normalfont\bfseries Столбец &
	\centering\arraybslash\bfseries Описание\\\hline
	\endhead
	\hline
	name & Имя таблицы\\\hline
	primary pointer page & Номер первой страницы косвенных указателей на страницы, хранящие данные таблицы\\\hline
	index root page & Номер страницы, которая является первой страницей указателей на индексы таблицы \\\hline
	pointer pages & Общее количество страниц косвенных указателей на страницы, хранящие
	данные таблицы\\\hline
	data pages & Общее количество страниц, в которых хранятся данные таблицы. Этот
	счетчик включает страницы, хранящие неподтвержденные версии записей и мусор, потому что \ttt{gstat} не может их отличить друг от друга\\\hline
	data page slots & Количество указателей на страницы базы данных, содержащихся на
	страницах указателей. Должно равняться числу страниц данных\\\hline
	primary pages & Количество страниц, равное (\ttt{Data pages - Secondary pages})\\\hline
	secondary pages & Количество страниц, на которых не хранятся первичные версии записей\\\hline
	swept pages & Количество страниц, которые имеют только первичные версии записей, и все они были созданы подтвержденными транзакциями. Такие страницы данных должны быть пропущены процедурой \ttt{sweep}\\\hline
	empty pages & Количество страниц, на которых нет записей\\\hline
	full pages & Количество полностью заполненных страниц\\\hline
	big record pages & Количество страниц, которые полностью заняты только одной
	записью\\\hline
	blob pages & Количество страниц с блобами\\\hline
	average record length & Средний размер сжатой записи в байтах\\\hline
	total records & Общее количество строк в таблице\\\hline
	average version length & Среднее значение длины старых версий в байтах\\\hline
	total versions & Общее количество старых версий в таблице\\\hline
	max versions & Максимальная цепочка старых версий для записи\\\hline
	average fill & Обобщающая гистограмма распределения использования памяти для всех страниц, выделенных в таблице\\\hline
	total formats & Общее количество форматов в таблице \ttt{RDB\$FORMATS}\\\hline
	used formats & Количество используемых форматов\\\hline
	average fragment length & Средний размер фрагмента в байтах\\\hline
	total fragments & Количество всех фрагментов во всех записях\\\hline
	max fragments & Максимальное количество фрагментов в одной записи\\\hline
	average unpacked length & Средний размер записи в байтах (без сжатия)\\\hline
	compression ratio & Отношение средней длины ключа без учета сжатия (\ttt{Average prefix length + Average data length}) к средней длине ключа с учетом сжатия (\ttt{Average key length}).\\\hline
	blobs & Количество всех блобов (0, 1 и 2 уровня) \\\hline
	total length & Размер, занимаемый блобами, в байтах\\\hline
	level <n> & Количество блобов каждого уровня\\\hline
	\textcolor{blue}{table size(without blobs)} & Размер таблицы в байтах (\ttt{data pages * page size}). Вычисляемый столбец\\\hline
	\textcolor{blue}{size with blobs} & Совокупный размер данных таблицы (\ttt{table size(without blobs) + blobs total length}). Вычисляемый столбец\\\hline
	\textcolor{blue}{size with blob pages} & (\ttt{table size(without blobs) + blob pages * page size}). Вычисляемый столбец\\\hline
	\textcolor{blue}{size with indices} & Совокупный размер данных таблицы и совокупный размер всех индексов по таблице. Вычисляется по формуле:
	$$X = S + \sum_{i=0}^{n} E_i $$
	 
	 где \ttt{S} -- \ttt{size with blob pages};
	 
	 \ttt{n} -- количество индексов в таблице;
	 
	 \ttt{E} -- предполагаемый размер индекса в байтах (\ttt{estimated full size}).
	\\\hline
	range <процент заполненности> & Количество страниц, которые заполнены на указанный процент\\\hline
\end{longtable}

\paragraph{Вкладка Индексы}

Вкладка \ttt{Индексы} отображает статистику по индексам:

\begin{figure}[H]
	\centering
	\includegraphics[width = 1\linewidth]{img/gstat_index.png}
	\caption{Статистика по индексам}
\end{figure}

Описание значений столбцов вкладки \ttt{Индексы}. Вычисляемые столбцы выделены синим.
\begin{longtable}[r]{|>{\ttfamily}m{4.6cm}|m{10.9cm}|}
	\hline
	\centering\normalfont\bfseries Столбец &
	\centering\arraybslash\bfseries Описание\\\hline
	\endfirsthead
	\centering\normalfont\bfseries Столбец &
	\centering\arraybslash\bfseries Описание\\\hline
	\endhead
	\hline
	name & Имя индекса\\\hline
	table name & Имя таблицы\\\hline
	\textcolor{blue}{real selectivity} & Вычисляется по формуле \ttt{1 / (nodes - totalDup)}\\\hline
	average data length & Средняя длина каждого ключа в байтах. Она, скорее всего, меньше, чем фактическая сумма размеров столбцов, поскольку Ред База Данных использует индексное сжатие для уменьшения объема данных, хранящихся на странице листа индекса\\\hline
	total dup & Общее количество строк дубликатов индекса \\\hline
	max dup & Количество дублирующих узлов в "цепочке", имеющих наибольшее
	количество дубликатов. Всегда будет нулем для уникальных индексов.
	Если число велико по сравнению с числом в \ttt{total dup}, то это признак
	плохой селективности \\\hline
	root page & Номер корневой страницы индекса \\\hline
	depth & Количество уровней в странице индексного дерева. Если глубина дерева индексной страницы превышает 3, то доступ к записям через индекс не будет максимально эффективным. Для уменьшения глубины дерева индексной страницы увеличьте размер страницы. Если увеличение размера страницы не уменьшает глубины, снова увеличьте размер страницы \\\hline
	leaf buckets & Количество страниц самого низкого уровня (листовых) в дереве индекса. Это страницы, которые содержат указатели на записи. Страницы высокого уровня содержат косвенные связи. \\\hline
	\textcolor{blue}{leaf full size} & Размер страниц самого низкого уровня (листовых) в дереве индекса (leaf buckets * page size). Вычисляемый столбец.\\\hline
	\textcolor{blue}{estimated full size} & Приблизительный размер индекса в байтах. Вычисляется по формуле:
	
	 $$\ttt{N} = \ttt{L} * (1+\frac{\ttt{A}}{\ttt{P}})^{\ttt{D - 1}}*\ttt{P},$$
	 
	 где \ttt{N} -- приблизительный размер индекса в байтах;
	 
	 \ttt{L} -- количество страниц самого низкого уровня (листовых) в дереве индекса (\ttt{leaf buckets});
	 
	 \ttt{A} -- средний размер узлов в байтах (\ttt{average node length});
	 
	 \ttt{P} -- размер страницы базы данных (\ttt{page size});
	 
	 \ttt{D} -- количество уровней в странице индексного дерева (\ttt{depth}).
	 \\\hline
	nodes & Общее количество записей, индексированных в дереве. Должно быть
	равно количеству индексированных строк в дереве, хотя отчет \ttt{gstat} может включать узлы, которые были удалены, но не вычищены в процессе сборки мусора. Может также включать множество элементов для записей, у которых был изменен индексный ключ\\\hline
	average node length & Средний размер узлов в байтах\\\hline
	average key length & Средний размер ключа в байтах с учетом сжатия. К длине каждого
	ключа прибавляется от 1 до 5 байт в зависимости от размера ключа и префикса. Затем высчитывается средний размер ключа\\\hline
	compression ratio & Средняя длина ключа и префикса в байтах \\\hline
	average prefix length & Средний размер, занимаемый префиксами узлов, байтах\\\hline
	average data length & Средняя длина каждого ключа в байтах. Она, скорее всего, меньше, чем фактическая сумма размеров столбцов, поскольку Ред База Данных использует индексное сжатие для уменьшения объема данных, хранящихся на странице листа индекса\\\hline
	clustering factor & Это мера того, насколько много операций ввода-вывода будет осуществлять база данных, если бы ей пришлось читать каждую строку таблицы посредством индекса, в порядке индекса. То есть она показывает, насколько упорядочены строки в таблице по значениям индекса. Если значение близко к общему количеству страниц, значит таблица очень хорошо упорядочена. В этом случае записи индекса на одной странице листа индекса обычно указывают на строки, находящиеся в одних и тех же страницах данных. Если значение близко к общему количеству
	строк, значит, таблица весьма неупорядочена. В этом случае маловероятно, что записи индекса на одной странице листа индекса указывают на те же страницы данных\\\hline
	ratio & Отношение \ttt{Clustering factor} к общему количеству узлов в индексе\\\hline
	full size & Размер индекса в байтах. Рассчитывается по формуле: \ttt{(leaf\_buckets * (1 + (avg\_node\_length / page\_size)) \^(depth - 1))*page\_size} \\\hline
\end{longtable}

\paragraph{Вкладка Табличные пространства}

Вкладка \ttt{Табличные пространства} отображает статистику по табличным пространствам:

\begin{figure}[H]
	\centering
	\includegraphics[width = 1\linewidth]{img/gstat_tablespace.png}
	\caption{Статистика по табличным пространствам}
\end{figure}


Описание значений столбцов вкладки \ttt{Табличные пространства}:
\begin{longtable}[r]{|>{\ttfamily}m{4.6cm}|m{10.9cm}|}
	\hline
	\centering\normalfont\bfseries Столбец &
	\centering\arraybslash\bfseries Описание\\\hline
	\endfirsthead
	\centering\normalfont\bfseries Столбец &
	\centering\arraybslash\bfseries Описание\\\hline
	\endhead
	\hline
	name & Имя табличного пространства\\\hline
	full path & Путь к файлу табличного пространства\\\hline
	table count & Количество таблиц, находящихся в табличном пространстве\\\hline
	index count & Количество индексов, находящихся в табличном пространстве\\\hline
\end{longtable}

\anonsubsection{Сравнение результатов анализа двух баз}

Для сравнения статистик двух баз данных выполните сбор статистики для второй базы и нажмите кнопку \ttt{Сравнить}:

\begin{figure}[H]
	\centering
	\includegraphics[width = 0.2\linewidth]{img/gstat_compare_button.png}
	\caption{Кнопка сравнить}
\end{figure}

Во вкладке \ttt{Текст} будет отображена разница между текстовыми выводами статистики:

\begin{figure}[H]
	\centering
	\includegraphics[width = 1\linewidth]{img/gstat_compare.png}
	\caption{Сравнение}
\end{figure}

Во вкладках \ttt{Таблицы} и \ttt{Индексы} отображается сравнение результатов анализа таблиц и индексов соответственно. В значениях столбцов отображается разница между результатом первой и второй базы.  

\newpage

\section{Системные журналы состояния}

\anonsubsection{Системный журнал состояния вывода}\label{sec:logout}
В журнал записывается вся информация, которая выводится в стандартный поток вывода, ошибки и предупреждения. В \textit{Настройках} можно задать уровень вывода для ограничения содержащейся в логе информации.

Системный журнал вывода хранится в \ttt{\$HOME/.redexpert/logs} и его можно открыть любым текстовым редактором. Но его также можно посмотреть и с помощью Red Expert, выбрав пункт меню  \textit{Инструменты}$\rightarrow$\textit{Системные журналы состояния}$\rightarrow$\textit{Системный журнал состояния вывода}. На центральной панели будет открыта новая вкладка с содержимым журнала.

Чтобы очистить содержимое журнала, нажмите на кнопку 
\begin{tikzpicture}
\pgftext{\includegraphics[width=12pt]{img/Delete16.png}} at (0pt,0pt)
\end{tikzpicture}
панели инструментов или выберите пункт меню \textit{Инструменты}$\rightarrow$\textit{Системные журналы состояния}$\rightarrow$\textit{Очистить журналы состояний}$\rightarrow$\textit{Системный журнал состояния вывода}.

\newpage


\section{Настройки}
Для настройки приложения выберите пункт меню \textit{Инструменты}$\rightarrow$\textit{Настройки}.
\begin{redremark}
	Большинство изменений в настройках вступают в силу после перезапуска Red Expert.
\end{redremark}
\anonsubsection{Общие}
\begin{longtable}[r]{|>{\ttfamily}m{6cm}|m{9cm}|}
	\hline
	\centering\normalfont\bfseries Параметр &
	\centering\arraybslash\bfseries Комментарий\\\hline
	\endfirsthead
	\hline
	\centering\normalfont\bfseries Параметр &
	\centering\arraybslash\bfseries Комментарий\\\hline
	\endhead
	\rowcolor[gray]{.9}\multicolumn{2}{|m{15.25cm}|}{\bfseries Общие}\\\hline 
	Отображение заставки при запуске
	&  \\\hline
	Развернуть окно при запуске 
	& Окно приложения раскроется на весь экран при запуске\\\hline
	Количество отображаемых последних файлов
	& Количество сохраненных файлов журнала логов. При достижении журнала (system.log) размера более 1МБ, запись продолжается в новый файл (system.log.1) и т.д. Когда количество файлов будет равно числу, указанному в этом параметре, файлы будут перезаписаны. Если указано 0, то количество файлов неограниченно. \\\hline
	Разделитель строк
	&  \\\hline
	Спрашивать о сохранении открытых документов
	&  \\\hline
	Проверять обновления при запуске
	& \\\hline
	Загружать нестабильные версии
	&  \\\hline 
	Проверять мажорные версии
	&  \\\hline
	Кодировка файлов по умолчанию
	&   \\\hline
	Путь к исполняемому файлу Java
	& Путь к исполняемому файлу Java. Можно указывать как абсолютный путь, так и относительный. \\\hline
	\rowcolor[gray]{.9}\multicolumn{2}{|m{15.25cm}|}{\bfseries Настройка прокси}\\\hline 
	Использовать прокси 
	&  Использовать прокси-сервер для подключения к интернету \\\hline
	Прокси сервер
	&  Имя или IP прокси-сервера \\\hline
	Прокси порт
	& Номер порта прокси-сервера \\\hline 
	Прокси пользователь
	&  Имя пользователя для авторизации на прокси-сервере (если требуется)\\\hline
	Прокси пароль
	& Пароль пользователя для авторизации на прокси-сервере (если требуется)\\\hline  
\end{longtable} 

\anonsubsection{Портативизация настроек}

Чтобы хранить \ttt{Red Expert} и его настройки в одной папке выполните следующие действия:
\begin{enumerate}
	\item Закройте программу;
	\item Откройте файл \ttt{...\textbackslash каталог c \ttt{Red Expert}\textbackslash config\textbackslash redexpert\_config.ini};
	\item Для параметра \ttt{eq.user.home.dir} укажите значение \ttt{../.redexpert} и сохраните изменения;
	\item В адресной строке проводника введите путь \ttt{\%homepath\%};
	\item Переместите папку \ttt{.redexpert} в корень каталога, в который установлен \ttt{Red Expert};
	\item Запустите \ttt{Red Expert}.
	
\end{enumerate}
	
\anonsubsection{Местоположение и язык}

Выберите подходящий язык для кнопок, диалогов и сообщений (в том числе об ошибках) из списка.

Выберите часовой пояс для вашего региона.


\anonsubsection{Настройки отображения}

\begin{longtable}[r]{|>{\ttfamily}m{6cm}|m{9cm}|}
	\hline
	\centering\normalfont\bfseries Параметр &
	\centering\arraybslash\bfseries Комментарий\\\hline
	\endfirsthead
	\hline
	\centering\normalfont\bfseries Параметр &
	\centering\arraybslash\bfseries Комментарий\\\hline
	\endhead
	\rowcolor[gray]{.9}\multicolumn{2}{|m{15.25cm}|}{\bfseries Общие}\\\hline 
	Строка состояния
	& Показать/скрыть панель инструментов \\\hline
	Системная консоль
	& Показать/скрыть консоль вывода \\\hline
	Подключения
	&  \\\hline
	Драйверы
	&  \\\hline
	Ключевые слова SQL
	&  \\\hline
	Коды состояния SQL 
	&  \\\hline
	Палитра свойств системы
	& Показать/скрыть панель с системными свойствами \\\hline
	\rowcolor[gray]{.9}\multicolumn{2}{|m{15.25cm}|}{\bfseries Внешний вид}\\\hline
	Тема &  \\\hline 
	Использовать сглаженные шрифты &   \\\hline
	Фон рабочего стола &  \\\hline
	Стилизовать диалоги & Декорирование рамок диалоговых окон \\\hline
	Стилизовать окна & Декорирование рамки окна приложения\\\hline
\end{longtable} 

\anonsubsection{Горячие клавиши}

Здесь перечислены действующие горячие клавиши для различных команд при работе с Red Expert. Их можно изменить, удалить или добавить новые. Для этого просто щелкните мышью по нужной ячейке -- откроется диалоговое окно:
\begin{figure}[H]
	\centering
	\includegraphics[width = 0.45\linewidth]{img/select_shortcut.png}
\end{figure}
Введите поочередно комбинацию клавиш на клавиатуре и нажмите кнопку \textit{OK}.

\anonsubsection{Плагины стилей}

В этом разделе настроек можно установить плагин с новой темой приложения. И в дальнейшем активировать ее в разделе настроек \textit{Настройки отображения}.

\anonsubsection{Панель инструментов}\label{sec:pref_toolbar}

Вся панель инструментов разделена на несколько логически связанных частей. Они подробно описаны в разделе \hyperref[sec:toolbar]{Панель инструментов}. В настройках можно указать какие кнопки будут присутствовать на панели, а какие скрыты. 

\anonsubsection{Редактор}
Настройки для Редактора запросов.
\begin{longtable}[r]{|>{\ttfamily}m{6cm}|m{9cm}|}
	\hline
	\centering\normalfont\bfseries Параметр &
	\centering\arraybslash\bfseries Комментарий\\\hline
	\endfirsthead
	\hline
	\centering\normalfont\bfseries Параметр &
	\centering\arraybslash\bfseries Комментарий\\\hline
	\endhead
	\hline
	\rowcolor[gray]{.9}\multicolumn{2}{|m{15.25cm}|}{\bfseries Общие}\\\hline 
	Преобразовывать табуляцию в пробелы
	&  \\\hline
	Количество пробелов в табуляции
	&  \\\hline
	Автозавершение ключевых слов 
	& Включить/отключить подсказку ключевых слов по нажатию \ttt{Ctrl+Space}\\\hline
	Автозавершение объектов БД 
	& Включить/отключить подсказку объектов базы данных по нажатию \ttt{Ctrl+Space}\\\hline
	Количество возможных отмен
	& Количество возвратов к предыдущим состояниям. Команда возврата активируется по нажатию \ttt{Ctrl+Z} (Windows) \\\hline 
	Количество сохраняемых запросов в истории
	& \\\hline 
	Авто-коммит для редактора по умолчанию
	&  \\\hline 
	Хранить метаданные результирующего набора 
	& \\\hline
	Открывать результирующие наборы данных в одной вкладке
	& Если опция не включена, то каждый результирующий набор помещается на отдельную вкладку Редактора запросов. Иначе каждая последующая выборка перезаписывает предыдущую на одной и той же вкладке \\\hline  
	Удалять комментарии перед выполнением 
	& Удалять SQL комментарии в запросах во время исполнения и отправлять на сервер запрос без них \\\hline
	Максимум возвращаемых строк по умолчанию
	& Максимальное число строк, показанное в таблице результирующего набора \\\hline 
	Выводить весь запрос на панель вывода 
	& \\\hline 
	Вывод лога в файл
	&  Сохранять записи панели вывода Редактора запросов в файл\\\hline 
	Путь к директории с логами
	&  Определяет абсолютный или относительный путь к директории, в которой будет создана папка \ttt{logs}, содержащая файлы логов. Для использования относительного пути в начале строки необходимо дописать \ttt{\%re\%}.\\\hline
	Количество сохраняемых файлов с логами
	& Количество сохраненных файлов с логами. При достижении файла с логами размера более 1МБ, он переименовывается в файл logfile.log.1 и т.д. И запись снова продолжается в файл logfile.log. Когда количество файлов будет равно числу, указанному в этом параметре, файлы будут перезаписаны. Если указано 0, то количество файлов неограниченно.\\\hline
	Открывать новый редактор для нового открытого подключения
	&  \\\hline
	\rowcolor[gray]{.9}\multicolumn{2}{|m{15.25cm}|}{\bfseries Настройки отображения}\\\hline     
	Строка состояния
	&  \\\hline
	Панель инструментов
	&  \\\hline
	Номера строк
	&  \\\hline
	Панель результатов
	& \\\hline
	Подсветка текущей строки
	& Подсветка строки, на которой стоит курсор \\\hline
	Поле справа
	& Отступ справа в поле ввода Редактора\\\hline
	Количество столбцов до поля справа
	& Ширина отступа\\\hline
	Цвет поля справа
	& Цвет линии отступа \\\hline      
\end{longtable}

\anonsubsubsection{Шрифты}

Измените шрифт и его размер для поля ввода запросов.

\anonsubsubsection{Цвета}

\begin{longtable}[r]{|>{\ttfamily}m{6cm}|m{9cm}|}
	\hline
	\centering\normalfont\bfseries Параметр &
	\centering\arraybslash\bfseries Комментарий\\\hline
	\endfirsthead
	\hline
	\centering\normalfont\bfseries Параметр &
	\centering\arraybslash\bfseries Комментарий\\\hline
	\endhead
	\hline
	\rowcolor[gray]{.9}\multicolumn{2}{|m{15.25cm}|}{\bfseries Цвета редактора запросов}\\\hline 
	Фон редактора &  Фон поля ввода запросов \\\hline
	Цвет шрифта редактора &  \\\hline
	Подсветка текущей строки &  \\\hline
	Цвет фона столбца с номерами строк &   \\\hline
	Цвет курсора &   \\\hline
	Цвет выделения & Цвет фона выделенного текста \\\hline
	Фон панели результатов &  Фон панели вывода и Result Set \\\hline
	Text selection alternative background &  \\\hline
	Цвет номеров строк &  \\\hline
	Цвет выделенного текста &  \\\hline
	Альтернативный фон редактора &  \\\hline
	\rowcolor[gray]{.9}\multicolumn{2}{|m{15.25cm}|}{\bfseries Синтаксические стили}\\\hline 
	Обычный текст &  \\\hline
	Ключевые слова &  \\\hline
	Строка в кавычках &  \\\hline
	Однострочный комментарий & \\\hline
	Многострочный комментарий &   \\\hline
	Число & \\\hline
	Оператор & Цвет операторов (+, -, =, <, > и т.д.) \\\hline
	Скобки &  \\\hline
	Литерал &  \\\hline
	Подсветка парных скобок &  \\\hline
	Ошибочная скобка & Цвет подсвечивания скобки с отсутствием пары \\\hline
	Объекты БД &  Цвет имен объектов базы данных \\\hline
\end{longtable}


\anonsubsection{Дерево подключений}

Этот раздел служит для редактирования содержимого панели \textit{Подключения}.
Параметр \textit{Высота узла} задает высоту одного узла дерева подключений (или, другими словами, определяет величину расстояния между узлами). В подразделе \textit{Шрифты} можно установить стиль и размер шрифта текста в окне \textit{Подключения}.

\anonsubsection{Таблица набора данных}

Настройки для таблицы результатов выборки.
\hypertarget{tab:datapattern}{}
\begin{longtable}[r]{|>{\ttfamily}m{5cm}|m{10.5cm}|}
	\hline
	\centering\normalfont\bfseries Параметр &
	\centering\arraybslash\bfseries Комментарий\\\hline
	\endfirsthead
	\hline
	\centering\normalfont\bfseries Параметр &
	\centering\arraybslash\bfseries Комментарий\\\hline
	\endhead
	\hline
	Изменяемость ширины столбцов &  Возможность менять ширину столбцов таблицы результатов \\\hline
	Перестановка столбцов & Возможность перемещать столбцы таблицы результатов\\\hline
	Показывать номера строк & Отображать ли номера строк в таблице результатов\\\hline
	Ширина столбца & Ширина столбца таблицы результатов в пикселях\\\hline
	Высота столбца &  Высота столбца таблицы результатов в пикселях\\\hline
	Сохранять ширину столбцов между запросами & Запоминать настроенную самостоятельно ширину столбцов в следующих результатах запросов \\\hline
	Шаблон даты & Шаблон даты для типа данных \ttt{DATE}. В таблице ниже рассмотрены подробно параметры шаблона\\\hline
	Шаблон времени & Шаблон времени для типа данных \ttt{TIME}. В таблице ниже рассмотрены подробно параметры шаблона\\\hline
	Шаблон даты-времени & Шаблон даты-времени для типа данных \ttt{TIMESTAMP}. В таблице ниже рассмотрены подробно параметры шаблона\\\hline
	Текст в ячейке со значением null & Текст, который будет отображаться в ячейке со значением null\\\hline
	Открывать визуализатор данных по двойному клику на ячейку & Двойной клик по ячейке таблицы открывает диалоговое окно \textit{Record Data Item Viewer} с подробной информацией о типе данных, размере и содержимом ячейки \\\hline
	Транспонировать результат с одной строкой & Транспонировать таблицу, когда строки становятся столбцами и наоборот, если результат из одной строки \\\hline
	Выравнивание числовых значений по правому краю & Выравнивать ли числовые значения по правому краю\\\hline
	Использовать форму для добавления/удаления строк & При добавлении или удалении записи из таблицы на вкладке \textit{Данные} (при редактировании таблицы базы данных) будет появляться диалоговое окно с подтверждением операции. \\\hline
	Использовать другой цвет для значений с null при добавлении/удалении записей & Подсвечивать значения равные \ttt{NULL} при добавлении/удалении записей на вкладке \textit{Данные} (при редактировании таблицы базы данных) \\\hline
	Количество подгружаемых записей за раз & Число первоначально загружаемых строк, остальные строки подгружаются при скролле \\\hline
\end{longtable}

\begin{longtable}[r]{|>{\ttfamily\centering}m{1.8cm}|m{11.7cm}|m{1.8cm}|}
	\hline
	\centering\normalfont\bfseries Символ &
	\centering\normalfont\bfseries Описание &
	\centering\arraybslash\bfseries Пример\\\hline
	\endfirsthead
	\hline
	\centering\normalfont\bfseries Символ &
	\centering\normalfont\bfseries Описание &
	\centering\arraybslash\bfseries Пример \\\hline
	\endhead
	\hline
	G &  эра (в английской локализации - AD и BC) & н.э. \\\hline
	y & год (4-х значное число) & 2012 \\\hline
	yy & год (последние 2 цифры) & 12 \\\hline
	yyyy & год (4-х значное число) & 2012 \\\hline
	M & номер месяца без лидирующих нулей & 2 \\\hline
	MM & номер месяца (с лидирующими нулями если номер месяца < 10) & 02 \\\hline
	MMM & четырех буквенное сокращение месяца в русской локализации и трех буквенное - в английской (Feb) & фев \\\hline
	MMMM & полное название месяца (в английской локализации - February) & Февраль \\\hline
	w & неделя в году без лидирующих нулей & 7 \\\hline
	ww & неделя в году с лидирующими нулями & 07 \\\hline
	W & неделя в месяце без лидирующих нулей & 2 \\\hline
	WW & неделя в месяце с лидирующим нулем (если это необходимо) & 02 \\\hline
	D & день в году & 38 \\\hline
	d & день месяца без лидирующих нулей & 7 \\\hline
	dd & день месяца с лидирующими нулями & 07 \\\hline
	F & день недели в месяце без лидирующих нулей & 1 \\\hline
	FF & день недели в месяце с лидирующими нулями & 01 \\\hline
	E & день недели (сокращение) & Вт \\\hline
	EEEE & день недели (полностью) & вторник \\\hline
	a & AM/PM указатель & AM \\\hline
	H & часы в 24-часовом формате без лидирующих нулей & 6 \\\hline
	HH & часы в 24-часовом формате с лидирующим нулем & 06 \\\hline
	k & количество часов в 24-часовом формате & 18 \\\hline
	K & количество часов в 12-часовом формате & 6 \\\hline
	h & время в 12-часовом формате без лидирующих нулей & 6 \\\hline
	hh & время в 12-часовом формате с лидирующим нулем & 06 \\\hline
	m & минуты без лидирующих нулей & 32 \\\hline
	mm & минуты с лидирующим нулем & 32 \\\hline
	s & секунды без лидирующих нулей & 11 \\\hline
	ss & секунды с лидирующим нулем & 11 \\\hline
	S & миллисекунды (только первая цифра) & 1 \\\hline
	SS & миллисекунды (только первые две цифры)	 & 11 \\\hline
	z & часовой пояс & EET \\\hline
	Z & часовой пояс в формате RFC 822 & +0200 \\\hline
	' & символ экранирования для текста & 'Date=' \\\hline
	'' & кавычка & 'o''clock' \\\hline	
\end{longtable}

\anonsubsubsection{Цвета}

\begin{longtable}[r]{|>{\ttfamily}m{6cm}|m{9.5cm}|}
	\hline
	\centering\normalfont\bfseries Параметр &
	\centering\arraybslash\bfseries Комментарий\\\hline
	\endfirsthead
	\hline
	\centering\normalfont\bfseries Параметр &
	\centering\arraybslash\bfseries Комментарий\\\hline
	\endhead
	\hline
	Фон ячейки по умолчанию &  \\\hline
	Фон ячейки с null & \\\hline
	Фон ячейки с null при добавлении & Цвет ячеек со значением \ttt{NULL} при добавлении записи \\\hline
	Фон ячейки с null при удалении & Цвет ячеек со значением \ttt{NULL} при удалении записи \\\hline
	Фон ячейки с измененным значением &  \\\hline
	Фон ячейки с удаляемым значением &  \\\hline
	Фон ячейки с новым значением &   \\\hline
	Фон ячейки со строковым значением &  \\\hline
	Фон ячейки с числовым значением &  \\\hline
	Фон ячейки со значением даты/времени & \\\hline
	Фон ячейки с логическим значением & \\\hline
	Фон ячейки с бинарным значением &   \\\hline
	Фон ячейки с прочим значением & Цвет ячеек с данными других типов \\\hline
	Фон ячейки на которой фокус & Цвет фона выделенной строки \\\hline
	Альтернативный фон ячейки & Цвет строк с чередованием \\\hline
\end{longtable}

\anonsubsection{Подключение}

\begin{longtable}[r]{|>{\ttfamily}m{5cm}|m{10.5cm}|}
	\hline
	\centering\normalfont\bfseries Параметр &
	\centering\arraybslash\bfseries Комментарий\\\hline
	\endfirsthead
	\hline
	\centering\normalfont\bfseries Параметр &
	\centering\arraybslash\bfseries Комментарий\\\hline
	\endhead
	\hline
	Автоматически подключаться при запуске & Подключаться к выбранной в пункте \textit{Подключение для автозапуска} базе данных при запуске Red Expert\\\hline
	Подключение для автозапуска &  Выберите подключение, которое будет установлено при запуске Red Expert\\\hline
\end{longtable}

\anonsubsection{Браузер базы данных}


\begin{longtable}[r]{|>{\ttfamily}m{5cm}|m{10.5cm}|}
	\hline
	\centering\normalfont\bfseries Параметр &
	\centering\arraybslash\bfseries Комментарий\\\hline
	\endfirsthead
	\centering\normalfont\bfseries Параметр &
	\centering\arraybslash\bfseries Комментарий\\\hline
	\endhead
	Показать только каталог или схему по умолчанию & \\\hline
	Подключение по двойному клику & Выполнять соединение с базой данных по двойному клику на подключении \\\hline
\end{longtable}

\anonsubsubsection{Панель данных таблицы}

\begin{longtable}[r]{|>{\ttfamily}m{5cm}|m{10.5cm}|}
	\hline
	\centering\normalfont\bfseries Параметр &
	\centering\arraybslash\bfseries Комментарий\\\hline
	\endfirsthead
	\centering\normalfont\bfseries Параметр &
	\centering\arraybslash\bfseries Комментарий\\\hline
	\endhead
	Максимальное количество возвращаемых записей  &  \\\hline
\end{longtable}

\anonsubsection{Консоль}

\begin{longtable}[r]{|>{\ttfamily}m{5cm}|m{10.5cm}|}
	\hline
	\centering\normalfont\bfseries Параметр &
	\centering\arraybslash\bfseries Комментарий\\\hline
	\endfirsthead
	\centering\normalfont\bfseries Параметр &
	\centering\arraybslash\bfseries Комментарий\\\hline
	\endhead
	Отображать системную консоль & \\\hline
	Уровень лога событий
	& Уровень вывода в \hyperref[sec:logout]{журнал состояния вывода} и в консоль вывода:
	\begin{enumerate}[leftmargin=10pt]
		\item \ttt{FATAL} -- фиксирует очень серьезные ошибки, которые, по-видимому, приведут к прекращению использования приложения;
		\item \ttt{ERROR} -- фиксирует серьезные ошибки, нуждающаяся в проверке, но не прекращающая работу приложения;
		\item \ttt{WARN} -- фиксирует потенциально опасные ситуации;
		\item \ttt{INFO} -- информирует о событиях работы приложения на высоком уровне;
		\item \ttt{DEBUG} -- информирует о событиях на более низком уровне для отладки приложения (для разработчиков);
		\item \ttt{TRACE} -- подробное ведение журнала для очень низкоуровневой отладки;
		\item \ttt{ALL} -- все вышеперечисленные уровни.
	\end{enumerate} Уровни упорядочены и каждый следующий включает в себя предыдущие уровни.\\\hline 
	Системный вывод в консоль
	& Выводить поток System.out в консоль\\\hline
	Вывод системных ошибок в консоль
	& Выводить поток System.err в консоль \\\hline
\end{longtable}

\anonsubsubsection{Шрифты}

Измените шрифт и его размер для текста консоли.

\newpage

\section{Обратная связь}

Вы можете связаться с нами через форму обратной связи в приложении (\textit{Справка} $\rightarrow$ \textit{Обратная связь}) или отправив письмо по адресу \href{mailto:rdb.support@red-soft.ru}{rdb.support@red-soft.ru}. В письме Вы можете оставить отзыв о работе программы, сообщить нам об ошибке или предложить новый функционал. 

Мы будем рады Вам помочь.


\appendix

\titleformat{\section}
{\huge\color{red}}
{\thesection}{1em}{}

\titleformat{\subsection}
{\LARGE\color{red}}
{\thesubsection}{1em}{}

\renewcommand{\thetable}{\Asbuk{application}.\arabic{table}} % нумерация таблиц с номером section

\renewcommand{\thefigure}{\Asbuk{application}.\arabic{figure}}

\renewcommand{\thesubsection}{\Asbuk{application}.\arabic{subsection}} % нумерация подразделов типа F.1


\renewcommand{\theredlisting}{\Asbuk{application}.\arabic{redlisting}}

\renewcommand{\thesubsection}{}



\appcount\section{Свойства базы данных} \apptest
\label{app:1}

В таблице отражена информация о возможностях конкретной СУБД в сочетании с конкретным JDBC драйвером. 
Представленные свойства являются методами интерфейса \href{https://docs.oracle.com/javase/7/docs/api/java/sql/DatabaseMetaData.html}{ DatabaseMetaData}.

\begin{longtable}[c]{|>{\ttfamily}m{5.1cm}|>{\ttfamily\centering}m{1.3cm}|m{9.2cm}|}
	\hline
	\centering\normalfont\bfseries Свойство &
	\centering\normalfont\bfseries Тип &
	\centering\arraybslash\bfseries Описание\\\hline
	\endfirsthead
	\hline
	\centering\normalfont\bfseries Свойство &
	\centering\normalfont\bfseries Тип &
	\centering\arraybslash\bfseries Описание\\\hline
	\endhead
	allProceduresAreCallable & boolean & Может ли текущий пользователь вызывать все процедуры, возвращаемые методом \ttt{getProcedures}.\\\hline 
	allTablesAreSelectable & boolean & Может ли текущий пользователь использовать все таблицы, возвращаемые методом \ttt{getTables} в выражении \ttt{SELECT}\\\hline 
	autoCommitFailureClosesAll 
	
	ResultSets& boolean & Закрываются ли все ResultSet'ы, если происходит исключение, когда параметр \ttt{autoCommit} установлен в \ttt{true} \\\hline 
	DatabaseMajorVersion & int &  Мажорная версия СУБД\\\hline 
	DatabaseMinorVersion & int & Минорная версия СУБД \\\hline 
	DatabaseProductName & String & Имя СУБД\\\hline 
	DatabaseProductVersion & String & Версия СУБД\\\hline 
	dataDefinitionCauses
	
	TransactionCommit & boolean & DDL-запросы в транзакции запускают её коммит\\\hline 
	dataDefinitionIgnoredIn
	
	Transactions & boolean & СУБД игнорирует DDL-запросы в транзакции \\\hline 
	DefaultTransactionIsolation & int & Уровень изоляции транзакции по умолчанию этой базы данных.\\\hline 
	doesMaxRowSizeIncludeBlobs & boolean & Учитываются ли типы данных JDBC  \ttt{LONGVARCHAR} и \ttt{LONGVARBINARY} для результата функции \ttt{getMaxRowSize}.\\\hline 
	DriverMajorVersion & int & Мажорная версия драйвера \\\hline 
	DriverMinorVersion & int & Минорная версия драйвера \\\hline 
	DriverName & String & Имя драйвера \\\hline 
	DriverVersion & String & Версия драйвера\\\hline 
	ExtraNameCharacters & String & Символы, которые могут использоваться в именах без кавычек\\\hline 
	generatedKeyAlwaysReturned & boolean & \\\hline 
	hashCode & --- & \\\hline 
	IdentifierQuoteString & String & Строка, используемая для цитирования идентификаторов SQL\\\hline 
	isCatalogAtStart & boolean & \ttt{True}, если имя каталога обозначается в начале полного имени таблицы; \ttt{false}, если в конце\\\hline 
	isReadOnly & boolean & База данных в режиме "только для чтения"\\\hline 
	JDBCMajorVersion & int & Мажорный номер версии JDBC для этого драйвера\\\hline 
	JDBCMinorVersion & int & Минорный номер версии JDBC для этого драйвера\\\hline 
	locatorsUpdateCopy & boolean & Сделаны ли изменения в LOB на копии или непосредственно на LOB.\\\hline 
	MaxBinaryLiteralLength & int &  Максимальная длина (в шестнадцатеричных символах) бинарного литерала; 0 означает, что нет предела или лимит не известен\\\hline 
	MaxCatalogNameLength & int & Максимальное количество символов, разрешенных в имени каталога; 0 означает, что нет предела или лимит не известен\\\hline 
	MaxCharLiteralLength & int & Максимальное количество символов символьного литерала; 0 означает, что нет предела или лимит не известен\\\hline 
	MaxColumnNameLength & int & Максимальное количество символов в имени столбца; 0 означает, что нет предела или лимит не известен \\\hline 
	MaxColumnsInGroupBy & int & Какое максимальное число столбцов может быть использовано в предложении \ttt{GROUP BY}; 0 означает, что нет предела или лимит не известен\\\hline 
	MaxColumnsInIndex & int & Максимальное число столбцов, на которых может быть построен индекс; 0 означает, что нет предела или лимит не известен \\\hline 
	MaxColumnsInOrderBy & int & Какое максимальное число столбцов может быть использовано в предложении \ttt{ORDER BY}; 0 означает, что нет предела или лимит не известен\\\hline 
	MaxColumnsInSelect & int & Максимальное число столбцов, которое может быть в результирующем наборе оператора \ttt{SELECT}; 0 означает, что нет предела или лимит не известен\\\hline 
	MaxColumnsInTable & int & Максимальное число столбцов в таблице; 0 означает, что нет предела или лимит не известен \\\hline 
	MaxConnections & int & Максимальное количество одновременных подключений к этой базе данных; 0 означает, что нет предела или лимит не известен\\\hline 
	MaxCursorNameLength & int & Максимальное количество символов в имени курсора; 0 означает, что нет предела или лимит не известен\\\hline 
	MaxIndexLength & int & Максимальное допустимое количество байтов, которые можно индексировать; этот предел включает в себя совокупность всех составляющих частей индекса; 0 означает, что нет предела или лимит не известен\\\hline 
	MaxLogicalLobSize & int & Максимально допустимый размер LOB \\\hline 
	MaxProcedureNameLength & int & Допустимое количество символов в имени процедуры; 0 означает, что нет предела или лимит не известен\\\hline 
	MaxRowSize & int & Максимальное число байт в одной строке; 0 означает, что нет предела или лимит не известен \\\hline 
	MaxSchemaNameLength & int & Максимальное количество символов, которое допустимо в имени схемы\\\hline 
	MaxStatementLength & int & Максимальное количество символов, которое допускается в операторе SQL; 0 означает, что нет предела или лимит не известен\\\hline 
	MaxStatements & int & Максимальное количество активных операторов, которые могут быть открыты одновременно; 0 означает, что нет предела или лимит не известен\\\hline 
	MaxTableNameLength & int & Допустимое количество символов в имени таблицы; 0 означает, что нет предела или лимит не известен\\\hline 
	MaxTablesInSelect & int & Максимальное количество таблиц, разрешенных в операторе \ttt{SELECT}; 0 означает, что нет предела или лимит не известен\\\hline 
	MaxUserNameLength & int & Допустимое количество символов в имени пользователя; 0 означает, что нет предела или лимит не известен\\\hline 
	nullPlusNonNullIsNull & boolean & Будет ли результат равным NULL при конкатенации значений NULL и не-NULL\\\hline 
	nullsAreSortedAtEnd & boolean & Будут ли значения NULL отсортированы в конце независимо от порядка сортировки\\\hline 
	nullsAreSortedAtStart & boolean & Будут ли значения NULL отсортированы в начале независимо от порядка сортировки\\\hline 
	nullsAreSortedHigh & boolean & При сортировке значения NULL будут выводиться в начале\\\hline 
	nullsAreSortedLow & boolean & При сортировке значения NULL будут выводиться в конце\\\hline 
	NumericFunctions & String & Математические функции\\\hline 
	ProcedureTerm & String & Термин, обозначающий хранимую процедуру\\\hline 
	ResultSetHoldability & int & Значение параметра holdability для объектов ResultSet по умолчанию: \ttt{HOLD\_CURSORS\_OVER\_COMMIT} (курсоры ResultSet остаются открытыми при вызове метода \ttt{commit}) или \ttt{CLOSE\_CURSORS\_AT\_COMMIT} (курсоры ResultSet закрываются при вызове метода \ttt{commit})\\\hline 
	SearchStringEscape & String & Строка для экранирования специальных символов. Таких как  \rr{\ttt{\_}} или \rr{\ttt{\%}}, используемых в шаблонах \\\hline 
	SQLKeywords & String &  Список всех ключевых слов SQL СУБД, которые НЕ являются также ключевыми словами SQL:2003\\\hline 
	SQLStateType & int & Тип \ttt{SQLSTATE: sqlStateXOpen} или \ttt{sqlStateSQL}\\\hline 
	storesLowerCaseIdentifiers & boolean & Идентификаторы, не заключенные в кавычки, не чувствительны к регистру и сохраняются в нижнем регистре \\\hline 
	storesLowerCaseQuoted
	
	Identifiers & boolean & Идентификаторы, заключенные в кавычки, не чувствительны к регистру и сохраняются в нижнем регистре\\\hline 
	storesMixedCaseIdentifiers & boolean &  Идентификаторы, не заключенные в кавычки, не чувствительны к регистру и сохраняются в смешанном регистре\\\hline 
	storesMixedCaseQuoted
	
	Identifiers & boolean & Идентификаторы, заключенные в кавычки, не чувствительны к регистру и сохраняются в смешанном регистре\\\hline 
	storesUpperCaseIdentifiers & boolean & Идентификаторы, не заключенные в кавычки, не чувствительны к регистру и сохраняются в верхнем регистре\\\hline 
	storesUpperCaseQuoted
	
	Identifiers & boolean & Идентификаторы, заключенные в кавычки, не чувствительны к регистру и сохраняются в верхнем регистре\\\hline 
	StringFunctions & String & Строковые функции\\\hline 
	supportsAlterTableWithAdd
	
	Column & boolean & Можно ли добавлять новый столбец при изменении существующей таблицы\\\hline 
	supportsAlterTableWithDrop
	
	Column & boolean & Можно ли удалять столбец при изменении существующей таблицы\\\hline 
	supportsANSI92EntryLevelSQL & boolean & Поддерживает ли эта СУБД SQL-стандарт ANSI92 начального уровня\\\hline 
	supportsANSI92FullSQL & boolean & Поддерживает ли эта СУБД SQL-стандарт ANSI92 полностью\\\hline 
	supportsANSI92Intermediate
	
	SQL & boolean & Поддерживает ли эта СУБД SQL-стандарт ANSI92 среднего уровня\\\hline 
	supportsBatchUpdates & boolean & Поддерживает ли эта база данных пакетные обновления.\\\hline 
	supportsCatalogsInData
	
	Manipulation & boolean & Можно ли использовать имя каталога в DML-запросе\\\hline 
	supportsCatalogsInIndex
	
	Definitions & boolean & Можно ли использовать имя каталога в операторе создания индекса\\\hline 
	supportsCatalogsInPrivilege
	
	Definitions & boolean & Можно ли использовать имя каталога в операторе назначения привилегий\\\hline 
	supportsCatalogsInProcedure
	
	Calls & boolean &  Можно ли использовать имя каталога в операторе вызова процедуры \\\hline 
	supportsCatalogsInTable
	
	Definitions & boolean & Можно ли использовать имя каталога в операторе создания таблицы \\\hline 
	supportsColumnAliasing & boolean & Можно ли давать псевдонимы столбцам \\\hline 
	supportsConvert & boolean & Поддерживает ли эта СУБД скалярную функцию JDBC \ttt{CONVERT} для преобразования одного типа JDBC в другой. Типы JDBC - это общие типы данных SQL, определенные в \ttt{java.sql.Types}\\\hline 
	supportsCorrelatedSubqueries & boolean & Поддерживает ли эта СУБД коррелированные подзапросы.\\\hline 
	supportsDataDefinitionAnd
	
	DataManipulationTransactions & boolean & Поддерживает ли СУБД и DDL- и DML-запросы в транзакциях\\\hline 
	supportsDataManipulation
	
	TransactionsOnly & boolean & Поддерживает ли СУБД только DML-запросы в транзакциях\\\hline 
	supportsDifferentTable
	
	CorrelationNames & boolean & Псевдонимы таблиц (если они поддерживаются СУБД) должны отличаться от настоящих имен таблиц \\\hline 
	supportsExpressionsInOrderBy & boolean & Поддерживает ли эта база данных выражения в списках \ttt{ORDER BY}\\\hline 
	supportsFullOuterJoins & boolean & Поддерживает ли СУБД полные внешние соединения\\\hline 
	supportsGetGeneratedKeys & boolean & Retrieves whether auto-generated keys can be retrieved after a statement has been executed\\\hline 
	supportsGroupBy & boolean & Поддерживает ли эта СУБД какую-либо форму предложения \ttt{GROUP BY}\\\hline 
	supportsGroupByBeyondSelect & boolean & Можно ли использовать столбцы, не включенные в \ttt{GROUP BY}, если все столбцы в операторе \ttt{SELECT} включены в \ttt{GROUP BY}\\\hline 
	supportsGroupByUnrelated & boolean & Можно ли использовать столбец, который не находится в инструкции \ttt{SELECT} в предложении \ttt{GROUP BY}\\\hline 
	supportsIntegrityEnhancement
	
	Facility & boolean & Поддерживает ли СУБД SQL Integrity Enhancement Facility\\\hline 
	supportsLikeEscapeClause & boolean & Поддерживает ли СУБД предложение \ttt{LIKE ... ESCAPE}\\\hline 
	supportsLimitedOuterJoins & boolean & Поддерживает ли СУБД ограниченные внешние соединения(\ttt{true}, если \ttt{supportsFullOuterJoins=true})\\\hline 
	supportsMixedCaseIdentifiers & boolean &  Идентификаторы, не заключенные в кавычки, чувствительны к регистру и сохраняются в смешанном регистре\\\hline 
	supportsMixedCaseQuoted
	
	Identifiers & boolean & Идентификаторы, заключенные в кавычки, чувствительны к регистру и сохраняются в смешанном регистре\\\hline 
	supportsMultipleOpenResults & boolean & Возможно одновременное использование нескольких объектов ResultSet, возвращаемых из объекта \ttt{CallableStatement} \\\hline 
	supportsMultipleResultSets & boolean & СУБД поддерживает получение нескольких объектов ResultSet из одного вызова метода \ttt{execute}\\\hline 
	supportsMultipleTransactions & boolean & СУБД позволяет сразу открывать несколько транзакций (в разных соединениях)\\\hline 
	supportsNamedParameters & boolean & Поддерживает ли СУБД именованные параметры\\\hline 
	supportsNonNullableColumns & boolean & Столбцы могут быть обозначены как не-NULL\\\hline 
	supportsOpenCursorsAcross
	
	Commit & boolean & СУБД сохраняет курсоры открытыми между коммитами\\\hline 
	supportsOpenCursorsAcross
	
	Rollback & boolean & СУБД сохраняет курсоры открытыми между откатами\\\hline 
	supportsOpenStatementsAcross
	
	Commit & boolean & СУБД сохраняет операторы открытыми между коммитами\\\hline 
	supportsOpenStatementsAcross
	
	Rollback & boolean & СУБД сохраняет операторы открытыми между откатами\\\hline 
	supportsOrderByUnrelated & boolean & Можно ли использовать столбец, который не находится в инструкции \ttt{SELECT} в предложении \ttt{ORDER BY} \\\hline 
	supportsOuterJoins & boolean &  Поддерживает ли СУБД внешние соединения\\\hline 
	supportsPositionedDelete & boolean & Поддерживает ли эта база данных позиционные операторы \ttt{DELETE} \\\hline 
	supportsPositionedUpdate & boolean & Поддерживает ли эта база данных позиционные операторы \ttt{UPDATE} \\\hline 
	supportsRefCursors & boolean & ---\\\hline 
	supportsSavepoints & boolean & Поддерживает ли СУБД точки сохранения\\\hline 
	supportsSchemasInData
	
	Manipulation & boolean & Можно ли использовать имя схемы в DML-операторе \\\hline 
	supportsSchemasInIndex
	
	Definitions & boolean & Можно ли использовать имя схемы в операторе создания индекса\\\hline 
	supportsSchemasInPrivilege
	
	Definitions & boolean & Можно ли использовать имя схемы в операторе назначения привилегий \\\hline 
	supportsSchemasInProcedure
	
	Calls & boolean & Можно ли использовать имя схемы в операторе вызова процедуры\\\hline 
	supportsSchemasInTable
	
	Definitions & boolean & Можно ли использовать имя схемы в операторе создания таблицы \\\hline 
	supportsSelectForUpdate & boolean & Поддерживает ли СУБД инструкции \ttt{SELECT FOR UPDATE}\\\hline  
	supportsStatementPooling & boolean & Поддерживает ли СУБД кэширование запросов \\\hline 
	supportsStoredFunctionsUsing
	
	CallSyntax & boolean & this database supports invoking user-defined or vendor functions using the stored procedure escape syntax.\\\hline 
	supportsStoredProcedures & boolean & Поддерживает ли СУБД хранимые процедуры\\\hline 
	supportsSubqueriesIn
	
	Comparisons & boolean & Поддерживает ли СУБД подзапросы в выражениях сравнения\\\hline 
	supportsSubqueriesInExists & boolean & Поддерживает ли СУБД подзапросы в выражениях \ttt{EXISTS} \\\hline 
	supportsSubqueriesInIns & boolean & Поддерживает ли СУБД подзапросы в выражениях \ttt{IN} \\\hline 
	supportsSubqueriesIn
	
	Quantifieds & boolean & Поддерживает ли СУБД подзапросы в кванторах \\\hline 
	supportsTableCorrelationNames & boolean &  СУБД поддерживает создание псевдонимов для таблиц \\\hline 
	supportsTransactions & boolean & Поддерживает ли СУБД транзакции. Если нет, вызов метода \ttt{commit} - холостой, а уровень изоляции - \ttt{TRANSACTION\_NONE}\\\hline 
	supportsUnion & boolean & Поддерживает ли СУБД объединение таблиц\\\hline 
	supportsUnionAll & boolean &  Поддерживает ли СУБД объединение таблиц с дубликатами строк на выходе (\ttt{UNION ALL})\\\hline 
	SystemFunctions & String & Список системных функций\\\hline 
	TimeDateFunctions & String & Список функций для работы с датой и временем\\\hline 
	toString & String & --- \\\hline 
	URL & String & URL-адрес подключения\\\hline 
	UserName & String & Имя подключившегося пользователя\\\hline 
	usesLocalFilePerTable & boolean & Использует ли СУБД файл для каждой таблицы\\\hline 
	usesLocalFiles & boolean & Хранит ли СУБД таблицы в локальном файле\\\hline 
\end{longtable}	

\end{document}